% =====================================================================
% N-DHT — Architecture, Implementation, and Deployment
%
%   The complete reference. Bible of the protocol.
%   The technology is shaped by the mission.
%
%   Compile with: tectonic N-DHT-whitepaper.tex
%   Or:           pdflatex N-DHT-whitepaper.tex (twice for cross-refs)
%
%   Requires: tufte-book, tikz, pgfplots, booktabs, microtype, xcolor
% =====================================================================

\documentclass[twoside,nobib]{tufte-book}

\usepackage[utf8]{inputenc}
\usepackage{tikz}
\usepackage{pgfplots}
\pgfplotsset{compat=1.18}
\usetikzlibrary{positioning, calc, arrows.meta, shapes.geometric, fit, backgrounds}
\usepackage{xcolor}
\usepackage{booktabs}
\usepackage{microtype}
\usepackage{amsmath, amssymb}
\usepackage{algorithm, algpseudocode}
\usepackage{multicol}
\usepackage{xspace}

% ── Color palette inherited from nh1-tufte.tex ────────────────────────
\definecolor{rust}{RGB}{185,78,54}
\definecolor{rustsoft}{RGB}{212,168,150}
\definecolor{inkmuted}{RGB}{102,102,102}
\definecolor{bridgegreen}{RGB}{56,118,80}

% ── Protocol-name shorthands ──────────────────────────────────────────
% Used throughout the document. NDHT (the generic family name) is the
% default; specific generations are named only when the chapter is
% generation-specific (Part III, Appendix A).
\newcommand{\NDHT}{N-DHT\xspace}
\newcommand{\NHone}{NH-1\xspace}
\newcommand{\NXseventeen}{NX-17\xspace}
\newcommand{\NXsix}{NX-6\xspace}
\newcommand{\KDHT}{K-DHT\xspace}
\newcommand{\GDHT}{G-DHT\xspace}

% ── Title metadata ────────────────────────────────────────────────────
% tufte-book's automatic \maketitle uses SOUL letterspacing on the
% title which fights with multi-line title content. We build a custom
% half-title page below and leave these as plain values for the PDF
% metadata only.
\title{N-DHT --- Architecture, Implementation, and Deployment}
\author{David A. Smith}
\date{2026-05-04}

% =====================================================================
\begin{document}
% =====================================================================

\frontmatter

% ── Custom half-title page ────────────────────────────────────────────
% Built by hand so we can control typography precisely without fighting
% tufte-book's letterspaced \maketitle.
\thispagestyle{empty}
\begin{fullwidth}
\null\vfill
{\fontsize{48}{52}\selectfont \textsc{N-DHT}\par}
\vspace{0.6em}
{\Large\itshape Architecture, Implementation, and Deployment\par}
\vspace{4em}
{\large David A.~Smith\par}
{\itshape\small YZ.social\par}
\vspace{2em}
{\small Version 0.3.49 \quad Simulator v0.70.05 \quad 2026-05-06\par}
\vfill
\end{fullwidth}
\newpage

% ── Epigraph ──────────────────────────────────────────────────────────
\thispagestyle{empty}
\null\vfill
\begin{fullwidth}
\begin{flushright}
\itshape The technology is shaped by the mission.\par
\end{flushright}
\end{fullwidth}
\vfill
\null\newpage

\tableofcontents

% =====================================================================
\mainmatter
% =====================================================================

% tufte-book applies SOUL letterspacing to running-header section
% titles via \MakeTextLowercase. SOUL's expansion fails on macro-bearing
% titles (including ours, which use \NDHT and protocol shorthands).
% Switching to a plain pagestyle suppresses the letterspaced running
% heads but keeps the page numbers and overall typography. Trade
% acknowledged.
\pagestyle{plain}


% #####################################################################
\part{Foundations}
% #####################################################################
% Why this protocol exists, where it sits in the lineage of distributed
% hash tables, and the vocabulary the rest of the document depends on.


% =====================================================================
\chapter{Introduction}
\label{ch:intro}
% =====================================================================

\begin{quote}
  \itshape The technology is shaped by the mission.
\end{quote}

% --------------------------------------------------------------------
\section{The Mission}
% --------------------------------------------------------------------

\newthought{This document} is the engineering reference for \NDHT, a
learning-adaptive distributed hash table. The protocol is designed for
a specific job: to carry the routing and publish/subscribe workload
of a privacy-first decentralized internet, on consumer devices, in
real browsers, against the actual physics of the network it runs on.
Everything else in this volume --- the architecture, the algorithms,
the simulator, the red team, the deployment path, the application
analysis --- is shaped by that job.

The job is hard for an unfamiliar reason. The internet of 2026 has
no shortage of clever distributed systems, but almost all of them
are \emph{distributed underneath, central on top}: a constellation of
servers that the user reaches through a single corporate
gateway.\sidenote{The pattern is so common we have stopped noticing
it. CDNs, federated identity providers, cloud auth, app stores, the
``serverless'' frameworks that run on three or four hyperscaler
clouds --- all examples of distribution-as-implementation-detail
sitting under a single point of decision.} A real peer-to-peer
internet --- one where two strangers can find each other, exchange a
message, and route around an outage without asking permission of any
intermediary --- requires a different kind of substrate. The substrate
must be \emph{federation-free at every layer}: addressing, routing,
discovery, group communication, eventual delivery under churn.

\textbf{A distributed hash table is the addressing layer of that
substrate.} Without one, every other peer-to-peer primitive
collapses back into a centralized service: there is no way to find
the user you want to talk to, no way to discover the file you want to
download, no way to subscribe to the channel you want to follow,
without somebody telling you where it lives. With one, those
operations are mechanical: the network itself routes you to the right
peer.

The mission is to build the addressing layer that a real peer-to-peer
internet needs. The protocol described here is what that addressing
layer looks like.

% --------------------------------------------------------------------
\section{What a DHT Is, in One Page}
% --------------------------------------------------------------------

\newthought{A distributed hash table} is the
machinery\marginnote{DHTs are the addressing layer behind BitTorrent's
trackerless swarms, IPFS, parts of Ethereum's peer
discovery, and the hidden services on Tor. Every node gets a random
ID; every piece of content gets an ID; lookups route from the source
ID to whoever holds the destination ID, hop by hop, with no
authoritative directory in the middle.} that lets a swarm of
independent computers collectively store and find data without a
central directory. Every peer is assigned an identifier from a large
keyspace. Every piece of content (a file, a topic name, a public key)
is also mapped into the same keyspace through a hash function. To
find a piece of content, a peer asks the closest peer it knows, who
in turn asks a peer closer still --- repeat until the responsible
peer is reached. The trick is making ``closer'' mean something the
network can compute locally, with no global view.

The dominant design --- Kademlia, 2002 --- defines closeness as XOR
distance in a 160-bit ID space.\sidenote{Maymounkov \& Mazi\`eres,
``Kademlia: A Peer-to-Peer Information System Based on the XOR
Metric.''  IPTPS 2002. The XOR metric is symmetric and has a useful
recursive property: a peer that knows one peer in each XOR-bucket
covering the space can route to anyone in the network.} Each lookup
hop at least halves the remaining XOR distance, so any target is
reachable in $\log_2 N$ steps. With a million-node network that is
twenty hops, every time. The math is clean, the algorithm is
provably correct, and Kademlia is in production all over the
internet, twenty-four years after publication.

There is just one problem.

% --------------------------------------------------------------------
\section{The Problem with the Status Quo}
% --------------------------------------------------------------------

\newthought{Twenty hops} sounds fast. But each hop is a message
between two computers somewhere on Earth, and Kademlia node IDs are
random. Two random points on the planet's surface are, on average,
about half the planet apart. A round trip between them is roughly
$100$\,ms.\marginnote{The actual one-way median latency between
random pairs of internet nodes was about $42$\,ms in 2004 and is in
the same neighborhood today (Dabek et al., NSDI 2004,
\S5.1). Round trip is roughly $2\delta$.} Twenty hops at 100\,ms
per hop is two seconds. Every lookup. Every time.

For real-time applications --- voice, video, multiplayer, live
messaging, push notifications, anything a user expects to feel
instant --- two seconds is unusable. The math says routing is
logarithmic, which is fast. The physics says each step is slow.

That tension is the latency tax, and Kademlia pays it in full. There
is a deeper structural cost too: Kademlia treats a lookup to the next
city the same as a lookup to the opposite hemisphere. The XOR metric
is geographically blind. A message destined for a node twenty
kilometers away may bounce through Tokyo, S\~ao Paulo, and Helsinki
before arriving. Real-world workloads are dominated by local
traffic\sidenote{A messaging app's neighbor-to-neighbor chats,
content-sharing services' regional replication, multiplayer games'
lobby discovery --- the vast majority of lookups are between peers
in the same city, the same country, or at worst the same continent.
Long-haul traffic is the minority, but Kademlia weights it equally.} ---
friends in the same city, regional content caches,
same-organization peers --- and every one of those local lookups
pays the global tax.

That is the pothole the next-generation DHT has to fill.

% --------------------------------------------------------------------
\section{What This Protocol Adds}
% --------------------------------------------------------------------

\newthought{The protocol} described in this volume sits on top of
Kademlia's correctness guarantee --- specifically, the property
that a peer with one well-chosen entry per XOR bucket can reach
any target in the network --- and makes three substantive changes
to how routing tables are populated and which peer to ask first.

\textbf{1. Geographic identity.} The first eight bits of every node
ID encode the node's S2 cell prefix --- a geographic
fingerprint.\sidenote{S2 is Google's spherical-geometry library,
which divides Earth's surface into cells along a Hilbert curve.
Hilbert curves are space-filling: consecutive cells on the curve
remain neighbors in space. Two cells whose IDs share a prefix are
geographically close. \S\ref{ch:locality} describes the construction
in detail.} XOR distance now approximately tracks physical
distance. The 20-hop world tour collapses into a 4--6-hop journey
around the neighborhood, where each hop costs single-digit
milliseconds rather than 100. We call this layer \emph{\GDHT},
discussed at length in Chapter~\ref{ch:lineage}.

\textbf{2. Adaptive routing tables.} Each peer maintains a
\emph{synaptome}: a set of weighted connections that strengthens on
successful use and decays through disuse. Routing tables are no
longer populated once and lazily repaired --- they are
\emph{learned}. Successful lookups deposit shortcuts at every
intermediate node along the path; observed transit pairs introduce
themselves directly; geographic neighbors propagate the news. The
network's structure is shaped by the traffic it carries, not by
rules an engineer guessed at. The biology is direct, not
metaphorical: the synaptome is the digital cousin of the synaptic
matrix in a biological brain, and the rule that governs which
connections persist is the same rule (long-term potentiation) that
your brain uses to consolidate memory.

\textbf{3. In-network publish/subscribe.} The protocol delivers a
genuine pub/sub primitive --- a publisher sends to a topic, every
subscriber receives. The implementation is an \emph{axonal tree}: a
delivery structure rooted at a publisher-prefixed topic ID, grown by
routed subscribe messages, healed continuously by the same routed
re-subscribe path that built it. Under 5\% per-cycle node churn the
tree delivers $100\%$ of messages with zero special machinery.

The combined effect is a routing layer that hits within a small
constant of the analytical $3\delta$ latency floor (Dabek et al.,
NSDI 2004) on global lookups, that converges to single-hop
performance on regional traffic, that delivers pub/sub at the
robustness ceiling against churn, and that distributes its
bandwidth load broadly across the population --- with no node
saturating disproportionately at the scales we have been able to
measure. Chapter~\ref{ch:eval-perf} reports the numbers; the rest of
the document explains how they are obtained, tested, and what they
do and do not justify.

% --------------------------------------------------------------------
\section{What This Document Is}
% --------------------------------------------------------------------

\newthought{This is the engineering reference} for the protocol --- the
document from which a working production system will be built.
That puts a high bar on the contents. Three commitments shape what
follows.

\textbf{Completeness.} Every algorithm, every parameter, every
constant, every default, every test, every measurement, every
deferred concern. If a future implementer needs to know it to build
the system, it lives here. The current document is approximately
\textit{[final page count]} pages, with vector figures throughout and
no length limit imposed on clarity.

\textbf{Honesty.} The simulator that generates every quantitative
claim in this volume has limits. Its omissions are catalogued
explicitly\sidenote{See Chapter~\ref{ch:sim-fidelity} for what the
simulator does and does not model, and Chapter~\ref{ch:redteam} for
the third-party adversarial review of the protocol's deployment
risks.} --- there is no claim made on the basis of an unmodeled
behavior, and no measurement reported without surfacing the
mechanism that generated it. The simulator is part of the design
discipline (\S\ref{ch:sim-philosophy}), not a tool we used and put
away.

\textbf{Lineage.} A protocol design without a clear genealogy is
either a copy or a bluff. Every mechanism in \NDHT has an ancestor
in twenty-five years of distributed-systems literature
(Chapter~\ref{ch:lineage}), and most have an ancestor in a hundred
years of neuroscience (Chapter~\ref{ch:substrate}). Where we have
followed prior art, we cite it. Where we have departed, we explain
why and what we tried that did not work
(Appendix~\ref{app:nx-evolution}).

% --------------------------------------------------------------------
\section{Audience and How to Read It}
% --------------------------------------------------------------------

\newthought{The Tufte design} is deliberate. The body of every
chapter is written for a reader who wants to understand the system
without first having read three textbooks --- a high-school student
with a strong math background, an engineer in an adjacent field, a
journalist, a policymaker. The margins carry the technical
precision\marginnote{This is one. Margin notes carry definitions,
units, citations, and the technical asides that would otherwise
interrupt a sentence. Reading them is optional; ignoring them costs
the reader nothing in continuity.} that makes the body's claims
defensible: definitions, units, citations, the technical asides
that would otherwise interrupt the prose. The reader chooses how
deep to go.

The chapters are organized so that successive parts assume the
previous parts. \emph{If you only have an hour:} read
Chapter~\ref{ch:intro} (this one), \ref{ch:operations} (the five
operations), and \ref{ch:eval-perf} (the headline performance
numbers). \emph{If you are an engineer planning a deployment:} add
Part~III (the NH-1 implementation), Part~IV (the simulator), and
Part~VI (deployment and hardening). \emph{If you are evaluating the
protocol against an alternative:} read Chapter~\ref{ch:lineage} (the
DHT lineage) and Chapter~\ref{ch:eval-bandwidth} (the bandwidth
concentration analysis). \emph{If you are building an application
on top of \NDHT:} read Part~VII (existing applications: a deep
dive). \emph{If you want to understand why every design decision is
the way it is:} read the whole document.

% --------------------------------------------------------------------
\section{Scope of the Document}
% --------------------------------------------------------------------

\newthought{What this document covers}: the routing layer
(Chapters~\ref{ch:substrate}--\ref{ch:locality}), the pub/sub layer
(Chapter~\ref{ch:pubsub}), the consolidated NH-1 implementation of
both (Chapters~\ref{ch:nh1-rules}--\ref{ch:nh1-pubsub}), the
simulator that we use to test them
(Chapters~\ref{ch:sim-philosophy}--\ref{ch:sim-tests}), measured
performance against three comparison protocols
(Chapters~\ref{ch:eval-method}--\ref{ch:eval-wins}), the deployment
pathway and operational requirements
(Chapters~\ref{ch:redteam}--\ref{ch:operability}), and a deep
analysis of how existing applications would change if they migrated
to this routing substrate (Chapter~\ref{ch:apps}).

\textbf{What it does not cover}: forward-looking applications enabled
by performance characteristics that this routing layer makes
possible but that depend on additional design work
(streaming-class media, real-time multiplayer, federated AI training,
V2V coordination). Those analyses are deferred to a separate
volume because each requires substantial protocol-specific
groundwork on top of the routing layer described here. We do not
make claims about them in this document. Treating them as future
volumes also lets the present document remain what it is: the
engineering reference for the routing and pub/sub substrate, and
nothing else.

% --------------------------------------------------------------------
\section{Summary of Contributions}
% --------------------------------------------------------------------

The protocol contributes the following:

\begin{enumerate}
\item \textbf{A unified vitality model} for routing-table
maintenance that consolidates eighteen rules from prior generations
(NX-1 through NX-17) into twelve, governed by a single equation:
\textsc{vitality} $=$ \textsc{weight} $\times$ \textsc{recency}.
This is described in Chapters~\ref{ch:operations}
and~\ref{ch:nh1-rules}.

\item \textbf{Five-operation organizational frame} (NAVIGATE, LEARN,
FORGET, EXPLORE, STRUCTURE) under which every routing rule is
classified, making the rule set comprehensible and ablation-amenable
rather than an inscrutable tangle. Chapter~\ref{ch:operations}.

\item \textbf{Performance within a small constant of the analytical
$3\delta$ floor} (Dabek et al., NSDI 2004) at every tested network
size from 5{,}000 to 50{,}000 nodes, and \emph{decreasing} latency
on regional workloads as the network scales --- the opposite of
Kademlia's behavior. Chapter~\ref{ch:eval-perf}.

\item \textbf{An axonal pub/sub primitive} delivering $100\%$ of
messages at baseline and recovering to $100\%$ within three refresh
intervals after $5\%$ instantaneous churn, without K-closest
replication, gossip, parent tracking, or any failure-detection
machinery beyond the routed re-subscribe that built the tree.
Chapter~\ref{ch:pubsub}.

\item \textbf{Empirical resolution of the bandwidth-concentration
question} that has historically dogged adaptive routing protocols.
At every tested scale, \NDHT distributes load broadly (Gini
$\sim$0.55--0.69) while Kademlia and \GDHT concentrate it
catastrophically (Gini $\sim$0.87--0.94, with $56$--$62$ nodes
processing more than $100\times$ the network mean at $50{,}000$
nodes). The volume cost of \NDHT's load distribution is a single
tunable parameter (\texttt{annealRateScale}) with a clean knee.
Chapter~\ref{ch:eval-bandwidth}.

\item \textbf{An open-source simulator} of approximately
\textit{[final line count]} lines of JavaScript, instrumented at
every routing chokepoint, capable of reproducing every measurement
in this volume from a clean checkout in under \textit{[final
runtime]}.\sidenote{The simulator is described in
Chapters~\ref{ch:sim-philosophy}--\ref{ch:sim-tests}; the
methodology is given in Chapter~\ref{ch:eval-method}; the source
is available at the URL listed in the References.}

\item \textbf{A red-team analysis} that catalogues thirteen
deployment-relevant gaps between the simulator's environment and
the production internet, ranked by deploy-blocker status and
time-to-fix, with the bandwidth-concentration concern
empirically reclassified.  Chapter~\ref{ch:redteam}.
\end{enumerate}

% --------------------------------------------------------------------
\section{A Note on Naming}
% --------------------------------------------------------------------

\newthought{Throughout this volume} the family name \emph{\NDHT}
refers to the neuromorphic-DHT family of protocols in general:
the architecture, the operations, the vitality model. When a chapter
discusses a specific generation under detailed scrutiny ---
typically NH-1, occasionally NX-17 or earlier members of the lineage
--- we name the generation explicitly. \GDHT and \KDHT (Kademlia)
appear when they serve as comparison protocols. The
generation-specific chapters live in Part~III; the rest of the body
discusses the family as a whole.

The fact that the most recent generation, NH-1, fits the family-level
description \emph{most cleanly} is not coincidence: NH-1 is the
generation in which the consolidation work converged. Earlier
generations carried more bespoke machinery; future generations will
inherit NH-1's structure and continue to refine it. We expect this
naming convention to remain useful for several iterations.


% =====================================================================
\chapter{The DHT Lineage}
\label{ch:lineage}
% =====================================================================

\begin{quote}
  \itshape A protocol design without a clear genealogy is either a
  copy or a bluff.
\end{quote}

\newthought{The neuromorphic DHT} has two parents. One is twenty-five
years of distributed-systems research --- the line of work that runs
from Plaxton's prefix routing in the late 1990s through Chord, CAN,
Pastry, Tapestry, Kademlia, Coral, Vivaldi, and a dozen named
mechanisms in between. The other is a hundred years of neuroscience,
which we discuss in Chapter~\ref{ch:substrate}. This chapter walks
the first lineage. Every mechanism in \NDHT has an ancestor in the
literature, and most have several. Where we have followed prior art,
we cite it; where we have departed, we explain in terms of what came
before what we did and what we tried that did not work.

The chapter is organized chronologically by paper publication. Each
protocol gets the same five-element treatment: \emph{the routing
primitive}, \emph{the problem the protocol solved}, \emph{the
problem it left open}, \emph{what \NDHT inherits from it}, and
\emph{where \NDHT continues from it}. The chapter closes with a
table that maps each protocol's mechanisms to the five operations of
\NDHT (Chapter~\ref{ch:operations}).

% --------------------------------------------------------------------
\section{Chord (2001)}
% --------------------------------------------------------------------

\newthought{Chord}\sidenote{Stoica, Morris, Karger, Kaashoek,
Balakrishnan, ``Chord: A Scalable Peer-to-Peer Lookup Service for
Internet Applications.'' SIGCOMM 2001.} placed every node on a
circular identifier space modulo $2^m$ and made distance asymmetric:
\emph{successor distance}. To find the node responsible for key $k$,
a Chord node consults its \emph{finger table} of $m$ pointers ---
the $i$-th finger points to the successor of $n + 2^i$ --- and
forwards to whichever finger is closest to $k$ without overshooting.

Each hop at least halves the remaining distance, so lookups complete
in $O(\log N)$. The original paper proved correctness, churn
resilience under stabilization, and load-balance properties under
random hashing.

\textbf{What it solved.} The first widely-cited DHT to give a
proof of $O(\log N)$ lookup in the average case, plus a maintenance
protocol (periodic stabilization) that kept the ring connected
under churn. Chord made the field's central claim --- that you
could build a global lookup service with no
servers\marginnote{The contemporaneous paper \emph{Looking up data
in P2P systems} (Balakrishnan, Kaashoek, Karger, Morris, Stoica,
\emph{CACM} 2003) framed the field. Chord, CAN, Pastry, Tapestry, and
Kademlia all appeared within an eighteen-month window in 2001--2002,
and the comparison literature treats them as a cohort.} ---
defensible.

\textbf{What it left open.} \emph{Geographic blindness}: Chord IDs
are random hashes, so the ring topology is unrelated to the physical
network. \emph{Maintenance overhead}: stabilization runs on a fixed
schedule whether the network is changing or not (the failure mode
that motivates Ghinita-Teo, below). \emph{Pub/sub}: not part of the
protocol. \emph{Adaptation}: the finger table is structurally
determined --- $i$-th finger is $2^i$ ahead, full stop. There is no
mechanism by which observed traffic shapes the routing structure.

\textbf{What \NDHT inherits.} The succession-style argument that a
peer with one well-chosen entry per power-of-two distance bucket
can route to anyone in the network. The XOR metric of Kademlia (and
of \NDHT) is the symmetric variant of Chord's asymmetric successor
distance.

\textbf{Where \NDHT continues.} The fixed finger table becomes a
vitality-scored synaptome whose entries are chosen by traffic, not
by structural mandate. The fixed-schedule stabilization becomes
event-triggered repair (temperature reheat on dead-peer discovery,
LTP reinforcement on successful traffic).

% --------------------------------------------------------------------
\section{CAN (2001)}
% --------------------------------------------------------------------

\newthought{CAN}\sidenote{Ratnasamy, Francis, Handley, Karp,
Schenker, ``A Scalable Content-Addressable Network.'' SIGCOMM 2001.}
embedded peers in a $d$-dimensional Cartesian
toroidal\marginnote{A torus in this sense is a $d$-dimensional cube
with opposite faces identified. The routing argument depends on the
torus topology so that there are no edges to fall off; in practice
$d$ is small (4--10).} space. Each node owned a hyperrectangular
zone; lookups were forwarded to the neighbor closest to the target
along any one axis, and routing terminated in $O(d \cdot N^{1/d})$
hops. With $d = \log_2 N$ the protocol matches Chord's logarithmic
bound.

\textbf{What it solved.} A constructive argument that a DHT could be
built on geometric intuition rather than ID-space arithmetic. CAN
also gave the first clean account of \emph{zone splitting} as a
join procedure --- a new node finds an overloaded neighbor and
takes half its zone --- which influenced every subsequent
load-balance discussion.

\textbf{What it left open.} For practical $d$, the lookup hop count
$O(N^{1/d})$ is much worse than Chord's $O(\log N)$; for $d = \log
N$ the routing-table size exceeds Chord's. Geographic blindness
again. No pub/sub. No adaptation.

\textbf{What \NDHT inherits.} The notion that locality of
\emph{any} kind in the routing structure produces measurable
benefit. The \GDHT geographic prefix
(Chapter~\ref{ch:locality}) is in spirit a Cartesian-style
embedding, restricted to two dimensions on the surface of a sphere
and stamped only into the high bits of a node's identifier.

\textbf{Where \NDHT continues.} CAN's geometry was static; \NDHT's
locality structure is reinforced by traffic and decayed by disuse.
The S2 prefix is a starting condition, not a fixed truth.

% --------------------------------------------------------------------
\section{Pastry (2001) and SCRIBE (2002)}
% --------------------------------------------------------------------

\newthought{Pastry}\sidenote{Rowstron \& Druschel, ``Pastry:
Scalable, decentralized object location and routing for large-scale
peer-to-peer systems.'' Middleware 2001.} took a different turn:
\emph{prefix routing}. Each node maintained three structures.

\begin{itemize}\setlength\itemsep{0.2em}
  \item A \textbf{routing table $R$}: $\lceil \log_{2^b} N \rceil$
    rows, each row covering peers that share a prefix with this
    node up to that row's depth. At base $b = 4$ and $N = 10^6$,
    that is roughly seventy-five entries.
  \item A \textbf{leaf set $L$}: the sixteen to thirty-two
    \emph{numerically-closest} peers, used for the final routing
    hop and for ring-wraparound continuity.
  \item A \textbf{neighborhood set $M$}: the sixteen to thirty-two
    \emph{RTT-closest} peers --- not used for routing, but kept
    fresh as a candidate pool for repairing $R$.
\end{itemize}

A lookup matches one more digit of the destination ID per hop; the
final hop uses $L$ to find the absolute closest peer.

The contribution Pastry is most cited for is \emph{locality
preserving join}: when a new node $X$ joins through a nearby
existing node $A$, each node along the path $A \to \mathrm{root}(X)$
sends $X$ its tables, and \emph{$X$ takes row $n$ of its routing
table from the $n$-th node along the route}. A triangulation
argument preserves locality without exhaustive search.

\textbf{SCRIBE}\sidenote{Castro, Druschel, Kermarrec, Rowstron, ``SCRIBE: A large-scale and decentralized application-level
multicast infrastructure.''  IEEE JSAC 2002.}, layered on Pastry, is
the structural ancestor of \NDHT's axonal pub/sub. Subscribe routes
toward a topic ID; \emph{every node along the path records the
subscriber}; when a publisher publishes, the message follows the
\emph{reverse paths of all the subscribes} to form a multicast
tree. \NDHT's axonal trees use precisely the same construction
(Chapter~\ref{ch:pubsub}).

\textbf{What it solved.} The three-tier routing state introduced
the canonical separation between \emph{routing} (the table $R$),
\emph{terminal hop} (the leaf set $L$), and \emph{maintenance pool}
(the neighborhood set $M$). The locality-preserving join made
routing tables RTT-aware at startup.

\textbf{What it left open.} Locality is set once at join and lazily
repaired thereafter; there is no mechanism by which traffic
\emph{evolves} the locality structure. The leaf set is fixed-size
and ID-proximal; there is no notion of a peer earning its place
through usage. SCRIBE multicast trees are fixed by routing
topology; the only repair is re-subscribe, but it lands at the
same parent every time.

\textbf{What \NDHT inherits.} Three things, deeply. \textbf{(1)} The
three-tier routing state; in \NDHT the tiers collapse into a
single vitality-scored synaptome whose high-, mid-, and
low-vitality entries play the leaf-set, routing-table, and
neighborhood-set roles respectively. \textbf{(2)} The
locality-preserving join, generalized to XOR strata
(Chapter~\ref{ch:nh1-rules}). \textbf{(3)} SCRIBE's reverse-path
multicast --- this is exactly what an axonal tree is.

\textbf{Where \NDHT continues.} The leaf set's terminal-hop role is
played by traffic-earned high-vitality peers, not the
ID-numerically-closest. SCRIBE's multicast tree is fixed at
construction; \NDHT's axonal tree adapts via incoming-promotion and
re-subscribe rebinding (\S\ref{ch:pubsub}). Pastry's locality is
static-after-bootstrap; \NDHT's locality evolves via LTP, hop
caching, lateral spread, and triadic closure.

% --------------------------------------------------------------------
\section{Tapestry and DOLR (2004)}
% --------------------------------------------------------------------

\newthought{Tapestry}\sidenote{Zhao, Huang, Stribling, Rhea, Joseph,
Kubiatowicz, ``Tapestry: A Resilient Global-Scale Overlay for
Service Deployment.'' IEEE JSAC 2004.} was Pastry's near-twin from
the Berkeley side of the field. Same prefix-routing primitive (base
$\beta = 16$, $\log_{16} N$ hops); same general routing-table
shape. Two contributions distinguished it.

First, \emph{at each routing-table slot, the entry stored is the
closest node by RTT that matches the prefix at that level}. The
routing table is built at node insertion via iterative
nearest-neighbor search. This is a stronger commitment to locality
than Pastry's neighborhood-set repair: locality is in the routing
slots themselves, not just adjacent to them.

Second, the \textbf{DOLR} (Decentralized Object Location and
Routing) API. Where a DHT's interface is \texttt{put(key, value)} /
\texttt{get(key)}, DOLR offers \texttt{PublishObject(GUID,
server)}, \texttt{RouteToObject(GUID)},
\texttt{RouteToNode(nodeID)}. When a server stores object $O$ at
GUID, it routes a \emph{publish} message toward $O$'s deterministic
root; \emph{every node along the publish path stores $\langle GUID,
server \rangle$ as a soft-state location pointer}. Queries route
toward the root and \emph{intersect the publish path early}. The
``tree'' of replica locations emerges from the union of publish
paths, not from explicit construction.

Tapestry also introduced \textbf{surrogate routing}: if a
deterministic root is dead, traffic routes to the closest live ID
in its place, transparently. The location pointers along the
former root's path remain valid, and the network heals around the
loss.

\textbf{What it solved.} Three of the four design pillars of
\NDHT: locality (in the routing table), soft-state replication
(via publish-path location pointers), and resilience (via
surrogate routing and Patchwork-style background probes). The
median Relative Delay Penalty\marginnote{Tapestry's signature
metric: \emph{RDP} = (Tapestry route latency) $\div$ (direct IP
latency). Median $\approx 1$ in the wide area; 90th percentile
$\sim 2$--$3$ with location-pointer optimization. This is the
metric \NDHT's $3\delta$-floor analysis is the descendant of.}
in the wide-area was reported as $\approx 1$.

\textbf{What it left open.} Like Pastry, locality was set at node
insertion and not subsequently evolved. The publish-path location
pointers were soft-state with periodic repair, but the repair
schedule was static. There was no learning over which routes
actually carry traffic.

\textbf{What \NDHT inherits.} Three pillars. \textbf{(1)} Locality
in the routing primitive --- in \NDHT, AP scoring's latency
penalty plays the same role Tapestry's RTT-closest-per-slot played,
made dynamic. \textbf{(2)} Soft-state replication along the route
--- this is structurally identical to \NDHT's \emph{hop caching}
(\S\ref{ch:operations}): every intermediate node along a successful
path remembers the destination. \textbf{(3)} Surrogate routing as
the failure-mode handler --- this is exactly \NDHT's
\emph{iterative fallback} (Chapter~\ref{ch:operations}), inherited
through NX-4. Patchwork-style probing is the family of mechanism
that \NDHT's temperature reheat and dead-peer eviction belong to.

\textbf{Where \NDHT continues.} Tapestry's locality is static
after bootstrap; \NDHT's locality is continuously evolved via LTP,
triadic closure, hop caching, and lateral spread. Tapestry's
pub/sub was layered (Bayeux); \NDHT's is built into the protocol
(axonal trees). Tapestry's deployment was server-class on
PlanetLab; \NDHT targets browser-class WebRTC.

The Berkeley OceanStore family is the closest-shaped historical
ancestor of the entire \NDHT architecture. We are an evolutionary
step on that line, not a clean-room reinvention.

% --------------------------------------------------------------------
\section{Kademlia (2002)}
% --------------------------------------------------------------------

\newthought{Kademlia}\sidenote{Maymounkov \& Mazi\`eres, ``Kademlia:
A Peer-to-Peer Information System Based on the XOR Metric.'' IPTPS
2002.} took the simplest of the early designs and shipped it. Two
ideas; both productively conservative.

The first was the XOR metric. Distance between IDs $a$ and $b$ is
$d(a, b) = a \oplus b$, interpreted as an integer. The metric is
\emph{symmetric} (unlike Chord's successor distance) and obeys a
useful triangle-inequality variant. The symmetry meant that
maintenance was cheaper: $A$'s view of distance to $B$ is the same
as $B$'s view of distance to $A$, so entries learned in one
direction are equally informative in the other.

The second was \emph{$k$-buckets and iterative
$\alpha$-parallel \texttt{FIND\_NODE}}. Each node maintained $k = 20$
peers per XOR-distance bucket\marginnote{$k = 20$ is a tradeoff:
small enough that maintenance traffic stays bounded, large enough
that even with significant churn at least one peer per bucket is
likely live. The number is empirically chosen; it is the same
$k$ \NDHT inherits as the bucket size in \GDHT and as the synaptome
core size in NH-1's bootstrap.} (160 buckets at 160-bit IDs,
covering distance ranges $[2^i, 2^{i+1})$). A lookup issued $\alpha
= 3$ parallel asynchronous \texttt{FIND\_NODE} queries to the
closest known peers and iterated until the $k$ closest reachable
peers had been identified. Bucket eviction was LRU-with-old-bias:
\emph{live old contacts were kept}, a deliberate hedge on Saroiu's
empirical observation that longer uptime predicts longer uptime.

Kademlia is the routing substrate that \emph{every \NDHT node still
uses at its bottom layer}. The XOR metric, the bucket-per-stratum
structure, the $k = 20$ and $\alpha = 3$ defaults, the iterative
lookup --- these all carry through to NH-1 unchanged. What \NDHT
adds, it adds on top.

\textbf{What it solved.} A correctness proof that survived
contact with reality. Twenty-four years after publication,
Kademlia is in production on BitTorrent's Mainline DHT, IPFS /
libp2p, Ethereum devp2p, and Tor v3 hidden services. It is the
default the rest of the field is measured against.

\textbf{What it left open.} Geographic blindness, exhaustively.
Kademlia treats a 500\,km lookup the same as a 12{,}000\,km lookup;
no aspect of the protocol references the physical network. No
pub/sub. No adaptation. No load awareness. The bucket eviction
policy preserves stable contacts but does not respond to traffic.

\textbf{What \NDHT inherits.} The XOR metric. Iterative
$\alpha$-parallel lookup as the underlying routing primitive.
$k$-buckets as the synaptome's stratum structure. The $k = 20$ /
$\alpha = 3$ defaults. The Saroiu bias toward stable contacts ---
\NDHT's vitality function inherits this through the
\texttt{recency} factor, which gives recently-confirmed entries
elevated standing.

\textbf{Where \NDHT continues.} Bucket eviction by LRU-with-old-bias
becomes vitality-scored eviction (\S\ref{ch:operations}); the
``stable contacts'' notion generalizes to ``contacts the actual
traffic uses.'' The fixed bucket structure becomes a single
synaptome whose entries are placed by traffic. Two-hop AP scoring,
LTP reinforcement, hop caching, and lateral spread are all added
as separate mechanisms over the Kademlia substrate.

% --------------------------------------------------------------------
\section{Vivaldi (2004)}
% --------------------------------------------------------------------

\newthought{Vivaldi}\sidenote{Dabek, Cox, Kaashoek, Morris,
``Vivaldi: A Decentralized Network Coordinate System.'' SIGCOMM 2004.}
is not a DHT; it is a \emph{coordinate system} for one. Each node
continuously adjusts a synthetic position vector --- typically in
three-dimensional Euclidean space plus a small \emph{height}
component to capture last-mile asymmetry --- so that the
\emph{Euclidean distance} between two nodes' coordinates
approximates the \emph{measured round-trip time} between them.

The mechanism is a spring-force update: every observed RTT pulls
or pushes the local coordinate toward better consistency.
Confidence-weighted, decentralized, requires no coordinator. Over
many observations the system converges to a stable configuration
where pair-wise Euclidean distances $\approx$ pair-wise RTTs.

The result is that \emph{any node can predict RTT to any peer it
has never directly measured}, purely from coordinates. Routing
layers can then choose ``nearby'' peers without an active probe
round.

Vivaldi was deployed in Coral CDN, Azureus / Vuze BitTorrent, and
several content-addressable overlays. It is the most influential
predecessor for latency-aware peer-to-peer design.

\textbf{What it solved.} A learned, low-cost, decentralized way to
rank peers by latency without explicit measurement. Vivaldi is the
canonical answer to ``how does a peer-to-peer system know which of
its peers is closer without asking them?''

\textbf{What it left open.} Vivaldi gives you a coordinate; what
you do with it is up to the routing layer. Convergence requires
sustained traffic, and an attacker who falsifies their coordinate
disrupts neighbors' position estimates.

\textbf{What \NDHT inherits.} The principle that latency is a
queryable property of the network, not an external measurement
event. \NDHT's AP scoring (Chapter~\ref{ch:operations}) factors
observed per-synapse latency into routing decisions; the synapse's
\texttt{latency} field is the descendant of a Vivaldi-style
position estimate, computed from haversine distance in the
simulator and intended to be replaced by a Vivaldi-style learned
coordinate in a future implementation
(\S\ref{ch:future}).\sidenote{The \GDHT geographic prefix is a
\emph{structural} approximation of locality; a Vivaldi-style
coordinate would be the \emph{learned} version. \NDHT today uses
the structural form for simplicity; the migration path to a
learned form is described in Chapter~\ref{ch:production}.}

\textbf{Where \NDHT continues.} Vivaldi-as-replacement-for-S2 is a
clean future direction; in the meantime, the structural locality
of the geographic prefix gives us most of the benefit at a small
fraction of the implementation cost.

% --------------------------------------------------------------------
\section{Coral DSHT (2004)}
% --------------------------------------------------------------------

\newthought{Coral}\sidenote{Freedman, Freudenthal, Mazi\`eres,
``Democratizing Content Publication with Coral.'' NSDI 2004.}
coined the term \emph{distributed sloppy hash table}. Each node
measured RTT to peers it encountered and joined \emph{multiple
nested clusters} at increasing RTT thresholds (e.g., $< 20$\,ms,
$< 60$\,ms, global). Each cluster ran its own internal DHT.

Lookup was local-first: a node queried the tightest cluster it
belonged to; if no result was found, it expanded outward. Most
queries terminated in the local cluster. \emph{Cross-cluster
fan-out happened only when necessary.}

The ``sloppy'' semantics relaxed the strict DHT promise. A key
could be stored at \emph{any} node in the lookup cluster, and
queries returned the \emph{first} hit, not the canonical
XOR-closest one. This made Coral a CDN: the same content was
replicated at many local nodes, and clients pulled from the
nearest available.

\textbf{What it solved.} A working content-distribution overlay
that ran in production from 2004 to roughly 2015 (Coral CDN),
serving millions of cache requests per day at peak as a free CDN
for academic and small-publisher sites.

\textbf{What it left open.} Coral's hierarchical-cluster design
\emph{enforces} locality structurally; nodes \emph{know} what is
near because they measure it at join. There is no learning over
which peers are most useful within a cluster; cluster boundaries
are static thresholds.

\textbf{What \NDHT inherits.} The principle that local-first
routing dominates global routing in real workloads. Coral's
nested-cluster structure is the same observation that motivates
\NDHT's S2 prefix and AP-scoring's latency penalty.

\textbf{Where \NDHT continues.} Coral makes a hard structural
choice (multiple nested DHTs, sloppy semantics); \NDHT keeps a
single flat synaptome and \emph{learns} locality through usage.
Coral relaxes the DHT promise to win locality; \NDHT preserves the
strict DHT promise and adds shortcuts on top.

% --------------------------------------------------------------------
\section{Adaptive Stabilization (2006)}
% --------------------------------------------------------------------

\newthought{Ghinita and Teo}\sidenote{Ghinita \& Teo, ``An Adaptive
Stabilization Framework for Distributed Hash Tables.'' IPDPS 2006.}
addressed one of the longest-running problems in DHT design: how
often should a peer check that its routing table is still correct?

Most DHTs --- Chord, Pastry, CAN --- run \emph{periodic
stabilization}: a fixed-interval timer pings neighbors and refreshes
pointers. Their core claim was that a fixed rate is wrong almost
everywhere. Too low and lookup failure spikes during churn bursts
(their data: $23.7\%$ failure at $5$ kills/sec). Too high and
communication overhead reaches $400\%+$ even during quiet periods.

\textbf{Mechanism.} Each peer estimates its local node-failure rate
$\mu$ and node-join rate $\lambda$ in a rolling window. From those
plus an estimate of network size $N$, it computes the probability
that a given pointer's target has died, $P_{\mathrm{dead}} = 1 -
e^{-\mu \Delta t}$, and the probability that a new joiner has come
to sit between this pointer and its ideal target,
$P_{\mathrm{inacc}}$. Stabilization splits into two channels:

\begin{itemize}\setlength\itemsep{0.2em}
  \item A \emph{liveness check} (one ping, $O(1)$) fires when
    $P_{\mathrm{dead}} \cdot P_{\mathrm{fwd}}$ exceeds a threshold.
  \item An \emph{accuracy check} ($O(\log N)$ lookup) fires when
    $P_{\mathrm{inacc}}$ exceeds a threshold.
\end{itemize}

The system is operated against a single tunable knob: the target
lookup-failure rate $P_f$. The peer adjusts its own thresholds to
hit it.

\textbf{Headline result.} On Chord under variable churn, Adaptive
Stabilization achieved $2.2\%$ peak failure at $172\%$ overhead
versus periodic stabilization's $7.5\%$ failure at $420\%$
overhead. \emph{$3.4\times$ lower failure, $2.4\times$ lower
overhead simultaneously}, by self-tuning rather than hand-picking
a rate.

\textbf{What it solved.} The fixed-cadence problem. The instinct is
recognizable as the closest match in mechanism family of any
protocol \NDHT can be compared against.

\textbf{What it left open.} Ghinita-Teo tunes \emph{cadence} ---
when to check --- not \emph{content} --- which entries to keep.
Their stabilization brings a Chord routing table back toward its
canonical XOR-ideal. \NDHT goes the other direction: the routing
table evolves \emph{away} from the canonical ideal toward a
traffic-shaped optimum.

\textbf{What \NDHT inherits.} The control-loop architecture:
local observation $\to$ probabilistic model $\to$
threshold-triggered protocol response. \NDHT's temperature reheat
on dead-peer discovery, the per-lookup probabilistic anneal, and
the proposed load-aware AP scoring all belong to the same family.

\textbf{Where \NDHT continues.} \NDHT lacks the closed-form
analytical model and the explicit QoS knob; instead, vitality is a
heuristic correlate that consolidates many such signals into one
score. The synthesis would be \emph{NH-1 with Ghinita-style
statistical estimation of churn driving adaptive anneal/reheat
parameters} --- the metaplastic NH-1 we describe in
\S\ref{ch:future}.

% --------------------------------------------------------------------
\section{The Geographic DHT (this work)}
% --------------------------------------------------------------------

\newthought{The geographic DHT}\marginnote{This is our own prior
work, included in the lineage because \NDHT inherits and extends
it. Earlier variants of \GDHT (\texttt{geo}, \texttt{geoa}) are
retired; the current design (\texttt{geob}) is what we discuss
here and use as the comparator throughout this volume.} (\GDHT) is
a single-pass change to Kademlia: stamp the high \textsc{geoBits}
$=$ 8 bits of every node ID with the node's S2 cell prefix.

\[
\mathrm{nodeId} = \underbrace{\mathrm{S2~cell~prefix}}_{8~\mathrm{bits}}
\;\Vert\;
\underbrace{H(\mathrm{publicKey})}_{56~\mathrm{bits}}
\]

The XOR metric is unchanged. The bucket structure is unchanged.
The lookup algorithm is unchanged. But XOR distance now
\emph{approximates physical distance}, because S2 is a Hilbert
space-filling curve and consecutive cells along the curve are
adjacent on the Earth's surface (Chapter~\ref{ch:locality}).

\textbf{What it solved.} The latency tax for regional traffic.
Regional lookups (under 2{,}000\,km) drop from $\sim 510$\,ms in
plain Kademlia to $\sim 150$\,ms in \GDHT at $25{,}000$ nodes ---
a $3.4\times$ improvement, achieved by changing \emph{only the
high bits of node IDs}.

\textbf{What it left open.} The S2 prefix is self-declared. A
node can claim any prefix it wants; the protocol does not
verify. The benign failure mode is degraded routing (mis-located
peers misroute traffic). The adversarial failure mode is a Sybil
swarm in a target cell. Locality is structural and static; it
does not adapt to observed traffic. Pub/sub is still bolted on
via Kademlia-style $K$-closest replication, which drifts under
churn.

\textbf{What \NDHT inherits.} The S2 cell prefix, unchanged. The
stratified-bootstrap construction (Chapter~\ref{ch:locality}). The
geographic-prefix injection point in node IDs.

\textbf{Where \NDHT continues.} Locality becomes
\emph{traffic-shaped} via LTP and the four reinforcement
mechanisms (\S\ref{ch:operations}); pub/sub becomes a built-in
axonal-tree primitive instead of a layered K-closest replication
(\S\ref{ch:pubsub}). The verification gap on the S2 prefix
remains and is addressed in the deployment chapter
(Chapter~\ref{ch:redteam}).

% --------------------------------------------------------------------
\section{The NX Series (2024--2026)}
% --------------------------------------------------------------------

\newthought{NX-1 through NX-17} were our experimental sequence.
Each iteration added or removed one mechanism, ran the full
benchmark suite, and either kept the change or retired it.
Appendix~\ref{app:nx-evolution} documents the sequence in detail
and is the right reference for the question \emph{``why is this
mechanism the way it is?''} for any NH-1 mechanism. The capsule
summary: NX-1 was a routing-table Hebbian baseline; NX-3
introduced the geographic three-layer initialization; NX-4 added
iterative fallback; NX-5 added stratified bootstrap and
incoming-synapse promotion; NX-6 added churn-resilient routing
(temperature reheat, dead-synapse eviction); NX-7 through NX-9
were dendritic pub/sub experiments that we eventually retired in
favor of axonal trees; NX-10 introduced
routing-topology forwarding; NX-15 introduced the AxonManager
abstraction; NX-17 closed the publisher-prefix-addressing question
and was the state-of-the-art before the NH-1 consolidation.

\textbf{What the sequence did.} Established empirically that each
mechanism either justified its complexity in the benchmark numbers
or did not. Of the eighteen rules NX-17 carried, twelve survived
the consolidation into NH-1; six were retired because their
contributions could be absorbed into the unified vitality model
without measurable performance loss.

\textbf{What we kept.} Vitality-scored synaptome, AP scoring with
two-hop lookahead, LTP, hop caching, lateral spread, triadic
closure, iterative fallback, temperature reheat, dead-peer
eviction, stratified bootstrap, incoming-synapse promotion,
publisher-prefixed topic IDs.

\textbf{What we retired.} Two-tier (highway) synaptome --- absorbed
into the single vitality-scored tier; stratum-floor enforcement ---
absorbed into the diversity component of vitality (later removed
entirely as ablation showed it was harmful, see
Appendix~\ref{app:nx-evolution}); Markov pre-learning --- not
retained in NH-1; load-tracked AP scoring --- deferred to the
production-pathway chapter (\S\ref{ch:production}); some bespoke
relay-pinning logic --- absorbed into the standard axon-tree
recruit-policy; weight-scale parameter on AP scoring --- ablation
showed it was not significant.

\textbf{Where \NDHT continues from NX-17.} The consolidation
itself: the same routing behavior, fewer parameters, simpler
vitality function, improved tunability. Chapter~\ref{ch:nh1-rules}
describes the consolidated rule set; Appendix~\ref{app:nx-evolution}
gives the full justification.

% --------------------------------------------------------------------
\section{Mapping the Lineage to the Five Operations}
% --------------------------------------------------------------------

\NDHT organizes every routing rule under five operations
(Chapter~\ref{ch:operations}): \textsc{navigate}, \textsc{learn},
\textsc{forget}, \textsc{explore}, \textsc{structure}. The
following table maps the contributions of the protocols above
into that frame, both as a summary of the lineage walk and as a
preview of the architecture chapter that follows.

\begin{table*}[h]
\centering
\small
\caption{Mapping prior-art mechanisms into \NDHT's five operations.
Each row shows what a protocol contributed to the lineage; each
column shows where that contribution lives in \NDHT.}
\label{tab:lineage-map}
\begin{tabular}{@{}lp{0.16\linewidth}p{0.16\linewidth}p{0.14\linewidth}p{0.13\linewidth}p{0.16\linewidth}@{}}
\toprule
\textbf{Protocol} &
\textbf{\textsc{Navigate}} &
\textbf{\textsc{Learn}} &
\textbf{\textsc{Forget}} &
\textbf{\textsc{Explore}} &
\textbf{\textsc{Structure}} \\
\midrule
Chord & succession ring & --- &
periodic stabilization & --- & finger table \\
CAN & Cartesian zones & --- & --- & --- &
zone-splitting join \\
Pastry & prefix routing & --- &
neighborhood-set repair & --- &
locality-preserving join \\
Tapestry & prefix + RTT-closest slots &
soft-state location pointers along publish path &
Patchwork probes &
surrogate routing &
nearest-neighbor join \\
Kademlia & XOR + iterative $\alpha$-parallel &
\emph{(none, but bucket eviction LRU-with-old-bias hints)} &
LRU eviction & $\alpha$-parallel &
$k$-buckets \\
Vivaldi & latency as queryable & coordinate update from observed RTT & --- & --- & --- \\
Coral DSHT & local-first nested-cluster lookup &
RTT-cluster join &
sloppy semantics & --- &
nested cluster construction \\
Ghinita-Teo &
threshold-gated lookup &
local $\mu, \lambda$ estimation &
adaptive cadence &
threshold trigger &
--- \\
\GDHT (this work) & XOR with geo-prefix bias & --- & --- & --- &
S2 cell prefix in ID \\
\midrule
\textbf{\NDHT} &
\textbf{AP scoring + 2-hop lookahead + iterative fallback} &
\textbf{LTP, hop cache, lateral spread, triadic closure,
incoming-synapse promotion} &
\textbf{vitality eviction, decay, dead-peer reheat} &
\textbf{temperature anneal, $\epsilon$-greedy first hop} &
\textbf{S2 prefix, stratified bootstrap, sponsor-chain join,
axonal-tree pub/sub} \\
\bottomrule
\end{tabular}
\end{table*}

The bottom row is the destination of the lineage walk. It is also
the table of contents for Part~II, where each cell becomes a
section.


% =====================================================================
\chapter{Vocabulary}
\label{ch:vocab}
% =====================================================================

\newthought{The terms in this chapter} are used throughout the
volume. Each definition is followed by a chapter reference where the
concept is developed in full. The intent is that this chapter
serves both as a forward-reference for a reader new to the field
and as a back-reference for a reader who has skipped ahead and
needs to look something up.

The glossary is alphabetical. Where a term has a precise formal
definition (an equation, a parameter range, a unit) the formal
definition appears in the margin alongside the prose definition in
the body. The extended glossary in
Appendix~\ref{app:glossary} repeats this material with additional
cross-references to source code symbols, useful when reading the
implementation.

% --------------------------------------------------------------------
\section{Glossary}
% --------------------------------------------------------------------

\begin{description}\setlength\itemsep{0.45em}

\item[\textsc{action potential (AP) scoring}]
The mechanism \NDHT uses to choose the next routing hop. For each
candidate synapse, AP combines XOR-distance progress, the
synapse's learned weight, and the synapse's observed latency into
a single score; the highest-scoring candidate is chosen. The
latency factor halves the score every 100\,ms, so a synapse that
is mathematically a small step closer but physically a large step
closer wins. See \S\ref{ch:operations}.\marginnote{The simplified
default form, ignoring weight bias for clarity:
$\mathit{AP}(s, t) = \mathit{progress} \times 2^{-\ell/100}$,
where progress is the bit
reduction in XOR distance and $\ell$ is the synapse's latency in
milliseconds. Two-hop lookahead extends this with one additional
forward probe.}

\item[\textsc{annealing}]
The exploration mechanism by which a node periodically replaces a
weak synapse with a candidate from its 2-hop neighborhood, biased
toward under-represented strata. Triggered probabilistically per
hop, scaled by the node's current temperature. Inherited from NX-6;
see \S\ref{ch:operations} (\textsc{explore}).

\item[\textsc{annealRateScale}]
A multiplier (default 1.0) on the per-hop annealing trigger
probability. \texttt{annealRateScale = 0.10} reduces total
\NDHT message volume by approximately $5\times$ at no measurable
routing-quality cost (Chapter~\ref{ch:eval-bandwidth}). The
single most important tuning knob in NH-1.

\item[\textsc{axon}]
In neuroscience, the long branching cable that delivers a neuron's
output to many downstream targets. In \NDHT, the data structure
representing a publisher's delivery tree for a topic --- rooted at
the publisher's S2 cell, growing branches as subscribers attach.
See Chapter~\ref{ch:pubsub}.

\item[\textsc{axonal tree}]
The complete delivery structure for a pub/sub topic: a root, an
arbitrary-depth branching of sub-axons, and the subscribers
attached to leaves and intermediate nodes. The tree grows by
routed subscribe messages and heals by routed re-subscribe under
churn. Contrast with K-closest replication (e.g., Pastry SCRIBE
trees), which fix the tree shape at construction time. See
Chapter~\ref{ch:pubsub}.

\item[\textsc{bilateral cap}]
The constraint that any two peers must \emph{both} consent to a
connection: peer $A$ may want to add $B$ to its synaptome, but if
$B$'s connection slots are full, the connection cannot form. This
is the simulator's first-order analogue of WebRTC's per-browser
connection limit (typically 50--100 simultaneous data channels).
See \S\ref{ch:sim-fidelity}.

\item[\textsc{bootstrap}]
The process of populating a new node's synaptome with an initial
set of peers when the node joins the network. \NDHT uses a
\emph{sponsor-chain join} (inherited from Pastry's
locality-preserving join, generalized to XOR strata) where each
node along the lookup path of \texttt{newNodeId} contributes a row
of its own synaptome. See \S\ref{ch:operations} (\textsc{structure})
and \S\ref{ch:nh1-rules}.

\item[\textsc{churn}]
The continuous arrival and departure of nodes from the network. We
quantify churn as a percentage of alive nodes replaced per unit
time. The benchmark suite tests $5\%$ instantaneous churn (the
robustness threshold) and continuous variable-rate churn (the
sustained-load test). See \S\ref{ch:sim-tests} and
Chapter~\ref{ch:eval-perf}.

\item[\textsc{Dabek $3\delta$ floor}]
A 2004 result\sidenote{Dabek, Li, Sit, Robertson, Kaashoek,
Morris. ``Designing a DHT for low latency and high throughput.''
NSDI 2004.} establishing the analytical lower bound on lookup
latency in any recursive DHT: $3\delta$, where $\delta$ is the
median one-way RTT between random pairs of nodes on the network.
The proof is geometric: each hop covers half the remaining
distance, the resulting series sums to $2\delta$, and one final
delivery hop adds a further $\delta$. \NDHT achieves $1.18$--$1.28
\times$ the floor at $25{,}000$ nodes (Chapter~\ref{ch:eval-perf}).

\item[\textsc{decay}]
The slow erosion of synapse weight in the absence of reinforcement.
\NDHT applies a multiplicative decay (\textsc{weight} $\leftarrow$
\textsc{weight} $\times \gamma$) at fixed intervals where
$\gamma = 0.995$ by default. LTP-locked synapses (those within
their \emph{inertia} window) skip decay. See \S\ref{ch:operations}
(\textsc{forget}).

\item[\textsc{Gini coefficient}]
A scalar in $[0, 1]$ summarizing the inequality of a
distribution: $0$ is perfectly even, $1$ is perfectly
concentrated. We use it to characterize per-node traffic
distribution per cycle: \KDHT and \GDHT exhibit Gini $\sim
0.94$ at $50{,}000$ nodes, indicating heavy concentration; \NDHT
exhibits Gini $\sim 0.66$, broadly distributed. See
Chapter~\ref{ch:eval-bandwidth}.

\item[\textsc{hop}]
A single message between two adjacent peers in a routing path.
Latency of one hop is determined by the haversine distance between
the peers' coordinates plus a fixed simulator constant (10\,ms). A
typical \NDHT global lookup at $25{,}000$ nodes completes in 4--6
hops.

\item[\textsc{hop caching}]
The mechanism by which every intermediate node along a successful
lookup path remembers the destination as a directly-routable
peer. After traversing the path $A \to B \to C \to D$, both $B$
and $C$ install a synapse pointing to $D$, with weight $0.5$ and
fresh inertia. Future lookups originating from any peer who reaches
$B$ or $C$ benefit from the shortcut. See \S\ref{ch:operations}
(\textsc{learn}).

\item[\textsc{Hot $>$ 10$\times$ / $>$ 100$\times$ nodes}]
Counts of nodes whose per-cycle traffic exceeds $10\times$ or
$100\times$ the network mean, respectively. The $100\times$
threshold is our operational definition of a bandwidth-saturation
hazard: nodes processing two orders of magnitude above the mean.
\NDHT produces zero such nodes at any tested scale; \KDHT produces
$56$ at $50{,}000$ nodes (Chapter~\ref{ch:eval-bandwidth}).

\item[\textsc{incoming synapse}]
A synapse that was added to a peer's view because \emph{another
peer} reached it for the first time, not because the peer reached
out. Incoming synapses contribute candidates for routing but with
weight $0.1$; if their use count crosses the
\texttt{promoteThreshold} they are promoted to first-class
synapses with weight $0.5$. The mechanism makes the network notice
who is interested in a node, not just who the node is interested
in. See \S\ref{ch:operations} (\textsc{learn}).

\item[\textsc{inertia}]
The protection window after a synapse is reinforced. While
$\mathit{simEpoch} - \mathit{syn.inertia} <$ \textsc{INERTIA\_DURATION}
(default 20 ticks), the synapse is excluded from decay and
eviction. The neuroscience analogue is the synaptic-tagging
mechanism\sidenote{Frey, U., \& Morris, R. G. M. ``Synaptic tagging
and long-term potentiation.'' \emph{Nature} 385, 533--536 (1997).}
that protects newly-strengthened synapses from competing
plasticity. See \S\ref{ch:operations}.

\item[\textsc{iterative fallback}]
The routing behavior when no synapse offers forward XOR-progress.
The lookup falls back to the closest unvisited peer in the
synaptome (regardless of strict XOR-progress), trading
optimality for reachability. The mechanism is \NDHT's surrogate
routing (Tapestry's term); inherited from NX-4. See
\S\ref{ch:operations} (\textsc{navigate}).

\item[\textsc{$k$-bucket}]
The Kademlia bucket structure: $k = 20$ peers per XOR-distance
range $[2^i, 2^{i+1})$. \NDHT's synaptome is bucket-organized at
bootstrap (via stratified bootstrap, \S\ref{ch:locality}) but
allows entries to migrate between buckets through traffic.

\item[\textsc{lateral spread}]
The mechanism by which a successful hop-cache pickup propagates to
the originating node's geographic neighbors. After node $A$ caches
target $C$ via hop caching, $A$ tells its $k = 2$
geographically-nearest synaptome peers about $C$ as well, biased
by their existing weights. See \S\ref{ch:operations} (\textsc{learn}).

\item[\textsc{long-term potentiation (LTP)}]
A neuroscientific phenomenon\sidenote{Bliss, T. V. P., \& L\o mo,
T. ``Long-lasting potentiation of synaptic transmission in the
dentate area of the anaesthetized rabbit.'' \emph{J. Physiol.}
232, 331--356 (1973). Hebb, D. O. \emph{The Organization of
Behavior.} Wiley, 1949 (the conceptual ancestor: ``cells that fire
together, wire together'').} in which repeated stimulation of a
synapse increases its strength for hours to weeks. \NDHT applies
this directly: every hop in a successful lookup that completes
faster than the running EMA receives a weight increment ($+0.05$,
capped at 1.0) on every traversed synapse. See
\S\ref{ch:operations} (\textsc{learn}) and \S\ref{ch:substrate}.

\item[\textsc{long-term depression (LTD)}]
The complement of LTP: synaptic strength erodes through disuse.
In \NDHT this is realized as multiplicative decay of weight at
fixed intervals (see \textsc{decay}).

\item[\textsc{maintenance traffic}]
Messages exchanged for routing-table upkeep rather than for
end-user lookups: ping checks, bucket refreshes, sponsor-chain
queries on join. \NDHT generates very little of this in steady
state; the dominant traffic source is annealing's 2-hop probes.
See Chapter~\ref{ch:eval-bandwidth}.

\item[\textsc{node ID}]
A 64-bit integer addressing a node in the network. In \GDHT and
\NDHT the high \textsc{geoBits} = 8 bits encode the S2 cell
prefix; the low 56 bits are derived from a hash of the node's
public key. See \S\ref{ch:locality}.

\item[\textsc{publisher-prefix topic ID}]
The construction $\mathrm{topicId} = (\text{publisher's S2 prefix})
\Vert H_{56}(\mathrm{topicName})$. The high bits pin the topic's
axon root close to the publisher; subscribers and publishers must
agree on the prefix out-of-band as a topic-naming convention
(typically encoded as \texttt{@XX/domain/event}). Inherited from
NX-17; see Chapter~\ref{ch:pubsub}.

\item[\textsc{recency}]
The exponential-decay multiplier in vitality, capturing how
recently a synapse was reinforced. Within the inertia window
(\textsc{INERTIA\_DURATION} ticks), recency is $1.0$. Outside it,
recency decays as $e^{-\Delta t / \tau}$ where
$\tau =$~\textsc{RECENCY\_HALF\_LIFE}
where the half-life is 50 ticks by default. See
\S\ref{ch:operations}.

\item[\textsc{reinforce}]
The verb form of LTP application: increment a synapse's weight by
$+0.05$, set its inertia field to the current epoch (entering the
protection window), and increment its use count. See
\S\ref{ch:operations} (\textsc{learn}).

\item[\textsc{replay cache}]
The bounded ring buffer (default 100 messages) that each axon node
maintains for each topic, holding recently-published messages.
Late-joining subscribers can request replay from the cache to fill
the gap from missed publishes. See Chapter~\ref{ch:pubsub}.

\item[\textsc{S2 cell prefix}]
Google's S2 library partitions Earth's surface into a hierarchy of
cells along a Hilbert space-filling
curve.\sidenote{\url{https://s2geometry.io/}} The top
\textsc{geoBits} $=$ 8 bits of every \GDHT and \NDHT node ID encode
the cell containing the node's claimed lat/lng. Two cells whose
prefixes are bitwise close are physically adjacent on the Earth's
surface. See \S\ref{ch:locality}.

\item[\textsc{Slice World}]
The adversarial test in which the network is partitioned almost in
half (Eastern hemisphere from Western), connected only through a
single bridge node near Hawaii. The protocol is asked to route
between hemispheres. \KDHT achieves $0\%$ success; \NDHT achieves
$\sim 94\%$ via hop caching and triadic closure converting the
bridge into a seed crystal for new cross-hemisphere edges. See
\S\ref{ch:sim-tests} and Chapter~\ref{ch:eval-perf}.

\item[\textsc{stratified bootstrap}]
A bootstrap construction that allocates the $k$-bucket budget
across XOR strata so that each stratum receives roughly equal
representation, with explicit fill on the geographic-prefix strata
(top \textsc{geoBits}). Inherited from NX-5;
see \S\ref{ch:locality}.

\item[\textsc{stratum}]
An XOR-distance bucket: peers whose IDs differ from the local
node's by a number in $[2^i, 2^{i+1})$ are stratum $i$. \NDHT uses
the term interchangeably with ``bucket.''

\item[\textsc{synapse}]
The data structure for a single peer entry in a synaptome. Fields:
\texttt{peerId}, \texttt{latencyMs}, \texttt{stratum},
\texttt{weight} ($\in [0, 1]$), \texttt{inertia} (epoch of last
reinforcement), \texttt{useCount}, plus a few diagnostic fields.
See \S\ref{ch:substrate}.

\item[\textsc{synaptome}]
The set of synapses a node maintains: its routing table. \NDHT's
synaptome is a single tier (no separate ``highway'' or
``leaf-set'' tier); entries are differentiated by their vitality
score, not by structural placement. Default capacity 50 synapses,
matching the WebRTC connection budget of a typical browser. See
Chapter~\ref{ch:substrate}.

\item[\textsc{temperature}]
A per-node scalar in [\textsc{T\_MIN} $= 0.05$, \textsc{T\_INIT} $=
1.0$] that governs the probability of triggering an annealing
event each hop. Cools by \textsc{ANNEAL\_COOLING} = 0.9997 per
hop the node participates in; reheats to \textsc{T\_REHEAT} =
0.5 on dead-peer detection. See \S\ref{ch:operations}
(\textsc{explore}).

\item[\textsc{tick / simEpoch}]
The simulator's logical time unit: one lookup. The
\texttt{simEpoch} field of a node advances by one each time the
node participates in a lookup. Decay, inertia decay, and recency
are measured in ticks, not wall-clock time.

\item[\textsc{triadic closure}]
The mechanism by which a node observing the same transit pair $A
\to \mathrm{this} \to C$ multiple times \emph{introduces} $A$
directly to $C$. After \textsc{triadicThreshold} = 2 observed
transits of the pair, the introducer hints both $A$ and $C$ to
form a direct synapse. The name comes from social-network theory:
dense triangles emerge in any network shaped by repeated
interaction. See \S\ref{ch:operations} (\textsc{learn}).

\item[\textsc{two-hop lookahead}]
The AP-scoring mechanism that explores not just one hop but two:
for each candidate first hop, \NDHT probes the candidate's
synaptome to estimate where the second hop would land, and
scores the joint pair. The probe set is bounded by
\textsc{lookaheadAlpha} = 5 candidates per hop. See
\S\ref{ch:operations} (\textsc{navigate}).

\item[\textsc{vitality}]
The score by which synapses are admitted, retained, and evicted.
\textsc{vitality} $=$ \textsc{weight} $\times$
\textsc{recency}.\marginnote{The single equation that drives every
admission decision in \NDHT. Earlier generations had eighteen
rules and forty-four parameters; the consolidation around vitality
collapses the bookkeeping into one score. See
Chapter~\ref{ch:operations} for the per-synapse computation and
\S\ref{ch:nh1-rules} for how it replaces the prior rule set.}
When a new candidate synapse wants to enter a full synaptome, the
synapse with the lowest vitality among non-LTP-locked candidates
is evicted. See \S\ref{ch:operations} (\textsc{forget}).

\item[\textsc{weight}]
The strengthening factor in vitality, in $[0, 1]$. Increased by
$+0.05$ per LTP reinforcement, decayed multiplicatively by
$\gamma = 0.995$ per decay tick. Bootstrap synapses begin at
$0.5$; introduced peers (triadic closure, hop caching) start at
$0.5$; annealed peers start at $0.1$. See
\S\ref{ch:operations}.

\end{description}

% --------------------------------------------------------------------
\section{The Five Operations}
% --------------------------------------------------------------------

\NDHT classifies every routing rule under one of five operations,
introduced briefly here and developed in
Chapter~\ref{ch:operations}. They are the organizational frame for
the rest of the document.

\begin{description}\setlength\itemsep{0.4em}

\item[\textsc{navigate}] Choose the next hop for a lookup.
Mechanisms: AP scoring, two-hop lookahead, iterative fallback,
$\epsilon$-greedy first hop, direct-target shortcut.

\item[\textsc{learn}] Strengthen what works on a successful path.
Mechanisms: LTP reinforcement, hop caching, lateral spread,
triadic closure, incoming-synapse promotion.

\item[\textsc{forget}] Decay what is unused; evict by vitality.
Mechanisms: multiplicative weight decay, vitality-based admission,
LTP-locked exemption.

\item[\textsc{explore}] Inject randomness so the protocol does not
lock onto local optima. Mechanisms: temperature-driven annealing,
$\epsilon$-greedy first-hop randomness, dead-peer reheat.

\item[\textsc{structure}] Bootstrap and maintain the network's
basic shape. Mechanisms: S2-prefix node IDs, stratified bootstrap,
sponsor-chain join, axonal-tree pub/sub construction.

\end{description}

The frame is universal in the sense that \emph{any} adaptive
system needs all five: act on present input, learn from feedback,
forget the obsolete, occasionally try something new, maintain
basic structure. \NDHT is one realization of that pattern with
biology-inspired specifics.


% #####################################################################
\part{The N-DHT Architecture}
% #####################################################################
% Generic, not version-specific. The operations and mechanisms that
% define what an N-DHT is, independent of any particular generation.


% =====================================================================
\chapter{The Neuromorphic Substrate}
\label{ch:substrate}
% =====================================================================

\begin{quote}
  \itshape Neurons that fire together, wire together.\\[0.3em]
  \normalfont\textsc{\footnotesize Donald Hebb} \,$\cdot$\, 1949
\end{quote}

\newthought{The neuromorphic DHT} is named for an analogy. The
analogy is exact rather than ornamental: the engineering problem
faced by a peer in a peer-to-peer network is the engineering
problem faced by a neuron in a brain. Both have a hard upper bound
on how many connections they can maintain. Both face vastly more
candidate partners than they can afford to keep. Both must decide,
on the basis of local information only, which connections are
worth retaining.

The brain solved this problem hundreds of millions of years ago.
This chapter explains the biology in enough detail to follow the
translation, and then performs the translation in full. The
mechanism that emerges is not metaphorical. The data structure
that holds the routing table in \NDHT (the \emph{synaptome}) is
the digital cousin of the synaptic matrix of a biological neuron;
the rule that governs which entries persist (the \emph{vitality
function}) is the engineering counterpart of long-term potentiation;
the protected window we apply to recently-reinforced entries is
the synaptic-tagging mechanism described by Frey and Morris in
1997. The rest of the document depends on this translation, so we
do it carefully.

% --------------------------------------------------------------------
\section{The Engineering Problem the Brain Solved}
% --------------------------------------------------------------------

A single neuron in your cerebral cortex connects to roughly $10{,}000$
others through structures called \emph{synapses}.\sidenote{Drachman,
``Do we have brain to spare?'' \emph{Neurology} 64(12), 2004.
Estimates vary by region, age, and methodology; the general figure
of $10^4$ synapses per cortical neuron is broadly accepted.}
There are about $86$ billion neurons total in a human
brain.\sidenote{Azevedo \emph{et al.}, ``Equal numbers of neuronal
and nonneuronal cells make the human brain an isometrically
scaled-up primate brain.'' \emph{J. Comp. Neurol.} 513(5), 2009.}
Any given neuron \emph{could} connect to almost any other in
principle, but maintaining a synapse is metabolically
expensive --- the cell has to manufacture and traffic receptor
proteins, maintain the molecular machinery of vesicle release,
and pay the running cost of resting potentials at every
contact. So neurons live with a hard cap, surrounded by far more
potential partners than they can afford to keep.

How does the brain decide which connections to retain?

The same problem faces a peer in a real-world peer-to-peer
network. A typical web-browser deployment, restricted by WebRTC,
can maintain about $50$--$100$ simultaneous data
channels.\marginnote{Browser-specific: Chrome desktop $\sim 256$;
Firefox $\sim 190$; Safari $\sim 95$; mobile browsers as low as
$\sim 20$. The simulator and most deployment targets in this
document use a $50$-connection cap, calibrated to the lowest of
the desktop tier with headroom for the application's own
WebRTC traffic outside the DHT.} The network might have hundreds of
thousands or millions of nodes; the routing table can hold only a
small fraction of them. Which fraction?

The brain answers this question through a phenomenon called
\emph{long-term potentiation}, abbreviated LTP.

% --------------------------------------------------------------------
\section{Long-Term Potentiation in Five Pieces}
% --------------------------------------------------------------------

In 1973, Tim Bliss and Terje L\o mo published the foundational
experimental result.\sidenote{Bliss, T.~V.~P., \& L\o mo, T.
``Long-lasting potentiation of synaptic transmission in the
dentate area of the anaesthetized rabbit.'' \emph{Journal of
Physiology} 232, 331--356 (1973).} They stimulated a particular
pathway in a rabbit's hippocampus with a high-frequency burst of
electrical pulses --- a few hundred milliseconds of
$10$--$100$\,Hz stimulation --- and observed that, for the
\emph{following several weeks}, the same pathway responded more
strongly to subsequent inputs. The connection had physically
strengthened. It stayed strengthened.

Decades of follow-up research filled in the molecular detail.
Five elements are relevant to our analogy.

\textbf{1.\ The coincidence detector.} Synapses contain a special
kind of glutamate receptor called the
\emph{NMDA}\sidenote{\emph{N-methyl-D-aspartate}: the synthetic
agonist used to identify the receptor in the 1980s. The same
receptor type is found in essentially every excitatory synapse in
the cortex.} receptor, which has an unusual property. Glutamate
binding alone does not open the channel; the post-synaptic
membrane must \emph{also} be sufficiently depolarized at the same
moment. Both conditions in the same temporal
window\marginnote{The timing window is short: roughly $10$--$20$
milliseconds. Two events more than $\sim 50$\,ms apart do not
register as coincident.} are required. NMDA receptors are
biological AND-gates: they fire only when ``the input arrived''
and ``the cell was already active'' both happen.

\textbf{2.\ The strengthening signal.} When the AND condition is
met, calcium ions rush into the post-synaptic cell. Calcium
triggers a cascade of intracellular reactions whose net effect is
that the synapse inserts \emph{more} of a different kind of
receptor (called \emph{AMPA}). After the cascade resolves, the
same incoming signal produces a larger post-synaptic response. The
synapse has become physically more sensitive.

\textbf{3.\ Consolidation.} Over the next half hour to several
hours, the strengthening transitions from molecular to structural.
New proteins are synthesized; the synapse physically grows; in
some cases entirely new spines emerge on the dendrite. The change
becomes durable in a way the early-phase strengthening was not.
This is the difference between \emph{early-LTP} (lasting hours,
reversible) and \emph{late-LTP} (lasting weeks, structural).

\textbf{4.\ Decay.} Connections that don't get used regularly
weaken over time. The complementary process is called
\emph{long-term depression}, or LTD. The molecular mechanism is
roughly the inverse: a different intracellular cascade, triggered
by low-frequency activity, that removes AMPA receptors and
strengthens the inhibitory side of the synapse. Without LTD the
brain would saturate every connection at maximum strength; with
LTD, only the connections that prove useful retain their gain.

\textbf{5.\ Synaptic tagging.} The fifth element is the most
subtle and the one with the largest engineering implication.
Frey and Morris published in 1997\sidenote{Frey, U., \& Morris,
R.~G.~M. ``Synaptic tagging and long-term potentiation.''
\emph{Nature} 385(6616), 533--536 (1997).} that recently-strengthened
synapses receive a \emph{tag} --- a transient molecular state ---
that protects them from being overwritten by other plasticity
events for a brief window. Without the tag, the protein synthesis
required for consolidation would be \emph{first} captured by
whichever synapse fired most recently, regardless of whether it
had the coincidence-detection signal. With the tag, the proteins
are routed to the synapses that the AND-gate said deserved them.

The five together --- coincidence detection, strengthening, decay,
consolidation, tagging --- comprise the brain's solution to the
``which connections to keep'' problem. The pithy summary, which
predates the molecular detail by twenty-four years, is Hebb's:

\begin{quote}\itshape
Cells that fire together, wire together.
\end{quote}

Every artificial neural network you have ever heard
of\marginnote{Including all of contemporary deep learning. The
backpropagation algorithm is a digital descendant of Hebbian
learning: weight updates are proportional to the product of
pre-synaptic and post-synaptic activations. The form changes; the
underlying ``coincidence increases weight'' principle does not.}
ultimately runs on a digital descendant of this rule.

% --------------------------------------------------------------------
\section{From Neuron to Routing Table}
% --------------------------------------------------------------------

\newthought{The translation} from biology to network is direct.
Each row of the table maps a neural mechanism to a routing-protocol
mechanism that does the same engineering work.

\begin{table}[h]
\small
\begin{tabular}{@{}p{0.42\linewidth}p{0.52\linewidth}@{}}
\toprule
\textsc{In the brain} & \textsc{In N-DHT} \\
\midrule
A synapse (point of connection between two neurons) &
A connection to a peer (an entry in the synaptome) \\

Synaptic strength, in $[0, 1]$, gated by AMPA receptor density &
A learned weight on each connection, in $[0, 1]$ \\

LTP: coincidence increases weight &
Successful lookup increments weights along the path \\

LTD: disuse decreases weight &
Each connection's weight slowly decays each tick \\

Synaptic tagging (recent activity protected) &
Recently-used connections are protected from eviction
(\textsc{INERTIA\_DURATION} ticks) \\

Pruning of unused connections at metabolic limit &
Lowest-vitality connection evicted when synaptome is full \\

NMDA AND-gate: input \emph{and} post-synaptic activity &
AP scoring's combination of progress \emph{and} latency \emph{and}
weight \\

Spike-timing dependent plasticity: ``together'' is timed &
LTP applied only on paths fast enough relative to a moving
average of recent path latencies \\

Coincidence at the network level (groups firing together) &
Triadic closure: peers observed in transit pairs are introduced
to each other \\
\bottomrule
\end{tabular}
\end{table}

That table is what \NDHT \emph{is}. The synaptome is a routing
table whose entries have weights; weights go up on success, down
on disuse; recently-strengthened entries are protected briefly;
when the table is full, the lowest-scoring entry is evicted. The
network's ``memory'' is the routing table, and the routing table
evolves into whatever shape best serves the actual traffic.

We will spend the rest of this Part formalizing each row. The
formalization is the architecture: Chapter~\ref{ch:operations}
defines the vitality function and the five operations that
implement it; Chapter~\ref{ch:locality} adds the geographic
prefix that seeds the routing table at bootstrap; Chapter~\ref{ch:pubsub}
extends the same machinery to publish/subscribe via axonal trees.

% --------------------------------------------------------------------
\section{Why the Analogy is Literal}
% --------------------------------------------------------------------

\newthought{Two objections} are reasonable here. The first: ``The
analogy is too neat. Engineering problems borrowed from biology
usually don't survive contact with the engineering details.'' The
second: ``Even if it worked once, biological mechanisms are
shaped by evolutionary accidents that have nothing to do with the
specific problem of network routing.''

The defense is empirical. The vitality function (Chapter~\ref{ch:operations})
collapses eighteen routing rules from prior generations
(NX-1 through NX-17) into twelve, drops the parameter count from
forty-four to twelve, and produces routing within a small constant
of the analytical $3\delta$ floor (Chapter~\ref{ch:eval-perf}).
The \emph{simplification} is the evidence: a forty-four-parameter
collection of bespoke rules is not what an engineer designs from
first principles to solve a problem. It is what the engineer
arrives at when they accumulate fixes and never go back to
audit. A twelve-parameter system organized around one equation is
what you get when the underlying principle is correct and you
let it do the work.

The biology was the principle that let the consolidation happen.
That is the only claim the analogy needs to support.

A second defense is comparative. Other adaptive routing
systems --- Tapestry's soft-state pointers, Coral's nested
clusters, Vivaldi's spring-force coordinate updates,
Ghinita-Teo's adaptive stabilization --- are all attacking the
same problem from different angles. They are sibling solutions in
the same design space (Chapter~\ref{ch:lineage}). The neuromorphic
formulation is the one that comes out simplest. That is not an
accident: the brain has been working on this problem with intense
selective pressure for several hundred million years. We get to
borrow the answer.

% --------------------------------------------------------------------
\section{The Synapse and the Synaptome}
% --------------------------------------------------------------------

\newthought{The data structures} that carry the analogy across.
A \emph{synapse} is the protocol's unit of routing knowledge. A
\emph{synaptome} is the collection of synapses a node maintains.

\textbf{The synapse.} A synapse holds a peer's identity and the
weights and timestamps that govern how it is treated. The
implementation fields are the following.

\begin{itemize}\setlength\itemsep{0.2em}
  \item \texttt{peerId} --- the 64-bit identifier of the peer this
  synapse points to. Combined with the local node ID, this gives
  the synapse its XOR distance and stratum.
  \item \texttt{stratum} --- the count of leading-zero bits of the
  XOR difference; equivalently, the bucket index in a Kademlia-style
  routing table. Used by stratified-bootstrap and stratum-diversity
  rules.
  \item \texttt{latency} --- the synapse's observed round-trip
  latency to the peer, in milliseconds. Used directly by AP scoring
  (\S\ref{ch:operations}).\marginnote{In the simulator,
  \texttt{latency} is computed once at synapse creation as
  $\mathrm{haversine}(A, B) \cdot c^{-1} + \delta_{\text{node}}$,
  where $c$ is the speed of light in fiber and $\delta_{\text{node}}$
  is a fixed per-hop processing cost. In production, \texttt{latency}
  would be a Vivaldi-style learned coordinate or an EMA over
  observed RTTs (\S\ref{ch:future}).}
  \item \texttt{weight} --- a real number in $[0, 1]$. Set to $0.5$
  on bootstrap and on hop-cache pickup; set to $0.1$ on annealing;
  incremented by $+0.05$ per LTP reinforcement (capped at $1.0$);
  decayed multiplicatively per decay tick.
  \item \texttt{inertia} --- the simulator epoch at which this
  synapse was last reinforced. While $\mathit{simEpoch} -
  \mathit{inertia} <$~\textsc{INERTIA\_DURATION} (default $20$
  ticks), the synapse is exempt from decay and from eviction. This
  is the synaptic-tagging analogue.
  \item \texttt{useCount} --- the cumulative count of LTP
  reinforcements this synapse has received. Used by the
  incoming-synapse promotion rule (\S\ref{ch:operations}) and as
  a debug signal.
  \item \texttt{bootstrap} --- a flag indicating this synapse was
  installed by the bootstrap process rather than by traffic.
  Bootstrap synapses receive a small admission preference under
  vitality competition\sidenote{NX-6 parity. The preference is
  worth a few percentage points of routing quality at bootstrap
  cost vs. recovered convergence; ablation showed it was net
  positive at $25{,}000$ nodes.}: when the lowest-vitality
  candidate would otherwise be a bootstrap entry, eviction
  prefers the next-lowest non-bootstrap entry instead.
  \item \texttt{\_addedBy} --- a debug field carrying the rule that
  added the synapse (\texttt{bootstrap}, \texttt{hopCache},
  \texttt{lateralSpread}, \texttt{triadic}, \texttt{evictReplace},
  \texttt{anneal}, \texttt{promote}). Not consulted by routing;
  used in observability (\S\ref{ch:operability}).
\end{itemize}

\textbf{The synaptome.} The synaptome is the JavaScript \texttt{Map}
keyed by \texttt{peerId}, holding the live synapses. Default
capacity is $50$. When a new synapse wants to enter and the
synaptome is at capacity, the vitality-based admission gate
(\S\ref{ch:operations}, \emph{\textsc{forget}}) decides which
existing synapse to evict.

\textbf{Incoming synapses.} A second, parallel structure. When peer
$A$ uses the local node as an intermediate hop in $A$'s routing,
the local node does not necessarily already know about $A$. The
local node logs a record in its \emph{incoming synapses} map ---
``$A$ has reached me'' --- with weight $0.1$, in case $A$ proves
useful as a peer for outgoing traffic later. The mechanism is
called \emph{incoming-synapse promotion}: if the same incoming
peer is reached for routing more than \textsc{promoteThreshold}
times, the entry is upgraded to a regular synapse with weight
$0.5$ via vitality-based admission. This makes the network notice
who is interested in a node, not just who the node is interested
in.

The biological analogue is \emph{retrograde signaling} --- a
post-synaptic neuron signaling back across the synapse to the
pre-synaptic side, modulating its release behavior. The mechanism
exists because traffic in a neural circuit is not strictly one-way,
and a peer's view of who-cares-about-it is not strictly one-way
either.

% --------------------------------------------------------------------
\section{What This Substrate Buys}
% --------------------------------------------------------------------

\newthought{The substrate} described in this chapter is not yet a
routing protocol. It is the data structure on which a routing
protocol can be built. What \NDHT does on top of the substrate
constitutes the rest of Part~II:

\begin{itemize}\setlength\itemsep{0.3em}
  \item Chapter~\ref{ch:operations} defines five operations
    (\textsc{navigate}, \textsc{learn}, \textsc{forget},
    \textsc{explore}, \textsc{structure}) and the rules within
    each. The five operations together fully specify how the
    synaptome is manipulated during routing.
  \item Chapter~\ref{ch:locality} adds geographic structure to the
    bootstrap process: node IDs encode position via S2, the initial
    synaptome respects strata weighted toward locality, and routes
    converge to short physical paths in the common case.
  \item Chapter~\ref{ch:pubsub} extends the substrate to support
    publish/subscribe. The result is the axonal-tree primitive,
    which inherits the same vitality logic for tree maintenance.
\end{itemize}

By the end of Part~II, the substrate is operational: every
mechanism the protocol uses to route lookups, deliver pub/sub, and
heal under churn is fully specified. Part~III then describes the
specific NH-1 implementation choices that took the architecture
from eighteen rules and forty-four parameters in NX-17 to twelve
rules and twelve parameters in NH-1, and Part~V measures what the
combination delivers.


% =====================================================================
\chapter{Five Operations, One Vitality Model}
\label{ch:operations}
% =====================================================================

\begin{quote}
  \itshape Act, learn from feedback, forget the obsolete,
  occasionally try something new, maintain basic structure. \\[0.3em]
  \normalfont\textsc{\footnotesize the universal pattern}
\end{quote}

\newthought{This chapter} is the architectural heart of the
document. Every routing decision \NDHT makes falls into one of
five operations: \textsc{navigate} (act on present input),
\textsc{learn} (update from feedback), \textsc{forget} (decay the
unused), \textsc{explore} (occasionally try something new), and
\textsc{structure} (maintain the basic shape of the network). The
operations together constitute the protocol. The single equation
that drives admissions, retentions, and evictions is
\emph{vitality}, defined first.

The frame is universal: any adaptive system will need the same
five operations in some form. The specifics of each operation are
shaped by the substrate (Chapter~\ref{ch:substrate}) and by the
peer-to-peer setting. We describe each operation at the level of
mechanism --- what computation is performed, on what state, in
what order --- with formal pseudocode reserved for
Chapter~\ref{ch:nh1-routing} once the NH-1 specifics are
introduced.

% --------------------------------------------------------------------
\section{The Vitality Function}
\label{sec:vitality}
% --------------------------------------------------------------------

\textbf{Vitality} is a real number in $[0, 1]$ that scores every
synapse. The definition is one line:

\begin{equation}
  v(s) \;=\; \mathit{weight}(s) \;\times\; \mathit{recency}(s)
  \label{eq:vitality}
\end{equation}

\textbf{Weight} (the strengthening factor, in $[0, 1]$) accumulates
under reinforcement and erodes under disuse. The dynamics are
described in \S\ref{sec:learn} (additive reinforcement) and
\S\ref{sec:forget} (multiplicative decay). At any instant,
$\mathit{weight}(s)$ is the value the LTP and decay rules have
left.

\textbf{Recency} (the timing modulator, in $[0, 1]$) captures how
recently the synapse was reinforced. Within the
\emph{inertia window} of \textsc{INERTIA\_DURATION} ticks (default
$20$) following reinforcement, recency is locked at $1.0$ ---
this is the synaptic-tagging analogue from
\S\ref{ch:substrate}.\sidenote{The lock prevents the failure mode
where a freshly-strengthened synapse loses to an older, more
weight-saturated competitor in eviction. The fix preserves
\emph{recently learned} signals long enough to be reinforced
again.} Outside the window, recency decays exponentially with
half-life \textsc{RECENCY\_HALF\_LIFE} ticks (default $50$). The
formal expression:

\begin{equation}
  \mathit{recency}(s) \;=\;
  \begin{cases}
    1 & \text{if } \mathit{simEpoch} - s.\mathit{inertia} < \mathit{T_{lock}} \\
    \max(0.1, e^{-(\mathit{simEpoch} - s.\mathit{inertia}) / \tau})
      & \text{otherwise}
  \end{cases}
\end{equation}

where $\tau =$~\textsc{RECENCY\_HALF\_LIFE}. The floor of $0.1$
prevents very old synapses from contributing zero vitality even
when they have non-trivial weight; this matters in low-traffic
regions where a synapse may be useful but used rarely.

That is the entire vitality function. It is the only score the
admission gate (\S\ref{sec:forget}) consults.

\textbf{Why this product, and not a sum?} A sum would let a
high-weight, very-old synapse beat a freshly-introduced one with
moderate weight indefinitely. A product makes recency a
first-class veto: a synapse that has not been reinforced in many
ticks is automatically a weak candidate, regardless of its
historical weight. The brain works the same way --- consolidated
synapses still need recent activity to remain useful.

% --------------------------------------------------------------------
\section{NAVIGATE: Acting on Present Input}
\label{sec:navigate}
% --------------------------------------------------------------------

\newthought{Navigation} is the operation by which a node selects
the next routing hop for a given lookup. Five mechanisms operate
in priority order; the first one to succeed determines the next
hop.

\textbf{Mechanism 1: Direct shortcut.} If the synaptome contains
a synapse pointing directly to the target ID, use it. Trivial in
practice (the lookup is one hop), but matters when hop caching
(below) has installed the destination earlier.

\textbf{Mechanism 2: Epsilon-greedy first hop.} On the very first
hop of a lookup, with probability \textsc{EPSILON} (default
$0.05$), pick a candidate uniformly at random from those that
make XOR progress. This is the protocol's exploration-vs-exploitation
hedge: a purely greedy router would lock onto the first decent
shortcut and never discover better ones.

\textbf{Mechanism 3: Action-Potential (AP) scoring with two-hop
lookahead.} The default mechanism. AP scoring combines
XOR-distance progress with observed latency into a single score
per candidate; the highest-scoring candidate is chosen. The
formula uses progress in bits and latency in milliseconds, with
latency entering exponentially:

\begin{equation}
  \mathit{AP}(s, t) \;=\; \mathit{progress}(s, t) \;\times\;
  2^{-\ell(s) / 100}
  \label{eq:ap}
\end{equation}

where $\mathit{progress}(s, t)$ is the reduction in XOR distance
from the current node's distance to $t$ to $s$.peer's distance to
$t$ (interpreted as an integer; can be negative for non-progress
candidates), and $\ell(s)$ is the synapse's observed latency in
milliseconds. The exponential gives the score a half-life of
$100$\,ms in latency: a synapse that is mathematically a small
step closer but physically a large step closer wins, by a factor
proportional to how much closer.

\emph{Two-hop lookahead.} A pure single-hop AP score is greedy:
it picks the locally-best peer regardless of where the peer's
peers sit. For each of the top \textsc{lookaheadAlpha} candidates
(default $5$), \NDHT extends the search one step: it inspects the
candidate's synaptome (already known to the local node from
previous interactions or sponsor-chain bootstrap) for the
candidate's best forward-progress synapse, and computes a joint
two-hop AP score. The candidate whose two-hop projection is best
becomes the actual chosen first hop. The mechanism doubles the
effective routing budget at small cost (the candidate's synaptome
is local data, no extra round trip required).

\textbf{Mechanism 4: Iterative fallback.} If no synapse offers
forward XOR progress (the lookup has hit a local minimum), the
fallback is to pick the closest unvisited peer in the synaptome
\emph{regardless} of strict progress. This trades optimality for
reachability and is the mechanism that prevents catastrophic
failures at partition boundaries. Inherited from NX-4. Tapestry
called the same mechanism \emph{surrogate routing}.

\textbf{Mechanism 5: Termination.} If even the iterative fallback
yields no candidate (every synapse visited or every synapse
dead), the lookup fails. The simulator records the failure for
the success-rate calculation. In practice this happens
exceedingly rarely: at $25{,}000$ nodes under normal conditions
the success rate is $100\%$ across all our tested workloads.

The five mechanisms together handle every navigation case the
protocol encounters. Chapter~\ref{ch:nh1-routing} gives the full
pseudocode for a lookup tick, including the operation-by-operation
walk-through.

% --------------------------------------------------------------------
\section{LEARN: Updating from Feedback}
\label{sec:learn}
% --------------------------------------------------------------------

\newthought{Learning} happens after a successful lookup. \NDHT
applies five reinforcement mechanisms; the first three run on
every successful path, and the last two run on observed transit
patterns.

\textbf{1. Long-term potentiation (LTP).} Every synapse traversed
in a successful path that was \emph{faster than the running
average} of recent path latencies receives a reinforcement event:

\begin{itemize}\setlength\itemsep{0.15em}
  \item Increment $\mathit{weight}$ by $+0.05$, capped at $1.0$.
  \item Set $\mathit{inertia}$ to the current simulator epoch
    (entering the \textsc{INERTIA\_DURATION} protection window).
  \item Increment $\mathit{useCount}$ by $1$.
\end{itemize}

The ``faster than average'' filter is essential. If LTP fired on
\emph{every} successful path, every synapse on a slow recovery
path would also be reinforced, and the protocol would learn
\emph{whichever paths exist}, not \emph{which paths are fast}. The
filter is implemented as: maintain an exponential moving average
of recent path latencies (per node), apply LTP only on paths with
total latency at or below the current EMA. The EMA itself updates
with each successful path, so the threshold is adaptive.

\textbf{2. Hop caching.} Every intermediate node along a successful
path remembers the destination as a directly-routable peer. After
traversing $A \to B \to C \to D$, both $B$ and $C$ install a
synapse pointing to $D$ (via vitality-based admission, weight
$0.5$, fresh inertia). The next lookup originating anywhere ---
not just from $A$ --- that needs $D$ and reaches $B$ or $C$ will
take the shortcut. This is the source of the dramatic
hop-count reductions we see for concentrated-destination workloads.

The mechanism is the structural ancestor of Tapestry's soft-state
location pointers (Chapter~\ref{ch:lineage}); we make it active
on every successful path rather than only on
\texttt{PublishObject} events.

\textbf{3. Lateral spread.} A successful hop-cache pickup
propagates to the originating node's geographic neighbors. After
$A$ caches $D$, $A$ tells its $k =$~\textsc{LATERAL\_K} $= 2$
geographically-nearest synaptome peers about $D$ as well, biased
by their existing weights. The mechanism is named after the
biological observation that activated neurons influence their
spatial neighbors via diffusive signaling.\sidenote{The biological
analogue is loose; the engineering function is precise. Lateral
spread converts a one-to-one cache hit into a one-to-three (or
$k$+1) cache deposit, dramatically accelerating shortcut
formation in dense regional clusters.}

\textbf{4. Triadic closure.} A node observing the same transit
pair $A \to (\text{this node}) \to C$ on multiple lookups
introduces $A$ directly to $C$. The transit count is held in a
local cache; when it crosses
\textsc{TRIADIC\_THRESHOLD}~$=2$, the node sends an introduction
hint to $A$ pointing at $C$. The name is from social-network
theory: dense triangles emerge in any network shaped by
repeated interaction.

The biological analogue is the network-level coincidence
counterpart to NMDA receptors. NMDA detects coincidence at
\emph{single synapses}; triadic closure detects coincidence at the
\emph{network} level, between groups of peers that keep flowing
through the same intermediate.

\textbf{5. Incoming-synapse promotion.} An incoming synapse
(\S\ref{ch:substrate}) records that another peer reached this
node. If the same incoming peer is reached for routing more than
\textsc{PROMOTE\_THRESHOLD} $= 2$ times, the protocol promotes
the incoming entry to a first-class outgoing synapse with weight
$0.5$, via vitality-based admission. The mechanism turns observed
inbound interest into established routing capacity.

% --------------------------------------------------------------------
\section{FORGET: Decaying the Unused}
\label{sec:forget}
% --------------------------------------------------------------------

\newthought{Forgetting} is the operation that keeps the
synaptome's content responsive to current traffic rather than
historical accumulation. Three mechanisms.

\textbf{1. Periodic decay.} Every \textsc{DECAY\_INTERVAL} $= 100$
lookups in which the local node participates, every synapse not
under LTP-lock has its weight multiplied by $\gamma = 0.995$:

\begin{equation}
  \mathit{weight}(s) \;\leftarrow\; \mathit{weight}(s) \cdot \gamma
\end{equation}

After fifty unrefreshed decay ticks (5{,}000 lookups), a synapse
that started at weight $1.0$ has fallen to $0.78$; after a
hundred, $0.61$; after two hundred, $0.37$. The shape is gentle
enough that occasional-but-real synapses retain useful weight,
sharp enough that genuinely-unused synapses fade to
eviction-eligible territory within a few decay periods.

\textbf{2. Vitality-based admission.} When a new candidate synapse
wants to enter a full synaptome, the rule is:

\begin{enumerate}\setlength\itemsep{0.15em}
  \item Compute the vitality of every existing synapse not under
    LTP-lock.
  \item Among the non-bootstrap synapses, identify the lowest-vitality
    candidate $v_{\min}$.
  \item If $v_{\min}$ is below the new candidate's projected
    vitality, evict it and admit the new synapse. Otherwise,
    refuse admission.
\end{enumerate}

The bootstrap-synapse preference is a safety hedge: bootstrap
entries are never first-evicted as long as a non-bootstrap
candidate is available, because they encode the structural backbone
of the routing table. If \emph{every} candidate is a bootstrap
synapse, the gate falls back to the lowest-vitality bootstrap entry.

\textbf{3. LTP-lock exemption.} Synapses within the inertia
window ($\mathit{simEpoch} - \mathit{inertia} <$ \textsc{INERTIA\_DURATION})
are exempt from both decay and eviction. This is the synaptic-tagging
mechanism in action: recently-strengthened synapses are protected
from competing plasticity for a brief window.

The combination of the three mechanisms gives \NDHT its
characteristic behavior: \emph{the synaptome's content drifts toward
whatever the current traffic is using}, with enough hysteresis that
isolated bursts don't permanently disrupt the routing table.

% --------------------------------------------------------------------
\section{EXPLORE: Avoiding Local Optima}
\label{sec:explore}
% --------------------------------------------------------------------

\newthought{A purely greedy} adaptive system locks onto the
first decent solution it finds and stops looking. \NDHT injects
randomness through three mechanisms.

\textbf{1. Temperature.} Each node carries a scalar
$T \in [T_{\min} = 0.05, T_{\mathrm{init}} = 1.0]$ (the named
constants are \textsc{T\_MIN} and \textsc{T\_INIT}). After
each hop the node participates in, $T$ is multiplied by
\textsc{ANNEAL\_COOLING} $= 0.9997$ down to the floor. After
several thousand hops a typical node sits at the floor, $T = 0.05$.

\textbf{2. Probabilistic annealing.} On each hop, with probability
$T \cdot r$ (where $r$ is \textsc{annealRateScale}), the node attempts an
annealing step: pick a weak existing synapse, scan the 2-hop
neighborhood for an under-represented stratum, find a candidate,
and replace the weak synapse with a synapse to the candidate. The
\textsc{annealRateScale} parameter (default $1.0$) is the master
volume control on annealing rate;
Chapter~\ref{ch:eval-bandwidth} describes how it serves as the
bandwidth-vs-distribution tuning knob.

The biological analogue is exploratory dendritic outgrowth: a
neuron that has spare capacity occasionally extends a new branch
toward an under-represented region. The engineering rationale is
the same: keep some uncommitted slots available so that emerging
patterns of useful traffic can be picked up.

\textbf{3. Reheat on dead-peer detection.} When a routing decision
encounters a dead peer (a synapse pointing at a node that no
longer responds), the protocol \emph{reheats}: the node's
temperature is raised to \textsc{T\_REHEAT} $= 0.5$. Subsequent
hops will fire annealing more aggressively until the temperature
cools again. This is the mechanism by which \NDHT recovers
synaptome composition under churn without needing a periodic
maintenance schedule (Ghinita-Teo's adaptive-stabilization
problem solved differently; Chapter~\ref{ch:lineage}).

\textbf{4. Epsilon-greedy first hop.} Already mentioned in
\S\ref{sec:navigate}, but listed here too because it is also an
exploration mechanism. The epsilon hedge is small ($5\%$ of first
hops) but persistent, ensuring that the protocol always explores
\emph{some} fraction of off-greedy first-hop choices regardless
of synaptome state.

% --------------------------------------------------------------------
\section{STRUCTURE: Maintaining Network Shape}
\label{sec:structure}
% --------------------------------------------------------------------

\newthought{Structure} is the operation that maintains the
underlying graph topology: new nodes joining, the bootstrap
process, and the maintenance of basic connectivity invariants.

\textbf{1. Geographic node IDs.} Every node ID encodes the
node's S2 cell prefix in the high bits (Chapter~\ref{ch:locality}).
This is a structural seed: it makes XOR distance correlate with
physical distance from the moment the node joins, before any
traffic-driven learning has happened.

\textbf{2. Stratified bootstrap.} A new node's initial synaptome
is populated by a stratified algorithm: enough entries in each
XOR-distance bucket to satisfy the routing-table connectivity
guarantee, with explicit fill on the geographic-prefix strata
(top \textsc{geoBits}) so locality is operational immediately. The
budget allocation is described in detail in
Chapter~\ref{ch:locality}; the principle is that the synaptome
should be \emph{usable for routing} from the first lookup, not
require a learning warmup.

\textbf{3. Sponsor-chain join.} A new node joins through a
sponsor: it routes a self-lookup query through the sponsor toward
its own ID. Each node along the resulting path contributes a row
of its synaptome to the joining node, in the same locality-preserving
manner as Pastry's join (Chapter~\ref{ch:lineage}). The result is
that the new node arrives with synaptome entries that are
RTT-close to it, courtesy of the path's locality.

\textbf{4. Bilateral cap enforcement.} Every connection edge is
enforced bilaterally: peer $A$ can want $B$ in its synaptome, but
the connection only forms if $B$'s connection slots also have
room. The cap is the simulator's first-order analogue of WebRTC
connection limits and the structural constraint behind every
admission decision.

\textbf{5. Axon-tree construction.} Pub/sub topics build delivery
trees through the same mechanism: a routed subscribe message
follows the routing layer, and every intermediate node along the
path becomes a potential axon-tree branch point. The full
treatment is in Chapter~\ref{ch:pubsub}; here we note that the
axon tree is structurally part of the protocol, not a layered
add-on.

% --------------------------------------------------------------------
\section{The Universality of the Five-Operation Frame}
% --------------------------------------------------------------------

\newthought{The frame is universal.} Any system that adapts to
its environment must implement all five operations in some form.
\textsc{navigate} is the act-on-input operation: every system
must produce \emph{some} output for current input. \textsc{learn}
is the credit-assignment operation: every system must update its
internal state from feedback or it will not adapt.
\textsc{forget} is the entropy-resistance operation: every
system must shed obsolete state or it will saturate.
\textsc{explore} is the local-optima-avoidance operation:
without it, a system locks onto the first usable solution and
stops improving. \textsc{structure} is the topology-maintenance
operation: every system needs some account of how its parts are
connected.

The same frame describes how reinforcement learning systems work
(action selection, reward update, value-function decay,
$\epsilon$-greedy, environment model), how immune systems work
(antigen recognition, clonal expansion, naive-cell turnover,
hypermutation, repertoire diversification), and how social
networks form (contact, friendship reinforcement, drift, novelty
seeking, structural maintenance). \NDHT is one realization of
this pattern with biology-inspired specifics tuned to the
peer-to-peer routing problem.

The unification is not just an organizational convenience.
Earlier generations (NX-1 through NX-17) carried eighteen rules
across ad-hoc categories; the same rules organized into five
operations made it visible which mechanisms did the same job at
different abstraction levels, which made consolidation possible.
Twelve rules and twelve parameters, with one equation (vitality)
governing every admission decision, is the consolidation that
became NH-1 (Chapter~\ref{ch:nh1-rules}).


% =====================================================================
\chapter{Geographic Locality}
\label{ch:locality}
% =====================================================================

\newthought{Locality} in a peer-to-peer network is the property
that two nodes physically close on the Earth's surface can
reach each other with few hops and small per-hop latency. The
property is essential because real workloads are dominated by
local traffic. Kademlia does not have it; locality is its
single largest performance gap. \GDHT introduces it via a
geographic prefix in node identifiers, and \NDHT inherits the
construction unchanged. This chapter describes the construction,
its tradeoffs, and the way it interacts with the
learning mechanisms of Chapter~\ref{ch:operations}.

The principle, stated up front: \textbf{locality is a seed, not a
teacher.} The geographic prefix gives the routing structure
\emph{starting} affinity for short physical paths; the learning
mechanisms then \emph{strengthen} the affinity by reinforcing
shortcuts that traffic actually uses. Without the seed, learning
would have to discover regional structure from scratch (slow);
without the learning, the prefix alone would provide static
locality that doesn't respond to traffic (rigid). Together, they
produce a routing table that is regional from the first lookup
and increasingly traffic-shaped from then on.

% --------------------------------------------------------------------
\section{The S2 Cell Prefix}
% --------------------------------------------------------------------

\newthought{Google's S2}\sidenote{S2 documentation:
\url{https://s2geometry.io/}. The library is open source under
Apache 2.0; we use only the Hilbert-curve cell-ID computation,
which is a few hundred lines of arithmetic.} divides the
Earth's surface into a hierarchy of cells. The construction
projects the surface of a sphere onto the six faces of a
circumscribed cube, partitions each face along a Hilbert
space-filling curve, and assigns each cell a 64-bit integer
identifier called the \emph{cell ID}.

The Hilbert curve has a useful property: \emph{geographically
adjacent points have numerically close cell IDs}. Two cells
that share a 4-bit prefix are guaranteed to be on the same
quadrant of the same cube face; two cells that share an 8-bit
prefix are within roughly continental scale of each other.
Refining the prefix one bit further halves the cell area along
the curve.

\textbf{Top 8 bits.} \NDHT uses the top 8 bits of the S2 cell ID
as the geographic prefix. Three bits encode the cube face (six
faces, two unused encodings); five bits subdivide that face.
$6 \times 32 = 192$ tiles worldwide,\sidenote{Some tiles are
unused or sparsely populated (oceans, polar regions); the live
node distribution settles into roughly $80$--$120$ active tiles
on a uniform-land deployment.} each on the order of
continent-scale (e.g.\ ``western North America,''
``South-East Asia,'' ``northern Europe'').

\textbf{Sub-cell hierarchy.} Refining the prefix bit-by-bit
subdivides the tile in half along the Hilbert curve. At 16 bits
the cell is roughly state-sized. At 30 bits it is roughly
city-block-sized. At 40+ bits it is metres. \NDHT uses 8 bits
because that is the right granularity for routing decisions:
fine enough to distinguish continents, coarse enough that the
remaining 56 bits of node ID give plenty of intra-cell
diversity.\marginnote{Earlier experiments tried 4 and 16. Four
gave too few cells (only $24$ live tiles); routing decisions
within a tile were near-random. Sixteen gave too many tiny
cells (tens of thousands), most empty; routing within a tile
collapsed to a single peer. Eight is the empirical knee.}

% --------------------------------------------------------------------
\section{Node ID Construction}
% --------------------------------------------------------------------

\newthought{The node ID} is a 64-bit integer assembled from two
fields:

\[
\mathit{nodeId} \;=\;
\underbrace{\mathit{S2~cell~prefix}}_{\text{8 bits}}
\;\Vert\;
\underbrace{H_{56}(\mathit{publicKey})}_{\text{56 bits}}
\]

The high 8 bits come from the S2 cell containing the node's
claimed lat/lng coordinates (computed by the S2 library's
\texttt{cellIDFromLatLng()} routine, truncated to 8 bits). The
low 56 bits are a hash of the node's public key, providing
intra-cell uniqueness while keeping the ID bound to a stable
cryptographic identity.

\textbf{Why XOR distance approximates physical distance.} The
XOR metric Kademlia uses (Chapter~\ref{ch:lineage}) is bit-by-bit:
two IDs differ by an integer whose bit pattern reflects which bits
disagree. Two node IDs that share their top 8 bits XOR to a value
whose top 8 bits are zero --- a small integer in absolute terms.
Two node IDs that disagree in the top 8 bits XOR to a value whose
top 8 bits are nonzero --- a large integer.

Because the top 8 bits encode geographic position via the
Hilbert curve, ``share the top 8 bits'' translates to ``geographically
close,'' and ``differ in the top 8 bits'' translates to ``in
different continents.'' XOR distance, applied to S2-prefixed IDs,
\emph{approximately tracks} physical distance.

The approximation is loose at fine granularity (intra-cell, the
56-bit hash dominates and XOR distance is essentially random)
and tight at coarse granularity (cross-cell, XOR distance and
physical distance are very tightly correlated). This is exactly
the property a routing table wants: routes within a region need
all the intra-region diversity they can get; routes between
regions need to find each other quickly.

% --------------------------------------------------------------------
\section{Stratified Bootstrap}
% --------------------------------------------------------------------

\newthought{The bootstrap}\marginnote{Inherited from NX-5 with
minor refinements. The motivating observation: a uniformly random
synaptome at bootstrap fails to give good locality immediately,
because the geographic-prefix strata are statistically
under-represented. A pure-locality bootstrap fails the other way,
collapsing reachability when local traffic doesn't dominate.
Stratified bootstrap is the empirically-tuned middle ground.}
populates a new node's synaptome with $50$ initial entries
distributed across XOR strata such that:

\begin{itemize}\setlength\itemsep{0.15em}
  \item Every stratum from $0$ (closest to local) up to $63$
    (most distant) has at least one entry, so the routing-table
    connectivity guarantee is satisfied. This consumes $\le 64$
    entries (one per stratum), well within the $50$ budget for
    sparse strata.
  \item The geographic-prefix strata (top \textsc{geoBits}) are
    explicitly filled: the bootstrap pulls multiple entries from
    each of these strata up to a per-stratum cap, so that
    intra-region routing has multiple peers to choose from from
    the start.
  \item Remaining budget is allocated by sampling. The result is
    that the synaptome shape is \emph{biased} toward locality
    without being \emph{exclusive} to it.
\end{itemize}

The mechanism is deterministic given the bootstrap candidate
list, so we describe it as a budget-allocation algorithm:

\begin{algorithm}
\caption{Stratified bootstrap (sketch)}
\begin{algorithmic}[1]
\State $\mathit{budget} \gets 50$
\State $\mathit{geoStrata} \gets \{0, 1, \ldots, g - 1\}$
\Comment{$g$ = \textsc{geoBits}; the locality strata}
\State Allocate $k$ slots per geo stratum from $\mathit{budget}$
\State Allocate $1$ slot per non-geo stratum (skeleton)
\State Allocate remaining slots uniformly at random across all
strata
\State Materialize each slot by sampling a peer from the
candidate list at that stratum
\end{algorithmic}
\end{algorithm}

The candidate list is built by a sponsor-chain join (\S\ref{sec:structure}):
the new node routes a self-lookup through its sponsor; each node
on the path returns a row of its synaptome; the union becomes the
candidate list. This piggybacks on Pastry's locality-preserving
join (Chapter~\ref{ch:lineage}) to ensure that the candidates are
themselves RTT-close to the sponsor's region.

% --------------------------------------------------------------------
\section{Locality is a Seed, Not a Teacher}
% --------------------------------------------------------------------

\newthought{The S2 prefix} gives \NDHT \emph{starting} locality.
The learning mechanisms of Chapter~\ref{ch:operations} give it
\emph{evolving} locality. The interaction between the two is
the protocol's locality story.

\textbf{The seed.} The bootstrap synaptome is geographically
biased; AP scoring's latency penalty means that, even before any
LTP has fired, the protocol prefers nearby peers; XOR distance in
S2-prefixed ID space approximately tracks physical distance, so
$\log N$-bound iterative searches naturally converge on regional
peers in the common case. From the first lookup, regional traffic
already routes regionally.

\textbf{The teacher.} As traffic flows, LTP reinforces the
synapses that the actual workload uses. In a region where most
traffic is regional, LTP-strengthened synapses are nearby; in a
region with sustained cross-continental flows, LTP strengthens
some non-local synapses too. Hop caching installs targets that
were the destination of recent successful paths; lateral spread
propagates these to geographic neighbors. Over enough lookups,
the synaptome's content reflects the traffic pattern in the
node's region.

\textbf{The non-locality story.} The same protocol routes
non-local traffic perfectly well. Cross-continental lookups
follow the iterative-fallback path: first, route within the local
region toward a peer whose strata cross the geographic boundary;
then, the cross-boundary peer's synaptome (which carries some
peers in the destination region by virtue of its
under-represented-stratum bootstrap fill) provides the next hop.
The geometric-cost analysis ($3\delta$ floor,
Chapter~\ref{ch:eval-perf}) is unchanged by the locality bias;
locality is purely additive.

The combination is what gives \NDHT its scaling shape. Regional
latency is short and \emph{decreases} with $N$ as more nodes give
more shortcut opportunities. Global latency is bounded near the
$3\delta$ floor and grows logarithmically. Both regimes are
captured by the same routing protocol with no special-cased
modes.

% --------------------------------------------------------------------
\section{Cooperative Trust}
% --------------------------------------------------------------------

\newthought{The cost} of structural locality is a trust assumption.
The S2 prefix is \emph{self-declared}: a node can claim any
prefix it wants, regardless of where it actually is. The protocol
does not verify.

The benign failure mode is degraded routing. A node that lies
about its location ends up in a synaptome stratum that does not
match its actual physical position; routing through it costs more
latency than the synaptome's latency estimates would suggest. The
node still works as a routing primitive --- it can still receive
and forward messages --- but it does not contribute to locality.

The adversarial failure mode is a \emph{cell eclipse}: an attacker
generates many forged identities all claiming to land in the
same target cell, flooding that cell with attacker-controlled
peers. Subsequent honest peers in the cell are likely to bootstrap
from the attacker's nodes; the attacker can then drop traffic,
inject false peers, or redirect lookups within the cell.
Chapter~\ref{ch:redteam} discusses this attack class and the
mitigations.

\textbf{Mitigations.} Three non-exclusive directions are visible
in the literature:

\begin{itemize}\setlength\itemsep{0.15em}
  \item \emph{Verifiable RTT} via Vivaldi-style coordinates
    (Chapter~\ref{ch:lineage}): a peer's claimed location must be
    consistent with measured RTT to other peers. An attacker
    that misclaims must also fake RTTs, which is expensive at
    scale.
  \item \emph{Geographic proof-of-work}: bind the cell prefix to a
    cryptographic challenge that takes more time on a node not
    physically in the cell (e.g., a TURN-server-of-record
    handshake, an IP--ASN binding check).
  \item \emph{Trusted-witness schemes}: have a small number of
    well-known peers attest to the locations of other peers,
    similar to certificate-authority models.
\end{itemize}

The current protocol implements none of these; the prefix is
cooperative trust. \emph{The protocol does not depend on the
prefix being honest, but its locality benefits do.}
Chapter~\ref{ch:redteam} ranks the eclipse-attack risk as a
Tier-2 issue: real, but not blocking initial deployment, and
addressable through staged hardening.

\textbf{What this means for performance claims.} Every quantitative
claim in this document about regional latency assumes honest
prefixes. The simulator's nodes have correct prefixes by
construction. A real-world deployment with $5\%$ adversarial
prefixes would experience proportionally degraded locality, not
catastrophically degraded routing. The protocol's correctness is
unchanged; only the locality benefit is.

% --------------------------------------------------------------------
\section{Where Vivaldi Could Replace S2}
% --------------------------------------------------------------------

\newthought{A future direction}\sidenote{See \S\ref{ch:future}
for the full discussion.} is to replace the structural S2 prefix
with a \emph{learned} coordinate, in the spirit of Vivaldi
(Chapter~\ref{ch:lineage}). The trade is concrete: structural
locality is fast to bootstrap (no convergence period required,
prefix is known at addNode-time), trust-bound (the prefix is
self-declared), and parameter-free (no learning rate, no
coordinate dimension to choose). Learned locality is slow to
bootstrap (requires sustained traffic for convergence), trust-free
(coordinates can't be falsified without also falsifying RTTs),
and has a few tuning parameters.

For initial deployment, structural locality is the right choice:
it gives most of the benefit at a small fraction of the
implementation and operational cost. For mature deployment with
adversarial concerns, the migration path to learned locality is
clean --- the AP scoring's latency factor (\S\ref{sec:navigate})
already uses observed latency, so a Vivaldi-style coordinate
system would slot in under the existing scoring formula without
disturbing the rest of the protocol.


% =====================================================================
\chapter{Axonal Publish/Subscribe}
\label{ch:pubsub}
% =====================================================================

\begin{quote}
  \itshape Pub/sub belongs in the network, not on top of it.
\end{quote}

\newthought{Publish/subscribe} is the second primitive a real
peer-to-peer substrate must provide. A publisher sends a message
to a topic; every subscriber to that topic receives it. The
publisher does not know the subscribers; the subscribers do not
know each other; the network handles the fan-out.

Most contemporary peer-to-peer systems either layer pub/sub on
top of the routing primitive (Pastry's SCRIBE, IPFS's gossipsub)
or punt the problem to a central component (essentially every
production messaging app, however ``decentralized'' the
addressing layer). \NDHT builds pub/sub \emph{into} the protocol
via a structure called the \emph{axonal tree}, named for the
output branching of a biological neuron. The tree grows by
routed subscribe messages, heals continuously by routed
re-subscribe under churn, and delivers $100\%$ of messages at
baseline and $100\%$ of messages within three refresh intervals
under $5\%$ instantaneous churn (Chapter~\ref{ch:eval-perf}).

This chapter describes the construction. Chapter~\ref{ch:nh1-pubsub}
gives the per-tick pseudocode.

% --------------------------------------------------------------------
\section{Five Requirements}
% --------------------------------------------------------------------

\newthought{What pub/sub} must do, stated as five requirements
that any working implementation must meet.

\textbf{1. Topic addressing.} A publisher and a subscriber must
agree on a topic \emph{address} without coordinating. Both must
be able to derive the address from the topic name alone, with no
central registry.

\textbf{2. Subscriber discovery.} The publisher must be able to
reach all subscribers without knowing them individually. The
publisher's only inputs should be the topic address and the
message payload.

\textbf{3. Bounded fan-out.} A single node must not be required to
deliver to a number of peers proportional to the total subscriber
count of a popular topic. A topic with $10{,}000$ subscribers
cannot saturate the publisher's outbound bandwidth.

\textbf{4. Resilience to churn.} A subscriber that joins should
start receiving messages quickly. A node in the delivery path
that leaves should not silently break the topic for downstream
subscribers; recovery should be automatic.

\textbf{5. Replay for late joiners.} A subscriber that joins
\emph{slightly} after a publish should be able to receive the
recent past of the topic --- not the entire history, but enough
to fill the gap from missed messages.

The five requirements together are what lets pub/sub be
\emph{useful} as a real-time primitive: discovery is automatic,
fan-out is bounded, churn doesn't break anything, late joiners
recover.

% --------------------------------------------------------------------
\section{Topic Addressing: The Publisher Prefix}
% --------------------------------------------------------------------

\newthought{The address} of a topic is a 64-bit integer. \NDHT
constructs it as:

\[
\mathit{topicId} \;=\;
\underbrace{\mathit{publisher's~S2~prefix}}_{\text{8 bits}}
\;\Vert\;
\underbrace{H_{56}(\mathit{topicName})}_{\text{56 bits}}
\]

The high 8 bits come from the publisher's S2 cell prefix; the low
56 bits are a hash of the topic name. The address is deterministic:
given the topic name and the publisher's claimed location, both
publisher and subscribers can compute the same topic ID without
coordinating.

\textbf{Why publisher-prefix instead of plain hash.} If the topic
ID were a uniform hash of the topic name (NX-15 era), the topic's
delivery tree would root at a random cell on the planet. A chat
group of US-east subscribers might find its tree pinned to
central Asia. By embedding the publisher's S2 prefix in the
topic ID, the tree naturally roots near the publisher, and the
publisher's neighborhood (in S2) is also typically the
subscribers' neighborhood (real-world topic affinity is
geographically clustered). Inherited from NX-17.

\textbf{Naming convention.} Publishers and subscribers must agree
on the topic ID, which means agreeing on the publisher's prefix.
The convention we use is the topic-name form
\texttt{@CC/domain/event} where \texttt{CC} is the two-character
S2-cell shortcode for the publisher's region. The convention is
out-of-band; the protocol does not enforce it but treats topic
IDs as opaque 64-bit integers.

% --------------------------------------------------------------------
\section{Tree Construction}
% --------------------------------------------------------------------

\newthought{The axonal tree} is built by routed subscribe
messages.

\textbf{Subscribe.} A subscriber wishing to follow a topic
constructs the topic ID and routes a \texttt{pubsub:subscribe}
message toward it via the underlying \NDHT routing layer. The
message is processed at every intermediate node along the path:
each node checks whether it is already an axon for this topic.

\begin{itemize}\setlength\itemsep{0.15em}
  \item If the node is already an axon for the topic (it has
    open children for the topic), it accepts the new subscriber
    as one of its children and forwards no further. The subscribe
    has been intercepted; the subscriber is in the tree.
  \item If the node is not already an axon and the message has
    further to route, the node forwards. The message continues
    toward the topic ID.
  \item If the node is not already an axon and the message
    cannot route further (this is the topic's local root), the
    node opens a new role for the topic, becomes the
    \emph{axon root}, and accepts the subscriber as its first
    child.
\end{itemize}

\textbf{Publish.} A publisher constructs the same topic ID and
routes a \texttt{pubsub:publish} message toward it. The message
is intercepted by the same first-axon-along-the-path; from there,
it flows down the tree to all subscribers via the children-of-children
recursion.

\textbf{Branching on overflow.} When an axon's direct-child count
exceeds the configured \texttt{maxDirectSubs} threshold, the
overflow subscriber is delegated to a sub-axon (\emph{recruited}
from one of the existing children, or from an external synaptome
peer if the children are in the wrong direction for the new
subscriber). The mechanism keeps the maximum fan-out at any node
bounded by \texttt{maxDirectSubs}, regardless of total subscriber
count. A topic with $10{,}000$ subscribers might have $40$
direct-child slots filled at the root, and the rest of the
subscribers handed off recursively to depth-2 and depth-3
sub-axons.

% --------------------------------------------------------------------
\section{Self-Healing Under Churn}
% --------------------------------------------------------------------

\newthought{The mechanism} that heals the tree under churn is the
same routed re-subscribe that built it.

Every member of the tree --- subscribers and axons alike ---
periodically (every \texttt{refreshIntervalMs} milliseconds)
re-issues its subscribe message upstream. The re-subscribe is
exactly the same operation as the initial subscribe: a routed
message toward the topic ID.

\textbf{If the tree is intact}, the re-subscribe is intercepted
by the same axon as before. The tree state is unchanged; the
re-subscribe is a heartbeat.

\textbf{If a parent died}, the re-subscribe routes \emph{around}
the dead parent (the routing layer detects the dead peer and
takes an alternate hop, via the normal mechanism described in
Chapter~\ref{ch:operations}). The re-subscribe lands at whichever
live ancestor is now closest to the topic ID, which becomes the
subscriber's new parent. The orphan branch reattaches automatically.

\textbf{If the entire ancestor chain died}, the re-subscribe
routes all the way to the topic root. If the root is also dead,
the re-subscribe terminates at the closest-surviving node, which
opens a new role for the topic and becomes the new root. The
tree has been rebuilt from the re-subscriber's perspective.

The mechanism uses no special machinery. There are no heartbeats
beyond the re-subscribe itself; no failure detection separate
from the routing layer's dead-peer detection; no parent-tracking
state beyond the local children list. \emph{The same operation
that builds the tree repairs it.} This is the property that gives
\NDHT $100\%$ pub/sub delivery under $5\%$ churn at $25{,}000$
nodes (Chapter~\ref{ch:eval-perf}).

\textbf{TTL sweep.} As a complement to re-subscribe, axons
\emph{also} TTL-sweep their children: any child that has not
re-subscribed within \texttt{maxSubscriptionAgeMs} is dropped from
the children list. The TTL sweep cleans up after subscribers that
left without un-subscribing (e.g., the user closed the browser).

% --------------------------------------------------------------------
\section{The Replay Cache}
% --------------------------------------------------------------------

\newthought{Late-joining subscribers} are subscribers who joined
\emph{after} a publish. The naive behavior --- the late joiner
just misses the publish --- is acceptable for some workloads but
not for many.

\NDHT addresses this with a \emph{replay cache}: each axon node
maintains, for each topic it serves, a bounded ring buffer of
recently-published messages. The default size is $100$ messages.
When a new subscriber joins, the subscribe message can carry a
\texttt{lastSeenTs} parameter; the receiving axon then sends back
all messages in its replay cache more recent than that timestamp.
The new subscriber receives the publishes it missed during its
join transition.

\textbf{Bounded recovery.} The replay cache is a bounded
recovery, not infinite history. If a subscriber has been offline
for longer than the cache holds, the gap is irrecoverable. This
is acceptable for the operational target: ``a subscriber missing
a few seconds of publishes during a flaky reconnect should
recover'' is a much weaker requirement than ``a subscriber that
joins arbitrarily late should see all history'' --- the latter
needs a different system (a log, not a pub/sub bus).

\textbf{Per-topic sizing.} For high-rate topics, the cache size
in messages may correspond to a very short time window. The
cache can be configured to size by time rather than message
count, in which case the cache holds whatever fits in
\texttt{cacheLifetime} seconds. The protocol exposes both modes;
the default is message-count for predictable memory.

% --------------------------------------------------------------------
\section{Hot-Topic Migration (Future)}
% --------------------------------------------------------------------

\newthought{One known limitation}\sidenote{Discussed in
Chapter~\ref{ch:redteam} as a residual of the
bandwidth-saturation issue.} of the axonal-tree primitive as
described above is that a single popular topic with high publish
rate places its bandwidth cost on the topic root. The publisher
sends to the root; the root fans out to its
\texttt{maxDirectSubs} children; each child fans out further. The
root's bandwidth scales with the number of direct children, which
is bounded, but its compute cost (managing the children, sweeping
TTLs, processing re-subscribes) scales with the topic's churn rate
in subscribers. For very-high-rate Zipf-distributed pub/sub
workloads, this can saturate the root's CPU.

The mitigation is \emph{topic migration}: an axon root that
exceeds a load threshold can hand off the topic to a less-loaded
peer in the same S2 cell, leaving a redirect at the original
location. Subscribers re-subscribing to the original ID find the
redirect and follow it to the new root; the migration completes
when the redirect TTL expires. The mechanism is described in
detail in Chapter~\ref{ch:production} as part of the production
pathway. It is not implemented in the current \NDHT; the
foundational evidence (Chapter~\ref{ch:eval-bandwidth}) suggests
it is needed only for skewed pub/sub workloads, not for routing
load in general.

% --------------------------------------------------------------------
\section{Why ``Axonal''?}
% --------------------------------------------------------------------

\newthought{The naming} is exact. The output of a biological
neuron is the \emph{axon}: a long, branching cable that delivers
the neuron's spike output to many downstream targets. The branching
is structural --- new branches grow when the neuron's outputs need
to reach more cells; old branches retract when they fall into
disuse.

The axonal tree in \NDHT plays the same role in a peer-to-peer
network. A publishing node has an output (the topic message) that
must reach many downstream targets (the subscribers). The tree
branches are recruited as subscribers attach; sub-axons grow when
fan-out exceeds the local capacity; branches retract when
subscribers drop without renewing. The topology is not fixed at
construction; it evolves with the topic's load.

The translation, paralleling the table from Chapter~\ref{ch:substrate}:

\begin{table}[h]
\small
\begin{tabular}{@{}p{0.42\linewidth}p{0.52\linewidth}@{}}
\toprule
\textsc{In the brain} & \textsc{In N-DHT} \\
\midrule
Axon hillock (where the spike originates) & Publisher \\
Axon proper (the long delivering cable) & The route from publisher
  to topic root \\
Axonal arbor (the branching output structure) & The axon tree
  spanning the topic's subscribers \\
Synaptic terminal (where the axon meets a target dendrite) &
  An axon's child link to a subscriber \\
Branch growth on activity & Recruit-on-overflow \\
Branch retraction on disuse & TTL sweep of stale children \\
\bottomrule
\end{tabular}
\end{table}

The naming is not just decoration. The mechanism it describes is
structurally what biological axons do: deliver an output to many
targets through a branching cable that adapts its structure to
the demands of the workload it carries.


% #####################################################################
\part{NH-1: The Consolidated Generation}
% #####################################################################
% Generation-specific. The current production-ready implementation:
% which rules, which parameters, why each consolidation, full pseudocode.


% =====================================================================
\chapter{NH-1: Rules and Consolidation}
\label{ch:nh1-rules}
% =====================================================================

\begin{quote}
  \itshape Eighteen rules and forty-four parameters became twelve
  rules and twelve parameters. The behavior did not change.
\end{quote}

\newthought{Part~II} described the architecture in
generation-independent terms: the substrate, the five operations,
geographic locality, axonal pub/sub. This chapter describes the
specific implementation of that architecture in the current
generation, NH-1, and the consolidation work that produced it.

NH-1 is not a redesign of NX-17. It is the same protocol with
the rule set audited, deduplicated, and reorganized around the
vitality function. Every benchmark measurement in this volume
was produced by NH-1; every architectural claim in Part~II is a
description of what NH-1 does. The reason we treat the
generation as its own chapter is that the audit was the moment
when ``\NDHT'' became a coherent thing rather than a pile of
mechanisms each justified individually.

% --------------------------------------------------------------------
\section{The Twelve Parameters}
\label{sec:nh1-params}
% --------------------------------------------------------------------

NH-1 carries twelve tunable parameters. Each maps to one of the
five operations and controls a specific behavioral axis.
Table~\ref{tab:nh1-params} lists them in operational order.

\begin{table*}[h]
\centering
\small
\caption{The twelve tunable parameters of NH-1, with default
values, the operation each governs, and a one-line description
of what it controls.}
\label{tab:nh1-params}
\begin{tabular}{@{}lrlp{0.45\linewidth}@{}}
\toprule
\textbf{Parameter} & \textbf{Default} & \textbf{Operation} & \textbf{Controls} \\
\midrule
\texttt{maxSynaptome}      & $50$    & \textsc{structure} &
  Routing-table capacity (matches WebRTC connection budget) \\
\texttt{lookaheadAlpha}    & $5$     & \textsc{navigate}  &
  Two-hop probe breadth in AP scoring \\
\texttt{maxHops}           & $40$    & \textsc{navigate}  &
  Hard ceiling on lookup depth before failure \\
\texttt{epsilon}           & $0.05$  & \textsc{explore}   &
  First-hop random-choice probability \\
\texttt{annealCooling}     & $0.9997$ & \textsc{explore}  &
  Per-hop temperature decay multiplier \\
\texttt{annealRateScale}   & $1.0$   & \textsc{explore}   &
  Multiplier on per-hop annealing trigger probability \\
\texttt{decayGamma}        & $0.995$ & \textsc{forget}    &
  Per-decay-tick weight multiplier \\
\texttt{vitalityFloor}     & $0.05$  & \textsc{forget}    &
  Below this, a synapse is replaceable regardless of recency \\
\texttt{inertiaDuration}   & $20$    & \textsc{learn}     &
  LTP-lock window protecting freshly-reinforced synapses \\
\texttt{promoteThreshold}  & $2$     & \textsc{learn}     &
  Incoming-synapse use count needed to promote to first-class \\
\texttt{triadicThreshold}  & $2$     & \textsc{learn}     &
  Transit-pair count needed to fire triadic introduction \\
\texttt{recencyHalfLife}   & $50$    & \textsc{forget}    &
  Tick half-life of the post-inertia recency decay \\
\bottomrule
\end{tabular}
\end{table*}

The point of the tabulation is that the twelve parameters span
the entire behavioral envelope of the protocol. There is no
``hidden'' state that depends on a parameter not in this table.
Every other constant in the implementation
(\textsc{T\_INIT}, \textsc{T\_MIN}, \textsc{T\_REHEAT},
\textsc{DECAY\_INTERVAL}, \textsc{ANNEAL\_LOCAL\_SAMPLE},
\textsc{GEO\_BITS}, \textsc{STRATA\_GROUPS},
\textsc{LATERAL\_K}) is a fixed implementation choice not
intended for routine tuning; the values are documented in
Appendix~\ref{app:params}.

A reader who has internalized Chapter~\ref{ch:operations} can
read Table~\ref{tab:nh1-params} as a complete behavioral
specification: tell me the twelve numbers and I can predict the
protocol's behavior. That is the contribution of the consolidation.

% --------------------------------------------------------------------
\section{The NX-17 → NH-1 Mapping}
\label{sec:nh1-mapping}
% --------------------------------------------------------------------

NX-17 carried eighteen rules across forty-four parameters
(Appendix~\ref{app:nx-evolution} traces the lineage). The audit
that produced NH-1 categorized each NX-17 rule under one of three
dispositions: \emph{retained} (the rule survives, possibly with a
parameter rename or a small simplification), \emph{absorbed} (the
rule's function is performed by another rule plus the vitality
function, and the explicit rule can be retired without behavioral
loss), or \emph{retired} (the rule is removed; ablation showed it
was either harmful or ineffective at the scales we benchmark).

Table~\ref{tab:nh1-mapping} is the audit result.

\begin{table*}[h]
\centering
\small
\caption{The disposition of each NX-17 rule in NH-1. Twelve
retained, four absorbed into vitality, two retired. The behavioral
delta on the benchmark suite is within measurement noise; the
parameter count drops from forty-four to twelve.}
\label{tab:nh1-mapping}
\begin{tabular}{@{}p{0.32\linewidth}lp{0.42\linewidth}@{}}
\toprule
\textbf{NX-17 rule} & \textbf{Disposition} & \textbf{NH-1 status} \\
\midrule
1. AP scoring with two-hop lookahead & retained & Same mechanism;
  the latency penalty is the dominant scoring signal \\
2. LTP reinforcement on fast paths & retained & Same mechanism;
  applied only on paths under the EMA threshold \\
3. Hop caching on successful path & retained & Same mechanism;
  uses vitality-based admission \\
4. Lateral spread to geographic neighbors & retained & Same;
  $k = 2$ propagation depth \\
5. Triadic closure on transit pairs & retained & Same; threshold
  $= 2$ pairs \\
6. Iterative fallback on local minimum & retained & Same; takes
  closest unvisited regardless of XOR progress \\
7. Vitality-based admission gate & retained & The single equation
  governing all admissions \\
8. Vitality-based eviction with bootstrap preference & retained &
  Same hedge: bootstrap synapses preferred to non-bootstrap when
  vitality is tied \\
9. Periodic decay & retained & Decay gamma $= 0.995$ at intervals
  of $100$ lookups \\
10. Temperature annealing with reheat & retained & Reheat to
  $0.5$ on dead-peer detection \\
11. Stratified bootstrap & retained & Same construction; geographic
  strata explicitly filled \\
12. Sponsor-chain join & retained & Same; piggybacks on
  locality-preserving routing \\
\midrule
13. Two-tier (highway) synaptome & absorbed & The vitality-scored
  single tier subsumes both ``leaf set'' (high-vitality) and
  ``routing table'' (mid-vitality) roles. No need for an explicit
  highway tier \\
14. Stratum-floor enforcement & absorbed & Stratum diversity is
  preserved by the connection layer's \texttt{tryConnect} bilateral
  cap and the bootstrap's stratified fill. The explicit floor was
  redundant \\
15. Bespoke relay-pinning logic for pub/sub & absorbed & The standard
  axon-tree recruit-policy (synaptome-aware child selection) does
  the same work without a special-cased rule \\
16. \texttt{weightScale} parameter on AP scoring & absorbed &
  Ablation at $25{,}000$ nodes showed no measurable effect on
  routing or pub/sub. Removed; AP scoring is now pure
  progress/latency without LTP-weight bias \\
\midrule
17. Markov pre-learning of hot destinations & retired & The
  triadic-closure mechanism handles the same pattern at the
  network level. Removing Markov simplified the implementation
  without measurable performance loss \\
18. Load-tracked AP scoring (NX-12 experimental) & retired &
  Deferred to the production pathway (\S\ref{ch:production});
  not present in NH-1. The bandwidth-distribution analysis
  (Chapter~\ref{ch:eval-bandwidth}) showed it was not needed for
  the routing layer at our scales \\
\bottomrule
\end{tabular}
\end{table*}

\textbf{The two-tier collapse} is the largest single piece of the
consolidation. NX-15 and NX-17 carried both a synaptome (the
local-tier routing table, capacity $\sim 50$) and a separate
\emph{highway} tier (a small set of long-haul peers, capacity
$\sim 8$). The highway was managed by its own admission and
eviction rules, its own decay constants, and its own refresh
schedule. The audit observed that the highway-tier function ---
``a small number of long-haul peers maintained against churn''
--- is exactly what the high-vitality entries of the synaptome
do when traffic is dominated by long-haul lookups. The two-tier
machinery was producing a structural separation that the
underlying mechanism didn't actually need. Collapsing the two
tiers into one removed eight parameters and the entire
\texttt{\_refreshHighway} pathway.

\textbf{The stratum-floor absorption} is the second-largest. NX-17
carried an explicit rule that prevented any XOR stratum from
falling below a minimum population (default $2$ entries). The
rationale was reachability: if a stratum becomes empty, lookups
that need to traverse it have no candidate. The absorption
observation was that the connection-layer bilateral cap (every
candidate must pass \texttt{tryConnect}) plus the bootstrap's
stratified fill (every stratum gets at least one entry at
bootstrap) together provide the same guarantee. The explicit
floor rule was a safety net that the underlying mechanisms had
already made redundant.

\textbf{The Markov retirement} is the most interesting because it
documents a mechanism that worked but didn't justify its
complexity. NX-7 introduced \emph{Markov pre-learning}: each
node maintained a per-destination frequency table; destinations
above a frequency threshold received a pre-emptive synapse before
the first lookup actually used them. The mechanism produced
measurable improvement on hot-destination workloads, but only when
the same destinations were repeatedly queried within the Markov
window. The ablation showed that the same workload pattern was
captured by the triadic closure rule, with strictly less local
state. NX-17 still carried Markov as a separate rule for
backward compatibility; NH-1 retired it.

% --------------------------------------------------------------------
\section{What the Consolidation Did Not Do}
\label{sec:nh1-nondelta}
% --------------------------------------------------------------------

\newthought{The consolidation} did not change the protocol's
behavior. The benchmark suite at $25{,}000$ nodes shows NH-1 and
NX-17 within a few percent of each other on every metric:
$1.18\times$ vs $1.27\times$ the $3\delta$ floor on global
lookups; $32$\,ms vs $42$\,ms on $10\% \to 10\%$ concentrated
workload; $94.4\%$ vs $94.8\%$ Slice World success rate;
$100\%$ pub/sub baseline delivery for both. The full comparison
is in Chapter~\ref{ch:eval-perf}.

The point is that the audit was a \emph{simplification} not an
\emph{optimization}. NX-17 was already at the
state-of-the-art; the question was whether the eighteen rules
were each genuinely doing different work or whether the same
behavior could be expressed more economically. The answer was
the latter, by a margin of six rules and thirty-two parameters.

NH-1 is not faster than NX-17. NH-1 is the version of NX-17 that
\emph{can be reasoned about} --- twelve parameters spanning the
behavioral envelope, one equation governing every admission
decision, five operations under which every rule is classified.
That clarity is the contribution.

% --------------------------------------------------------------------
\section{Why NH-1, Specifically, for Production}
\label{sec:nh1-deploy-rationale}
% --------------------------------------------------------------------

\newthought{We selected NH-1} as the deployment target for two
reasons: \emph{maintainability} and \emph{tunability}.

\textbf{Maintainability.} An eighteen-rule, forty-four-parameter
protocol is feasible to implement once and very hard to maintain.
A twelve-rule, twelve-parameter protocol fits in a single
developer's head; the rule-to-mechanism map is small enough to
verify by inspection; ablation studies on a single parameter no
longer have to control for forty-three other interacting
parameters. The simulator and the rest of the codebase are easier
to reason about as a result. (NH-1's
\texttt{NeuromorphicDHTNH1.js} is approximately $1{,}500$ lines;
NX-17's was approximately $2{,}400$.)

\textbf{Tunability.} The bandwidth-distribution analysis in
Chapter~\ref{ch:eval-bandwidth} identified
\texttt{annealRateScale} as a single tunable knob with a clean
knee --- $5\times$ volume reduction at no routing-quality cost.
The same analysis was much harder to perform on NX-17 because
the equivalent ``annealing rate'' was determined by an
interaction between three different parameters
(\texttt{tInit}, \texttt{annealCooling}, and the
two-tier-specific \texttt{hubRefreshInterval}). Tuning one
without disturbing the others required a multi-dimensional sweep.
NH-1's twelve parameters are deliberately as orthogonal as we
could make them; sweeping a single axis produces an
interpretable result.

We continue to use NX-17 as the reference benchmark --- the bar
that NH-1 is asked to approach, and that future work should match
or surpass. The two protocols sit side-by-side in the test
harness; every benchmark in this volume reports both. The
eighteen-rule version remains as the empirical ceiling of the
``what can be done with this much machinery'' question; the
twelve-rule version is what we deploy.

% --------------------------------------------------------------------
\section{The State Each Node Carries}
\label{sec:nh1-state}
% --------------------------------------------------------------------

\newthought{For completeness}, here is the full per-node state of
an NH-1 node. Eight fields, all derived directly from the
operations described in Chapter~\ref{ch:operations}.

\begin{itemize}\setlength\itemsep{0.2em}
  \item \texttt{id} --- the 64-bit identifier
    (Chapter~\ref{ch:locality}).
  \item \texttt{lat, lng} --- claimed geographic coordinates,
    used only at bootstrap to compute the S2 prefix and at
    runtime for haversine latency estimation.
  \item \texttt{synaptome} --- a JavaScript \texttt{Map} keyed by
    \texttt{peerId}, holding live \texttt{Synapse} objects.
    Default capacity $50$.
  \item \texttt{incomingSynapses} --- a parallel \texttt{Map} of
    incoming-only synapses (peers that reached this node but that
    this node hasn't actively learned). Default capacity unbounded
    in practice; in \emph{de facto} bounded by the bilateral cap.
  \item \texttt{transitCache} --- a \texttt{Map} keyed by
    \texttt{(originId, nextId)} pairs, holding transit counts for
    triadic closure. Pruned on triadic firing.
  \item \texttt{temperature} --- a real in
    $[T_{\min}, T_{\mathrm{init}}]$, governing annealing
    probability.
  \item \texttt{simEpoch} --- the local count of lookups this
    node has participated in. Used as the time index for inertia
    and recency.
  \item \texttt{maxConnections} --- the bilateral cap, used by
    \texttt{tryConnect}.
\end{itemize}

That is everything an NH-1 node needs to operate. There is no
global state, no configuration that varies per node beyond the
identity and capacity fields, no specialized roles. Every node
runs the same code with the same parameters; specialization
emerges from the traffic pattern.


% =====================================================================
\chapter{The Full Routing Tick}
\label{ch:nh1-routing}
% =====================================================================

\newthought{This chapter} walks through the complete behavior of
a single lookup in NH-1, hop by hop, with each of the five
operations of Chapter~\ref{ch:operations} called out at the
moment it fires. The treatment is at the level of pseudocode:
detailed enough that a reader can implement it from this chapter
alone, abstract enough that the implementation language is not
assumed.

The reference implementation in JavaScript
(\texttt{src/dht/neuromorphic/NeuromorphicDHTNH1.js}) follows
this pseudocode line-for-line; cross-references in the margins
point to specific source-code symbols where useful.

% --------------------------------------------------------------------
\section{The Lookup Entry Point}
% --------------------------------------------------------------------

A lookup is invoked as \texttt{lookup(sourceId, targetKey)},
where \texttt{sourceId} is the originating node and
\texttt{targetKey} is the 64-bit key being sought. The lookup
returns a record \texttt{\{path, hops, time, found\}} where
\texttt{path} is the sequence of node IDs traversed, \texttt{hops}
is the path length minus one, \texttt{time} is the cumulative
latency, and \texttt{found} is true if the target was reached.

\begin{algorithm}
\caption{lookup(sourceId, targetKey) --- top level}
\begin{algorithmic}[1]
\State $\mathit{source} \gets \mathrm{nodeMap.get}(sourceId)$
\If{$\mathit{source} = \mathrm{null}$ \textbf{or}
    \textbf{not} $\mathit{source}.\mathit{alive}$}
  \State \textbf{return} $\mathrm{null}$
\EndIf
\State $\mathit{simEpoch} \mathrel{+}= 1$
\Comment{advance the local clock}
\If{$\mathit{lookupsSinceDecay} \geq D_{\mathrm{interval}}$}
  \Comment{$D_{\mathrm{interval}} =$ \textsc{DECAY\_INTERVAL}}
  \State $\mathrm{tickDecay}()$
  \Comment{periodic weight decay (FORGET)}
  \State $\mathit{lookupsSinceDecay} \gets 0$
\EndIf
\State $\mathit{path} \gets [sourceId]$;
       $\mathit{trace} \gets [\,]$;
       $\mathit{queried} \gets \{sourceId\}$
\State $\mathit{currentId} \gets sourceId$;
       $\mathit{totalTimeMs} \gets 0$;
       $\mathit{reached} \gets \mathrm{false}$
\For{$\mathit{hop} \gets 0$ \textbf{to} \texttt{maxHops} $- 1$}
  \State \emph{hop body --- see Algorithm 3}
\EndFor
\If{$\mathit{reached}$}
  \State \emph{post-lookup learning --- see Algorithm 4}
\EndIf
\State \textbf{return}
       $\{\mathit{path}, \mathit{hops}=|\mathit{path}|-1,
       \mathit{time}=\mathit{totalTimeMs}, \mathit{reached}\}$
\end{algorithmic}
\end{algorithm}

The simEpoch advance happens once per lookup. The per-decay-tick
weight multiplication happens every $100$ lookups (the
\textsc{DECAY\_INTERVAL} constant); the rationale is that decay
is cheap to amortize but doesn't need to fire every hop. The
trace structure --- a list of \texttt{(fromId, synapse)} pairs
--- is what the LTP reinforcement walk uses post-lookup
(\textsc{learn}); we accumulate it as we go.

% --------------------------------------------------------------------
\section{The Hop Body}
% --------------------------------------------------------------------

The body of the loop is where every operation other than periodic
decay actually fires. The structure is in five
phases.\sidenote{In the source, the phases are not always
presented in this exact order --- some interleave for performance
reasons --- but the pseudocode order is the conceptual one, and
the behavior is identical.}

\begin{algorithm}
\caption{hop body, run for each hop $\in [0, \texttt{maxHops})$}
\begin{algorithmic}[1]
\State $\mathit{current} \gets \mathrm{nodeMap.get}(\mathit{currentId})$
\If{$\mathit{current} = \mathrm{null}$ \textbf{or}
    \textbf{not} $\mathit{current}.\mathit{alive}$} \textbf{break}
\EndIf

\State $\mathit{currentDist} \gets \mathit{current}.\mathit{id} \oplus \mathit{targetKey}$
\If{$\mathit{currentDist} = 0$}
  $\mathit{reached} \gets \mathrm{true}$; \textbf{break}
\EndIf

\Statex \emph{Phase 1: \textsc{navigate} --- collect candidates}
\State $\mathit{candidates} \gets [\,]$;
       $\mathit{deadSynapses} \gets [\,]$
\For{\textbf{each} $s$ \textbf{in} $\mathit{current}.\mathit{synaptome}$}
  \If{$(s.\mathit{peerId} \oplus \mathit{targetKey}) \geq \mathit{currentDist}$}
    \textbf{continue} \Comment{no XOR progress}
  \EndIf
  \State $\mathit{peer} \gets \mathrm{nodeMap.get}(s.\mathit{peerId})$
  \If{$\mathit{peer}$ alive}
    $\mathit{candidates}.\mathrm{push}(s)$
  \Else
    $\mathit{deadSynapses}.\mathrm{push}(s)$;
    $s.\mathit{weight} \gets 0$
  \EndIf
\EndFor
\For{\textbf{each} $s$ \textbf{in} $\mathit{current}.\mathit{incomingSynapses}$}
  \If{$(s.\mathit{peerId} \oplus \mathit{targetKey}) < \mathit{currentDist}$
       \textbf{and} alive}
    $\mathit{candidates}.\mathrm{push}(s)$
  \EndIf
\EndFor

\Statex \emph{Phase 2: \textsc{forget}/\textsc{explore} ---
        dead-peer reheat \& evict-replace}
\If{$|\mathit{deadSynapses}| > 0$}
  \State $\mathit{current}.\mathit{temperature} \gets
         \max(\mathit{current}.\mathit{temperature}, T_{\mathrm{reheat}})$
  \For{\textbf{each} $\mathit{deadSyn}$}
    \State $\mathit{repl} \gets \mathrm{evictAndReplace}(\mathit{current},
           \mathit{deadSyn})$
    \If{$\mathit{repl} \neq \mathrm{null}$ \textbf{and} progress}
      \State $\mathit{candidates}.\mathrm{push}(\mathit{repl})$
    \EndIf
  \EndFor
\EndIf

\Statex \emph{Phase 3: \textsc{navigate} --- iterative fallback}
\If{$|\mathit{candidates}| = 0$}
  \State $\mathit{candidates} \gets [\mathrm{closestUnvisitedPeer}()]$
  \If{$|\mathit{candidates}| = 0$} \textbf{break} \EndIf
\EndIf

\Statex \emph{Phase 4: \textsc{navigate} --- select next hop}
\State $\mathit{nextSyn} \gets$ choose by direct-shortcut
       $\to$ epsilon-greedy $\to$ AP score with two-hop lookahead
\State $\mathit{nextId} \gets \mathit{nextSyn}.\mathit{peerId}$
\State $\mathit{nextNode} \gets \mathrm{nodeMap.get}(\mathit{nextId})$

\Statex \emph{Phase 5: \textsc{learn} \& \textsc{explore} ---
       on-the-fly updates}
\State \emph{(a) Incoming-synapse promotion if applicable
            (\textsc{learn})}
\State \emph{(b) Hop caching: install target in current
            (\textsc{learn})}
\State \emph{(c) Triadic-transit recording (\textsc{learn})}
\State \emph{(d) Anneal cooling \& probabilistic anneal
            (\textsc{explore})}
\State \emph{(e) Update path / trace / totalTimeMs}
\State $\mathit{currentId} \gets \mathit{nextId}$
\end{algorithmic}
\end{algorithm}

The hop body is deliberately structured so that each of the five
operations is identifiable as its own phase or sub-phase. A
reader who has internalized Chapter~\ref{ch:operations} can read
the algorithm as: \emph{at every hop, the protocol acts on
present input (Phase 4 chooses the hop), learns from the
choice (Phase 5a–c), forgets along the way
(Phase 2 evicts dead, periodic decay), and explores
occasionally (Phase 5d, plus the epsilon-greedy in Phase 4)}.
\textsc{structure} doesn't fire per hop; it fires at bootstrap
(\S\ref{sec:nh1-bootstrap-tick}).

% --------------------------------------------------------------------
\section{Selecting the Next Hop in Detail}
% --------------------------------------------------------------------

Phase 4 is the protocol's main act-on-input decision. The
selection runs through five mechanisms in priority order:

\begin{enumerate}\setlength\itemsep{0.15em}
  \item \emph{Direct shortcut.} If
    $\mathit{current}.\mathit{synaptome}$ contains a synapse for
    $\mathit{targetKey}$, use it. The lookup completes in one more
    hop.
  \item \emph{Epsilon-greedy first hop.} If $\mathit{hop} = 0$ and
    $\mathrm{Random}() < \texttt{epsilon}$, pick a candidate
    uniformly at random.
  \item \emph{AP scoring with two-hop lookahead.} For each of the
    top \texttt{lookaheadAlpha} candidates by single-hop AP score,
    extend the search one step into the candidate's known
    synaptome and compute a joint two-hop AP. Pick the candidate
    whose two-hop projection scores best.
  \item \emph{(Iterative fallback handled in Phase 3 ---
    invoked if Phase 1 produces no candidates.)}
  \item \emph{Termination.} If no candidate exists even after
    fallback, break the loop. The lookup returns
    $\mathit{found} = \mathrm{false}$.
\end{enumerate}

The two-hop AP formula is given in equation~\ref{eq:ap}; the
extension to two hops uses the candidate's known synaptome
entries (no extra round trip) to project the second-hop position.

% --------------------------------------------------------------------
\section{Phase 5 in Detail: On-the-Fly Updates}
% --------------------------------------------------------------------

The five sub-steps of Phase 5 are where most of the protocol's
\emph{adaptation} happens. We elaborate each.

\textbf{(a) Incoming-synapse promotion.} If $\mathit{nextSyn}$ is
an incoming synapse (came from
$\mathit{current}.\mathit{incomingSynapses}$), increment its
\texttt{useCount}. If the count crosses
\texttt{promoteThreshold}, promote it to a first-class synapse
in $\mathit{current}.\mathit{synaptome}$ via
$\mathrm{addByVitality}()$, weight $0.5$, fresh inertia.

\textbf{(b) Hop caching.} If $\mathit{currentId} \neq
\mathit{targetKey}$, the current node installs a synapse pointing
to $\mathit{targetKey}$ via $\mathrm{addByVitality}()$ with
weight $0.5$ and fresh inertia. Lateral spread also fires here:
the top \texttt{LATERAL\_K} geographically-nearest peers in
$\mathit{current}.\mathit{synaptome}$ each get a recursive
hop-cache invocation (depth-1 only; no further recursion).

\textbf{(c) Triadic-transit recording.} If $\mathit{currentId}
\neq \mathit{sourceId}$, increment the transit count for the
pair $(\mathit{sourceId}, \mathit{nextId})$ in
$\mathit{current}.\mathit{transitCache}$. If the count crosses
\texttt{triadicThreshold}, fire the introduction:
$\mathrm{introduce}(\mathit{sourceId}, \mathit{nextId})$ ---
which adds a synapse from \emph{source} to \emph{next}, weight
$0.5$, via $\mathrm{addByVitality}()$ at the source's synaptome.

\textbf{(d) Anneal cooling and probabilistic anneal.} Multiply
$\mathit{current}.\mathit{temperature}$ by \texttt{annealCooling},
floored at $T_{\min}$. With probability
$\mathit{current}.\mathit{temperature} \cdot \texttt{annealRateScale}$,
invoke $\mathrm{tryAnneal}(\mathit{current})$. The inner mechanism
of $\mathrm{tryAnneal}$ is described below.

\textbf{(e) Path / trace / time updates.} Append $\mathit{nextId}$
to $\mathit{path}$. Append $(\mathit{currentId}, \mathit{nextSyn})$
to $\mathit{trace}$. Increment $\mathit{totalTimeMs}$ by
$\mathit{nextSyn}.\mathit{latency}$. Mark $\mathit{nextId}$ as
queried.

% --------------------------------------------------------------------
\section{The Vitality-Based Admission Gate}
% --------------------------------------------------------------------

Many of the operations above call $\mathrm{addByVitality}()$.
This is the single function that admits new synapses to a
synaptome under the vitality rule.

\begin{algorithm}
\caption{addByVitality(node, newSyn)}
\begin{algorithmic}[1]
\State $\mathit{peer} \gets \mathrm{nodeMap.get}(\mathit{newSyn}.\mathit{peerId})$
\If{$\mathit{peer} \neq \mathrm{null}$ \textbf{and not}
    $\mathit{node}.\mathrm{tryConnect}(\mathit{peer})$}
  \State \textbf{return} false
  \Comment{bilateral cap refused; \textsc{structure}}
\EndIf
\If{$|\mathit{node}.\mathit{synaptome}| < \texttt{maxSynaptome}$}
  \State $\mathit{node}.\mathrm{addSynapse}(\mathit{newSyn})$
  \State \textbf{return} true
\EndIf

\Statex \emph{synaptome at capacity --- vitality-based eviction}
\State $\mathit{victim} \gets \mathrm{null}$;
       $\mathit{minV} \gets \infty$
\State $\mathit{victimAny} \gets \mathrm{null}$;
       $\mathit{minVAny} \gets \infty$
\For{\textbf{each} $s$ \textbf{in} $\mathit{node}.\mathit{synaptome}$}
  \If{$s.\mathit{inertia} > \mathit{simEpoch}$} \textbf{continue}
  \Comment{LTP-locked; protected}
  \EndIf
  \State $v \gets \mathrm{vitality}(\mathit{node}, s)$
  \If{$v < \mathit{minVAny}$}
    $\mathit{minVAny} \gets v$;
    $\mathit{victimAny} \gets s$
  \EndIf
  \If{\textbf{not} $s.\mathit{bootstrap}$ \textbf{and}
      $v < \mathit{minV}$}
    $\mathit{minV} \gets v$;
    $\mathit{victim} \gets s$
  \EndIf
\EndFor
\State $\mathit{evictTarget} \gets \mathit{victim} \neq \mathrm{null}\;?
       \;\mathit{victim} : \mathit{victimAny}$
\If{$\mathit{evictTarget} = \mathrm{null}$} \textbf{return} false
\EndIf
\State $\mathit{node}.\mathit{synaptome}.\mathrm{delete}(\mathit{evictTarget}.\mathit{peerId})$
\State $\mathit{node}.\mathrm{addSynapse}(\mathit{newSyn})$
\State \textbf{return} true
\end{algorithmic}
\end{algorithm}

The bootstrap-preference hedge (line 12 of Algorithm 5) is
worth dwelling on. Bootstrap synapses are the structural backbone
of the routing table; evicting them too aggressively breaks the
stratum-coverage guarantee that the bootstrap process worked to
establish. The hedge says: prefer to evict a non-bootstrap
synapse if one is available. Only if every candidate is a
bootstrap entry does the gate fall back to evicting the
lowest-vitality bootstrap entry. In practice this rarely fires;
new entries from triadic / hop-cache / lateral-spread quickly
populate the synaptome with non-bootstrap material, and bootstrap
entries become the longstanding anchors that never get evicted.

% --------------------------------------------------------------------
\section{The Anneal Step}
% --------------------------------------------------------------------

The anneal step is the protocol's main exploration mechanism. It
fires probabilistically, governed by the local temperature and
the \texttt{annealRateScale} master volume control.

\begin{algorithm}
\caption{tryAnneal(node)}
\begin{algorithmic}[1]
\State \emph{Find weakest non-LTP-locked synapse in node.synaptome}
\State $\mathit{victim} \gets \mathrm{null}$;
       $\mathit{minW} \gets \infty$
\For{\textbf{each} $s$ \textbf{in} $\mathit{node}.\mathit{synaptome}$}
  \If{$s.\mathit{inertia} > \mathit{simEpoch}$} \textbf{continue}
  \EndIf
  \If{$s.\mathit{weight} < \mathit{minW}$}
    $\mathit{minW} \gets s.\mathit{weight}$;
    $\mathit{victim} \gets s$
  \EndIf
\EndFor
\If{$\mathit{victim} = \mathrm{null}$} \textbf{return} \EndIf

\State \emph{Identify under-represented stratum group}
\State (count strata occupancy; pick the lo-hi range with min count)
\State $\mathit{candidate} \gets \mathrm{localCandidate}(\mathit{node},
       \mathit{lo}, \mathit{hi})$
\If{$\mathit{candidate} = \mathrm{null}$} \textbf{return} \EndIf

\State \emph{Atomic 1-for-1 exchange under bilateral cap}
\State $\mathit{node}.\mathit{synaptome}.\mathrm{delete}(\mathit{victim}.\mathit{peerId})$
\State $\mathit{node}.\mathit{connections}.\mathrm{delete}(\mathit{victim}.\mathit{peerId})$
\If{\textbf{not} $\mathit{node}.\mathrm{tryConnect}(\mathit{candidate})$}
  \textbf{return}
  \Comment{candidate's cap refused; abort}
\EndIf
\State $\mathit{node}.\mathrm{addSynapse}(\mathrm{Synapse}(\mathit{candidate},
       \mathrm{weight}=0.1))$
\end{algorithmic}
\end{algorithm}

The 1-for-1 exchange is intentional: anneal does not grow the
synaptome, it \emph{rotates} an entry. The peer count stays
constant; the under-represented-stratum bias gives the rotation
its diversifying effect over time. The starting weight of $0.1$
is low enough that the new entry is itself a candidate for
eviction by the next anneal --- only entries that the actual
traffic reinforces accumulate weight beyond the floor.

The \texttt{localCandidate}() function scans the 2-hop neighborhood
(the union of $\mathit{node}$'s synaptome peers' synaptomes) for a
peer in the target stratum range. It samples up to
\textsc{ANNEAL\_LOCAL\_SAMPLE} candidates and returns one
uniformly. This is the \texttt{local\_probe} traffic source
that dominates the bandwidth analysis in
Chapter~\ref{ch:eval-bandwidth}.

% --------------------------------------------------------------------
\section{Post-Lookup Learning}
% --------------------------------------------------------------------

If the lookup reached the target, NH-1 walks the trace and applies
LTP to every synapse on the path that beat the running EMA of
recent path latencies.

\begin{algorithm}
\caption{post-lookup learning, fires only if reached}
\begin{algorithmic}[1]
\State $\mathit{hopCount} \gets |\mathit{path}| - 1$
\State \emph{Update EMA of recent path latencies}
\State $\mathit{emaTime} \gets 0.9 \cdot \mathit{emaTime} +
       0.1 \cdot \mathit{totalTimeMs}$
\If{$\mathit{totalTimeMs} \leq \mathit{emaTime}$}
  \For{\textbf{each} $(\mathit{fromId}, \mathit{synapse})$
       \textbf{in} $\mathit{trace}$ from end to start}
    \State $\mathit{node} \gets \mathrm{nodeMap.get}(\mathit{fromId})$
    \State $\mathit{syn} \gets
           \mathit{node}.\mathit{synaptome}.\mathrm{get}(\mathit{synapse}.\mathit{peerId})$
    \If{$\mathit{syn} \neq \mathrm{null}$}
      \State $\mathit{syn}.\mathit{weight} \gets
             \min(1.0, \mathit{syn}.\mathit{weight} + 0.05)$
      \State $\mathit{syn}.\mathit{inertia} \gets
             \mathit{simEpoch}$
      \Comment{enter LTP-lock window}
      \State $\mathit{syn}.\mathit{useCount} \mathrel{+}= 1$
    \EndIf
  \EndFor
\EndIf
\end{algorithmic}
\end{algorithm}

The trace is walked end-to-start so that the synaptic-tagging
mechanism behaves correctly: the more recent reinforcements
(toward the destination) fire first, establishing their
LTP-lock, before earlier ones. The EMA threshold ensures that
LTP fires only on paths that beat the running average; if every
successful path triggered reinforcement, slow paths would fight
fast paths in the synaptome, and the learning would not converge
on the fast routes.

% --------------------------------------------------------------------
\section{The Bootstrap Tick}
\label{sec:nh1-bootstrap-tick}
% --------------------------------------------------------------------

A new node's bootstrap is the only time
\textsc{structure} operations dominate. The sequence:

\begin{enumerate}\setlength\itemsep{0.15em}
  \item Compute the node's S2 cell from its claimed lat/lng.
  \item Construct the node ID:
    $\mathit{nodeId} = \mathit{S2 prefix}\;\Vert\;
    H_{56}(\mathit{publicKey})$.
  \item Pick a sponsor (a known live peer; in production this is
    typically a published rendezvous list).
  \item Route a self-lookup through the sponsor toward
    $\mathit{nodeId}$. Each node along the path returns a row of
    its synaptome.
  \item Run the stratified-bootstrap allocator
    (\S\ref{ch:locality}) over the union of returned candidates.
    Allocate $50$ slots: $k$ per geographic stratum, $1$ per
    non-geographic stratum, remainder uniformly random.
  \item Materialize each slot by calling
    $\mathrm{addByVitality}()$ with weight $0.5$, fresh inertia,
    and the bootstrap flag set.
  \item Propagate the bidirectional handshake: each peer added
    learns of the new node as an incoming synapse.
\end{enumerate}

After this completes, the node is operational. Its very first
lookup will route through a synaptome shaped by its sponsor's
locality, with explicit fill on the geographic strata so that
intra-region traffic finds nearby peers immediately. The
learning mechanisms then take over on subsequent traffic.

The bootstrap pseudocode is more involved than the steady-state
hop body; the full algorithm is in Appendix~\ref{app:params}.


% =====================================================================
\chapter{The Full Pub/Sub Tick}
\label{ch:nh1-pubsub}
% =====================================================================

\newthought{The pub/sub layer} sits on top of the routing layer
of Chapter~\ref{ch:nh1-routing} and uses two new primitives:
\texttt{routeMessage} (route a typed message toward a topic ID,
intercepting at the first axon along the path) and
\texttt{sendDirect} (deliver a message to a specific known peer).
The axonal tree's subscribe / publish / refresh behaviors are
implemented entirely in terms of these two primitives plus a
per-node \emph{role} table that tracks topic membership.

% --------------------------------------------------------------------
\section{Per-Node Pub/Sub State}
% --------------------------------------------------------------------

NH-1 augments the per-node state of \S\ref{sec:nh1-state} with the
following pub/sub-specific fields. Each is initialized empty when
the node joins.

\begin{itemize}\setlength\itemsep{0.15em}
  \item \texttt{routedHandlers} --- a map from message type to
    handler function for routed messages received at this node.
    AxonManager populates this at axonFor() time.
  \item \texttt{directHandlers} --- the same for direct messages.
  \item \texttt{axonsByNode} --- a map from \texttt{NeuronNode}
    to its \texttt{AxonManager} instance. Lazy: an entry is
    created the first time this node calls \texttt{axonFor()}.
  \item Per topic, in the AxonManager: a \emph{role} record
    holding \texttt{children} (a map of subscriber IDs to
    metadata), \texttt{lastSeenTs} (the timestamp of the most
    recent publish), and the \texttt{replayCache} ring buffer.
\end{itemize}

The AxonManager is the per-node object that owns all pub/sub
state for that node. A single node can serve as an axon for many
topics simultaneously; each topic has its own role record inside
the same AxonManager.

% --------------------------------------------------------------------
\section{Topic ID Construction}
% --------------------------------------------------------------------

The topic ID is constructed from the topic name and the
publisher's S2 prefix, as described in
Chapter~\ref{ch:pubsub}. The \texttt{usesPublisherPrefix} flag on
the protocol activates the \texttt{topicIdForPrefixed} variant of
the construction: the topic name carries an explicit
\texttt{@CC/} prefix, and the protocol uses the first eight bits
of the publisher's claimed cell as the high bits of the topic ID.

\begin{algorithm}
\caption{topicId = topicIdForPrefixed(name)}
\begin{algorithmic}[1]
\State \emph{Parse \texttt{@CC/domain/event} into the cell shortcode and the rest}
\State $\mathit{cellPrefix} \gets \mathrm{S2.cellFromShortcode}(\mathit{CC})$
\State $\mathit{nameHash} \gets H_{56}(\mathrm{rest of name})$
\State \textbf{return} $(\mathit{cellPrefix} \ll 56) \mathbin{|}
       \mathit{nameHash}$
\end{algorithmic}
\end{algorithm}

The construction is deterministic. Both publisher and subscriber
compute the same topic ID without coordination, given only the
topic name string. The publisher's S2 prefix encoded in the high
bits causes routed subscribes to converge on a node geographically
close to the publisher --- typically also close to the subscribers,
since real-world topic memberships are geographically clustered.

% --------------------------------------------------------------------
\section{Subscribe}
% --------------------------------------------------------------------

A subscribe operation routes a typed message
\texttt{pubsub:subscribe} toward the topic ID. The
\texttt{routeMessage} primitive does the actual routing; an
intermediate node intercepts when it is already an axon for the
topic.

\begin{algorithm}
\caption{subscribe(node, topicId, subscriberMeta)}
\begin{algorithmic}[1]
\State $\mathit{result} \gets
       \mathrm{routeMessage}(\mathit{node},
       \mathit{topicId},
       \texttt{`pubsub:subscribe'},
       \mathit{subscriberMeta})$
\If{$\mathit{result}.\mathit{consumed}$}
  \State \emph{some axon along the path accepted; subscriber is
        in the tree}
\ElsIf{$\mathit{result}.\mathit{terminal}$}
  \State \emph{the routed walk reached a local terminal with no
        existing axon for this topic --- the terminal node
        becomes the new root}
  \State $\mathit{terminalNode} \gets
        \mathrm{nodeMap.get}(\mathit{result}.\mathit{atNode})$
  \State $\mathit{terminalNode}.\mathit{axonFor}(\mathit{topicId})
        .\mathrm{openRole}()$
  \State $\mathit{terminalNode}.\mathit{axonFor}(\mathit{topicId})
        .\mathrm{addChild}(\mathit{subscriberMeta})$
\Else
  \State \emph{exhausted --- failure mode; subscriber retries}
\EndIf
\end{algorithmic}
\end{algorithm}

\textbf{The interception logic.} When a routed message visits an
intermediate node, \texttt{routeMessage} delivers the message to
the node's registered routed-message handler for the message
type. The handler returns either \texttt{`forward'} (continue
routing) or \texttt{`consumed'} (stop here). The
\texttt{pubsub:subscribe} handler checks: \emph{is this node
already an axon for this topic?} If yes, it adds the subscriber
to its children, returns \texttt{`consumed'}, and routing
terminates. If no, it returns \texttt{`forward'} and the message
continues toward the topic ID.

The handler is registered via \texttt{onRoutedMessage} when the
node's AxonManager is initialized.

% --------------------------------------------------------------------
\section{Publish}
% --------------------------------------------------------------------

A publish operation routes a \texttt{pubsub:publish} message
toward the topic ID. The first axon along the path intercepts and
fans out.

\begin{algorithm}
\caption{publish(node, topicId, payload)}
\begin{algorithmic}[1]
\State $\mathit{result} \gets
        \mathrm{routeMessage}(\mathit{node},
        \mathit{topicId},
        \texttt{`pubsub:publish'},
        \mathit{payload})$
\Comment{first axon intercepts; fan-out happens inside the handler}
\end{algorithmic}
\end{algorithm}

\textbf{The fan-out logic, inside the handler at the intercepting axon:}

\begin{algorithm}
\caption{publish handler at an axon}
\begin{algorithmic}[1]
\State $\mathit{role} \gets \mathit{this}.\mathit{rolesFor}(\mathit{topicId})$
\State $\mathit{role}.\mathit{replayCache}.\mathrm{push}(\mathit{payload})$
\State $\mathit{role}.\mathit{lastSeenTs} \gets \mathrm{now}()$
\For{\textbf{each} $(\mathit{childId}, \mathit{childMeta})$ \textbf{in}
     $\mathit{role}.\mathit{children}$}
  \State $\mathrm{sendDirect}(\mathit{this},
         \mathit{childId},
         \texttt{`pubsub:deliver'},
         \mathit{payload})$
\EndFor
\State \textbf{return} \texttt{`consumed'}
\end{algorithmic}
\end{algorithm}

The fan-out uses \texttt{sendDirect} (point-to-point), not
\texttt{routeMessage} (routed). The axon already knows the
children's identities, so direct delivery is correct.

\textbf{Children that are sub-axons recurse.} If a child is
itself an axon (registered to handle the topic), its
\texttt{pubsub:deliver} handler does the same fan-out as above.
The result is a tree-structured fan-out where each level handles
up to \texttt{maxDirectSubs} children before delegating further.

% --------------------------------------------------------------------
\section{Refresh: The Self-Healing Mechanism}
% --------------------------------------------------------------------

The refresh tick fires on every member of the tree (subscribers
and axons alike) every \texttt{refreshIntervalMs} milliseconds. It
is the mechanism that heals the tree under churn.

\begin{algorithm}
\caption{refreshTick(node)}
\begin{algorithmic}[1]
\For{\textbf{each} topic this node is subscribed to or serves
     as axon for}
  \State \emph{Send a re-subscribe to keep the upstream link
       alive}
  \State $\mathrm{subscribe}(\mathit{node}, \mathit{topicId},
        \mathit{this node's meta})$
  \State \emph{TTL-sweep our own children}
  \For{\textbf{each} child $c$ in our role for $\mathit{topicId}$}
    \If{$\mathrm{now}() - c.\mathit{lastRefreshTs} >
        \texttt{maxSubscriptionAgeMs}$}
      \State drop $c$ from the children list
    \EndIf
  \EndFor
\EndFor
\end{algorithmic}
\end{algorithm}

\textbf{Why this is self-healing.} The re-subscribe is exactly
the same routed message that built the tree initially. If the
upstream parent is alive, the re-subscribe is intercepted by the
parent, which already has this child in its children list ---
the child's \texttt{lastRefreshTs} is updated, no other state
changes. If the parent is \emph{dead}, the re-subscribe routes
\emph{around} the dead parent (the routing layer's dead-peer
detection takes care of this), lands at whichever live ancestor
is now closest to the topic ID, and that ancestor accepts the
re-subscribe. The orphaned child is now connected to its new
parent.

The TTL sweep is the corresponding garbage-collection: a child
that has not refreshed within \texttt{maxSubscriptionAgeMs} is
assumed to have left without unsubscribing, and is dropped.

There is no separate failure-detection machinery. There are no
heartbeats independent of the re-subscribe. There is no parent-pointer
tracking that could become inconsistent. \emph{The same operation
that builds the tree repairs it.} This is the architectural
invariant that produces $100\%$ pub/sub recovery under $5\%$
instantaneous churn.

% --------------------------------------------------------------------
\section{Branching on Overflow: Recruit and Relay}
% --------------------------------------------------------------------

When an axon's direct-child count exceeds
\texttt{maxDirectSubs}, the overflow subscriber is delegated to a
\emph{sub-axon}. NH-1 inherits two policies for picking the
sub-axon, both implemented as callbacks on the AxonManager.

\textbf{Recruit policy (existing-child promotion).} If at least
one of the existing children is a synaptome peer of this axon
with high LTP weight, promote it to a sub-axon and have it adopt
the overflow subscriber. The high-vitality synaptome peer is the
most reliable choice for sustained delivery. Implementation:

\begin{algorithm}
\caption{pickRecruitPeer(role, meta, subscriberId)}
\begin{algorithmic}[1]
\For{\textbf{each} child $c$ in $\mathit{role}.\mathit{children}$}
  \If{$c$ is in our synaptome \textbf{and}
      $c$ has high LTP weight}
    \State \textbf{return} $c$
    \Comment{recruit existing child as sub-axon}
  \EndIf
\EndFor
\State \textbf{return} XOR-closest existing child to subscriber's ID
\end{algorithmic}
\end{algorithm}

\textbf{Relay policy (external promotion).} If no existing child
is suitable for promotion, pick an external synaptome peer
(XOR-closest to the new subscriber's ID, not yet in the
children) and recruit it. The recruited peer may then adopt
multiple existing children whose IDs point in the same direction.

\begin{algorithm}
\caption{pickRelayPeer(role, subscriberId, forwarderId)}
\begin{algorithmic}[1]
\State $\mathit{best} \gets \mathrm{null}$;
       $\mathit{bestDist} \gets \infty$
\For{\textbf{each} synapse $s$ in this node's synaptome}
  \If{$s.\mathit{peerId}$ already in role's children
       \textbf{or} $= \mathit{forwarderId}$
       \textbf{or} $= \mathit{subscriberId}$}
    \textbf{continue}
  \EndIf
  \State $d \gets s.\mathit{peerId} \oplus
        \mathrm{topicToBigInt}(\mathit{subscriberId})$
  \If{$d < \mathit{bestDist}$}
    $\mathit{bestDist} \gets d$;
    $\mathit{best} \gets s.\mathit{peerId}$
  \EndIf
\EndFor
\State \textbf{return} $\mathit{best}$
\end{algorithmic}
\end{algorithm}

The two policies together keep the tree's branching factor
bounded at any single node by \texttt{maxDirectSubs}, with the
new sub-axon chosen so that geographic / topological clustering
of subscribers ends up handled by the same sub-axon. Inherited
from NX-15.

% --------------------------------------------------------------------
\section{The Replay Cache}
% --------------------------------------------------------------------

Each axon's role for each topic carries a bounded ring buffer of
recently-published messages. Default capacity is $100$ messages.
The buffer is updated on every publish (line 2 of the publish
handler above) and consulted on every subscribe whose
\texttt{lastSeenTs} is set.

\begin{algorithm}
\caption{subscribe handler with replay request}
\begin{algorithmic}[1]
\State \emph{Standard subscribe handling --- add child, etc.}
\If{$\mathit{subscriberMeta}$ contains a \texttt{lastSeenTs}}
  \For{\textbf{each} cached message $m$ in $\mathit{role}.\mathit{replayCache}$}
    \If{$m.\mathit{ts} > \mathit{subscriberMeta}.\mathit{lastSeenTs}$}
      \State $\mathrm{sendDirect}(\mathit{this},
             \mathit{subscriberId},
             \texttt{`pubsub:replay'},
             m)$
    \EndIf
  \EndFor
\EndIf
\State \textbf{return} \texttt{`consumed'}
\end{algorithmic}
\end{algorithm}

The mechanism gives a late-joining subscriber a bounded recovery
of the recent past of the topic. ``Bounded'' is the operative
word: the ring buffer holds the last $100$ messages (or whatever
fits in \texttt{cacheLifetime} seconds), not the entire
history. A subscriber that has been offline for longer than the
cache holds is told nothing about the lost interval --- the gap
is irrecoverable from the pub/sub layer alone.\sidenote{This is
intentional. Pub/sub is a real-time message bus, not a log. If
clients need durable history they should use a different
primitive layered on top --- one that either pulls from a
log-backed source on resume, or carries a per-message ack
protocol that reports gaps.}

% --------------------------------------------------------------------
\section{The Two Underlying Primitives}
% --------------------------------------------------------------------

The entire chapter to this point has used two primitives:
\texttt{routeMessage} and \texttt{sendDirect}. We give the
sketches here for completeness; the full implementations are in
\texttt{NeuromorphicDHTNH1.js}.

\textbf{routeMessage(originNode, targetId, type, payload).} Walks
a typed message hop-by-hop toward \texttt{targetId} using greedy
single-hop routing (no two-hop lookahead, since pub/sub doesn't
benefit from it). At each hop, the node's registered routed-message
handler for \texttt{type} can intercept (\texttt{`consumed'}) or
forward (\texttt{`forward'}). On reaching a local terminal (no
hop offers further XOR progress), runs a 2-hop terminal-globality
check\sidenote{The terminal-globality check is necessary because
greedy routing can converge on different local minima from
different sources, fragmenting the tree. The 2-hop scan asks: ``is
there \emph{any} peer reachable in two hops that is closer to the
target than I am?'' If yes, take that path. If no, this is a
\emph{global} terminal, and an axon role can be opened here.} ---
if a 2-hop peer is closer to the target than the current node,
take it; otherwise this node is the global terminal. Returns
\texttt{\{consumed, atNode, hops, terminal?, exhausted?\}}.

\textbf{sendDirect(fromNode, peerId, type, payload).} Delivers
\texttt{payload} directly to a known peer's registered
direct-message handler. Used for fan-out from axons to children.
NH-1 implements an iterative drain queue under
\texttt{sendDirect}: a handler triggered by one
\texttt{sendDirect} can itself call \texttt{sendDirect}, and so
on; rather than building up a recursive call stack (which blows
past Node's $\sim 10\mathrm{K}$ stack limit on deep trees), the
nested calls enqueue items in a FIFO drain that the outer
\texttt{sendDirect} processes iteratively. From the caller's
perspective, all handlers triggered by a top-level publish
complete before \texttt{sendDirect} returns; the drain queue is
an implementation detail.

% --------------------------------------------------------------------
\section{Operational Tunables}
% --------------------------------------------------------------------

The pub/sub layer carries five operational parameters in addition
to the twelve routing parameters of \S\ref{sec:nh1-params}:

\begin{table}[h]
\small
\begin{tabular}{@{}lrl@{}}
\toprule
\textbf{Parameter} & \textbf{Default} & \textbf{Controls} \\
\midrule
\texttt{maxDirectSubs} & $32$ & Max direct children per axon
  before recruit-on-overflow \\
\texttt{minDirectSubs} & $4$ & Min direct children per axon
  before relinquishing role \\
\texttt{refreshIntervalMs} & $30000$ & Per-member refresh
  interval (re-subscribe) \\
\texttt{maxSubscriptionAgeMs} & $90000$ & TTL for child entries
  before sweep \\
\texttt{replayCacheSize} & $100$ & Ring buffer messages per
  axon-role \\
\bottomrule
\end{tabular}
\end{table}

These are operational tunables, not behavioral parameters in the
same sense as the routing parameters. They control timing and
capacity choices that depend on the deployed workload (publish
rate, expected churn, message size); they do not change what the
protocol does, only how often and how much.


% #####################################################################
\part{The Simulator}
% #####################################################################
% The simulator is essential to the design. Two jobs: provide an
% accurate analogy to the real world, and try to break the protocol.


% =====================================================================
\chapter{Philosophy: Accurate Analogy + Adversarial Laboratory}
\label{ch:sim-philosophy}
% =====================================================================

\begin{quote}
  \itshape The simulator is essential to the design.
  Its task is two-fold: provide an accurate analogy to the real
  world, and try to break the protocol.
\end{quote}

\newthought{The simulator} is not a thing we built and put away
once we had measurements. It is part of the design itself --- the
engineering bible's lab notebook --- and it has two jobs, both
ongoing.

Its first job is to provide an \emph{accurate analogy} of the
real-world conditions the protocol must operate under. Real
internet routing, real geographic latency, real WebRTC
connection caps, real churn. Where the analogy is faithful, every
quantitative claim in this volume rests on it; where the analogy
is incomplete, the gap is itself an artifact we describe
(Chapter~\ref{ch:sim-fidelity}). The simulator is honest by
construction --- it does what it does, and we report what it does.

Its second job is to \emph{try to break the protocol}. The
benign tests are not enough. Uniform global lookups against a
uniform global node distribution will succeed at any reasonable
DHT design; success on those tests is necessary but not
sufficient. The simulator's adversarial workload --- Slice World
partitions, sustained churn, Zipf-distributed pub/sub, the
bandwidth-instrumentation reality check --- is what gives the
protocol's claims their force. We pass tests that we
\emph{designed} to make the protocol fail.

This chapter is the philosophy. The next two chapters are the
specifics: what the simulator models with what fidelity
(Chapter~\ref{ch:sim-fidelity}) and what tests we run to break it
(Chapter~\ref{ch:sim-tests}).

% --------------------------------------------------------------------
\section{Why a Simulator and Not a Testbed}
% --------------------------------------------------------------------

\newthought{The two main alternatives} to a simulator for
distributed-systems research are a real testbed (PlanetLab,
EmuLab, or a constellation of cloud-rented VMs) and a closed-form
analytical model. Both have been used productively in this field.
We chose a simulator for a specific reason: it lets us run
experiments at scales where real testbeds become impractical
\emph{and} test conditions where analysis becomes intractable.

\textbf{Scale.} A real testbed at $50{,}000$ nodes is a
non-trivial undertaking: not just the cost of the machines, but
the operational overhead of orchestrating that many participants,
collecting metrics from each, and resetting state between runs.
Most published DHT research has been measured at the
$\sim 100$--$1{,}000$ node scale, with extrapolation to larger
networks. Our simulator runs the full
$50{,}000$-node benchmark suite to convergence in a few hours on
a laptop, with full state observability and exact reproducibility.
That is the scale at which the protocol's claims actually need to
hold.

\textbf{Adversarial conditions.} On a real testbed, you cannot
ask: ``what if half the nodes vanish at once?'' --- not without
disrupting other tenants, or paying for $50{,}000$ machines you
will discard immediately. On a simulator, you can. The Slice
World test (Chapter~\ref{ch:sim-tests}) takes a $25{,}000$-node
network and partitions it into hemispheres connected only by a
single bridge node, then asks the protocol to route between
hemispheres. This kind of test is essentially impossible on a
testbed; on a simulator it is one configuration flag.

\textbf{Reproducibility.} The simulator is deterministic given a
seed. Any measurement in this document can be reproduced exactly
from the open-source code, the parameters listed in the run, and
the seed. Real-testbed measurements have natural variance from
cross-traffic, scheduling jitter, and network conditions that no
researcher can fully control; on a simulator, we report the
mean over hundreds of runs and the variance is small enough to
ignore. \emph{The price is that the simulator only models what we
explicitly tell it to model, and the rest of this Part is honest
about that price.}

% --------------------------------------------------------------------
\section{The Engineering Bible's Lab Notebook}
% --------------------------------------------------------------------

\newthought{Every quantitative claim} in this volume comes from
the simulator. Average lookup latency, hop counts, success rates,
Gini coefficients, bandwidth distribution, churn recovery times
--- all measured by the same code base, against the same
\NDHT implementation, run with the parameters listed in each
chapter.

This puts the simulator at the heart of the document's
credibility. The whitepaper's accuracy is bounded by the
simulator's faithfulness. So Chapter~\ref{ch:sim-fidelity} has
to be brutal about what we \emph{do} model and what we
\emph{don't}: the missing models are part of the document, not a
footnote in the conclusion.

Three categories:

\textbf{Modeled accurately.} Geographic latency from haversine
distance plus per-hop processing cost; XOR routing; bilateral
WebRTC connection caps; topology effects from S2-prefix node
IDs; simulator-time advancement in lookup ticks; LTP / decay /
inertia at full resolution.

\textbf{Modeled approximately.} Churn (instantaneous removal and
re-insertion, no graceful-exit signaling); pub/sub topic
membership (uniformly-sampled subscriber populations, not
Zipf-distributed); message delivery (reliable; no packet loss
beyond explicit dead-peer detection).

\textbf{Not modeled.} WebRTC connection setup time
(ICE/STUN/TURN/DTLS overhead); RPC timeouts and asynchronous
black holes; latency jitter from queueing and bufferbloat;
bandwidth saturation and congestion (until v0.70.04 added
per-node send/receive instrumentation, this was the most
critical omission --- Chapter~\ref{ch:eval-bandwidth}); Byzantine
peer behavior (Sybil attacks, eclipse attacks, prefix forgery).

Each of the ``not modeled'' items is a deferred concern. They
are catalogued in Chapter~\ref{ch:redteam} (the red team) and
have explicit pathways for closure in Chapter~\ref{ch:production}
(the production pathway). The point is that they are
\emph{labeled}: we do not claim measurements that depend on them,
and where claims of behavior under such conditions appear (e.g.,
``the protocol will recover from a 30-second connection-setup
delay''), we mark them as projections, not measurements.

% --------------------------------------------------------------------
\section{Try to Break It as a Design Discipline}
% --------------------------------------------------------------------

\newthought{The second job} of the simulator is the harder one
to do right. A simulator that only runs benign tests will
produce optimistic numbers; a research field full of such
simulators (and we have been in such a field) will collectively
publish protocols that don't survive contact with reality. The
discipline that protects against this is to design tests
specifically intended to make the protocol \emph{fail}.

The test suite (Chapter~\ref{ch:sim-tests}) is organized around
this discipline. For each test we ask: \emph{what hypothesis am
I trying to falsify?}

\begin{itemize}\setlength\itemsep{0.2em}
  \item Slice World tests the hypothesis that the learning
    mechanisms cannot recover from a network partition. (The
    test passes; the hypothesis is falsified at $\sim 94\%$
    cross-partition success.)
  \item Sustained-churn tests the hypothesis that the protocol
    cannot maintain $100\%$ pub/sub delivery under continuous
    node turnover. (The test passes at $5\%$ instantaneous churn;
    fails gracefully past $20\%$ sustained.)
  \item Bandwidth instrumentation tests the hypothesis that
    adaptive routing concentrates load onto a small number of
    nodes (the success-disaster failure mode). (The test
    falsifies it: \NDHT distributes more broadly than Kademlia
    at every tested scale.)
  \item The \texttt{geoBits = 0} ablation tests the hypothesis
    that all of \NDHT's performance comes from the geographic
    prefix and not from the learning. (The test falsifies it:
    with zero geographic information, NX-17 still routes $26\%$
    faster than Kademlia.)
\end{itemize}

Each test was added to the suite after some specific moment in
the design process where we wanted to break the protocol. The
discipline is generative: every time we make a new claim, we
should ask what test would falsify it, and if no such test
exists, we should design one. This is the inverse of
``measurement-driven development''; we call it
\emph{falsification-driven development.}

\textbf{The bandwidth-concentration episode is the canonical
example.} For most of the protocol's development, the simulator
treated all nodes as having infinite bandwidth: every hop cost
$10$\,ms regardless of how many concurrent messages a relay was
processing. This was a known omission, and the red-team analysis
(Chapter~\ref{ch:redteam}) ranked ``bandwidth saturation'' as a
deploy blocker. The hypothesis the red team raised was that
\NDHT's adaptive routing would \emph{amplify} hotspot formation:
LTP would reinforce a fast peer until it saturated; AP scoring
would route around it when it slowed; LTP would re-reinforce it
when it recovered, creating oscillation. We had no way to
test that hypothesis until we instrumented the simulator with
per-node send/receive counters.

The instrumentation took two days to write
(\S\ref{ch:eval-bandwidth} describes it in detail). The first
test it produced \emph{empirically falsified the red-team's
hypothesis}: \NDHT distributes load more broadly than Kademlia at
every tested scale. The simulator's job was to break the
protocol; in this case, it broke the assumption that we needed
new mechanism to address Issue~\#3. The deploy blocker was
reclassified.

This is what we mean when we call the simulator part of the
design. It is not a tool we use after the protocol is finished;
it is the iterative loop in which the protocol is shaped.
Every claim in this volume passed through the loop at least
once; many of the rules in Chapter~\ref{ch:nh1-rules} survived
the consolidation specifically because the simulator showed they
mattered, and others were retired because the simulator showed
they didn't.

% --------------------------------------------------------------------
\section{The Open-Source Simulator}
% --------------------------------------------------------------------

\newthought{The simulator}\sidenote{Source available at the URL
listed in the References. Approximately $25{,}000$ lines of
JavaScript; runs in the browser via Vite-served Web Workers, or
as a Node.js process. Same code, both targets.} is part of the
\NDHT codebase. Anyone can clone the repository, run the test
suite, and reproduce every measurement in this document. The
parameters used for each measurement are recorded in the
benchmark CSV alongside the numbers, so a third party can
verify both the result and the conditions under which it was
obtained.

This commits us to a discipline that is easy to forget about: every
claim in this document is, in principle, falsifiable. A reader
who suspects a measurement is wrong can run the same test on the
same code with the same parameters and either confirm or
contradict the number. We have, on multiple occasions during the
development of this protocol, retracted findings because of
exactly this kind of self-audit. The bandwidth analysis is the
most recent example; Appendix~\ref{app:nx-evolution} catalogues
the full sequence.

The simulator is also a research tool independent of \NDHT. The
code base supports running comparison protocols (Kademlia, \GDHT,
NX-17) against the same workloads with the same instrumentation,
and adding a new protocol is a matter of implementing the
five \texttt{DHT} interface methods (\texttt{addNode},
\texttt{removeNode}, \texttt{buildRoutingTables},
\texttt{lookup}, \texttt{bootstrapJoin}) plus optional pub/sub
hooks. We have used this affordance to run Vivaldi-style
ablations and Coral-style nested-cluster experiments without
disturbing the main protocol.

The next two chapters describe the simulator concretely: what it
models (\S\ref{ch:sim-fidelity}) and what tests it runs
(\S\ref{ch:sim-tests}).


% =====================================================================
\chapter{What the Simulator Models, and What It Does Not}
\label{ch:sim-fidelity}
% =====================================================================

\newthought{This chapter} is the methodological honesty
chapter. Every quantitative claim in this volume rests on the
simulator; the simulator's faithfulness bounds the claims'
credibility. For each modeled phenomenon we describe the model;
for each \emph{not}-modeled phenomenon we describe the omission,
the deferred concern, and the pathway to closing the gap.

The categories we walk through, in order: latency model,
topology model, churn model, message-delivery model, connection
model, bandwidth model, and adversarial-behavior model. Each
gets a section.

% --------------------------------------------------------------------
\section{Latency: Haversine Plus Per-Hop Cost}
% --------------------------------------------------------------------

\textbf{What we model.} Per-hop latency is computed
deterministically from the haversine great-circle distance
between the two nodes' claimed coordinates, using a fiber-optic
propagation constant ($c \approx 200{,}000$ km/s, accounting for
refractive index) plus a fixed per-hop processing
cost.\sidenote{The $10$\,ms-per-hop processing constant is
calibrated from published WebRTC RTT measurements at small
distance --- two nodes in the same data center exhibit roughly
$10$\,ms RTT before great-circle distance becomes the dominant
factor. The constant is conservative; real-world processing
costs vary from $\sim 5$\,ms (server-class) to $\sim 30$\,ms
(mobile under poor radio conditions).}

\[
\mathit{latency}(A, B) \;=\;
2 \cdot \left( \frac{\mathit{haversine}(A, B)}{c} + \delta_{\mathrm{node}} \right)
\]

The factor of $2$ accounts for round trip; $\delta_{\mathrm{node}}$
defaults to $10$\,ms.

\textbf{What we do not model.} \emph{Latency jitter.} Real
internet round-trip times fluctuate $\pm 30\%$ from queueing,
bufferbloat, and asymmetric paths. Our latency is monotone and
clean. The protocol's exponential moving average of recent path
latencies (used as the LTP reinforcement threshold) has an easy
job in the simulator; high-frequency noise could prematurely
decay good synapses or promote lucky-but-unstable ones. The
deferred concern is documented in Chapter~\ref{ch:redteam} as
Tier 1 Issue \#4 and addressed in Chapter~\ref{ch:production} as
the jitter-injection extension to the simulator.

\textbf{What we do not model.} \emph{Asymmetric latency.} A real
internet path from $A$ to $B$ may take a different route than
the path from $B$ to $A$, with different queueing characteristics
on each. Our model is symmetric. The implication is that the
\texttt{forward} and \texttt{reverse} synapse fields collapse to
the same value; no protocol mechanism distinguishes them. Tier 2
Issue \#8 in the red-team analysis.

% --------------------------------------------------------------------
\section{Topology: Uniform Land Distribution}
% --------------------------------------------------------------------

\textbf{What we model.} Nodes are placed at uniformly-distributed
random coordinates over the Earth's land surface (a land-detection
function rejects ocean coordinates). The S2 cell prefix is
computed from the coordinates and stamped into the high bits of
the node ID. The total number of nodes is a configuration
parameter; experiments in this volume use $5{,}000$, $10{,}000$,
$25{,}000$, and $50{,}000$.

\textbf{What we do not model.} \emph{Population-weighted node
distribution.} Real-world peer-to-peer node distribution
correlates with population centers and broadband access; nodes
are not uniformly distributed across land. We expect the
qualitative behavior of the protocol to be insensitive to this
(locality benefits accrue regardless of the population
distribution), but the quantitative gap between regional and
global latency may differ. We have not measured this directly.

\textbf{What we do not model.} \emph{Node-class heterogeneity at
high resolution.} The simulator supports a single
``highway-percent'' parameter to bias a fraction of nodes toward
unrestricted connection caps (server-class deployment); this is
binary (browser-class vs. server-class). Real deployments will
exhibit a continuous spectrum of capacity and reliability tiers.
The current ablation (Chapter~\ref{ch:eval-perf}) is sufficient
for first-order conclusions but does not capture the full
heterogeneity.

% --------------------------------------------------------------------
\section{Churn: Instantaneous Removal and Re-Insertion}
% --------------------------------------------------------------------

\textbf{What we model.} A churn event removes a percentage of
live nodes and replaces them with fresh nodes that bootstrap
through random sponsors. ``Removed'' means: \texttt{alive} is
set to false, all incoming references are observable to the
remaining network as dead synapses on first attempt. ``Inserted''
means: a new node is created with fresh ID, the bootstrap
operation runs as in \S\ref{sec:nh1-bootstrap-tick}.

The simulator supports two churn modes:

\begin{itemize}\setlength\itemsep{0.15em}
  \item \emph{Instantaneous churn} --- a single event at a
    specific simulation time. ``$5\%$ instantaneous churn'' means
    one event removes and replaces $5\%$ of the live nodes, then
    the test continues against the new population.
  \item \emph{Continuous churn} --- a churn event fires every $5$
    simulation ticks at a configurable rate. ``$1\%$ continuous
    churn'' means each event replaces $1\%$ of the live nodes,
    so the cumulative turnover grows linearly with time.
\end{itemize}

\textbf{What we do not model.} \emph{Graceful exit signaling.} Real
clients leaving a network might announce their intent (``I am
disconnecting''); our churn is silent removal. The protocol's
recovery mechanism (dead-peer detection) covers both cases ---
a graceful exit and an abrupt disappearance look the same from
the routing layer's perspective --- but a real protocol can
exploit the graceful case for proactive handoff. We do not.

\textbf{What we do not model.} \emph{Clustered churn.} A real
network outage typically affects a geographic region (a power
failure, a transit cable cut, a regional ISP problem) rather
than a uniformly-sampled subset of nodes. Our churn is uniformly
sampled. We have run informal experiments with regional churn
(removing all nodes in a specific S2 cell) and the protocol
recovers --- the partition-recovery mechanism in Slice World
generalizes --- but the formal benchmark uses the uniform-sample
version.

% --------------------------------------------------------------------
\section{Message Delivery: Reliable, with Explicit Dead-Peer Detection}
% --------------------------------------------------------------------

\textbf{What we model.} A message sent from $A$ to $B$ is
delivered if $B$ is alive at the moment of delivery. The
\texttt{SimulatedNetwork.send} primitive is synchronous and
reliable: it returns immediately with the response, accounting
for the latency from the model in \S5.1.

\textbf{What we do not model.} \emph{Packet loss.} Real internet
delivery has a small but nonzero loss rate, typically
$0.1$--$1\%$. We do not lose packets except in the explicit
``the peer is dead'' case. The protocol's response to single-message
loss is the same as its response to a dead peer (timeout +
iterative fallback); we do not test the regime where the loss
rate is high but bounded.

\textbf{What we do not model.} \emph{Asynchronous black holes ---
RPC timeouts.} A real RPC to a peer that has gone silent can
stall for seconds before the caller knows the peer is dead. Our
simulator detects death instantaneously: \texttt{network.send}
to a dead peer returns immediately with a ``dead'' indication.
This is a fundamental simplification. The protocol's churn-recovery
measurements in this document assume instant detection; a real
deployment will see multi-second timeout windows during which
messages are lost. The deferred concern is Tier 1 Issue \#2.

\textbf{What we do not model.} \emph{Message corruption.} We
assume bytes delivered intact. Real-world deployments must use
authenticated and integrity-protected message formats; the
simulator does not exercise this layer.

% --------------------------------------------------------------------
\section{Connection Model: Bilateral Cap, Instant Setup}
% --------------------------------------------------------------------

\textbf{What we model.} Each node has a configurable
\texttt{maxConnections} cap (default $100$, with the synaptome
itself capped at $50$). When node $A$ wants to add $B$ as a
synapse, the simulator's \texttt{tryConnect}$(A, B)$ checks
whether \emph{both} sides have an open connection slot. If
either side is full, the connection is silently refused (the
caller's logic falls through to the next candidate). This is the
simulator's first-order analogue of WebRTC's per-browser
data-channel limit.

\textbf{What we do not model.} \emph{Connection setup time.} Real
WebRTC connections require ICE (negotiation), STUN (NAT
discovery), TURN (relay fallback), and DTLS (handshake). The
total setup time is typically $1.5$--$3$\,seconds in the wild,
and successful setup is not guaranteed --- ICE failure rates of
$10$--$15\%$ are common in the WebRTC literature.\sidenote{See
the WebRTC measurement studies, e.g.\
\href{https://webrtchacks.com}{webrtchacks.com}, for empirical
distributions of setup time and failure rate. The numbers vary
strongly by NAT type and by the availability of TURN relays.}

The implication: every architectural mechanism in \NDHT that
\emph{depends} on rapid synapse establishment --- Slice World's
re-stitching after partition; annealing's
$2$-hop-neighborhood replacement; hop caching's deposit-and-use;
the entire pub/sub re-subscribe self-heal --- looks faster in
the simulator than in production. The simulator says
``$94\%$ Slice World recovery in $500$ lookups''; a production
system facing $1$--$3$\,second connection setup per new
synapse would take much longer. Tier 1 Issue \#1 of the red team.

\textbf{What we do not model.} \emph{Connection failure.} Once
\texttt{tryConnect} succeeds, the connection is permanent
(until one side dies). Real connections drop intermittently and
need re-establishment. The protocol's mechanisms for handling
this are equivalent to dead-peer detection, but the rate is much
higher than churn alone would predict.

% --------------------------------------------------------------------
\section{Bandwidth Model: Per-Node Counters, No Saturation}
% --------------------------------------------------------------------

\textbf{What we model.} As of v0.70.04, every routing chokepoint
in the simulator is instrumented with per-node send/receive
counters. The instrumentation captures: every
\texttt{SimulatedNetwork.send} call (Kademlia, \GDHT); every
direct nodeMap-mediated message in NX-6, NX-17, and NH-1
(routing hops, two-hop probes, LTP feedback walks, hop caching,
lateral spread, triadic introductions, local-candidate scans for
annealing and dead-peer replacement, axonal tree
\texttt{routeMessage} forwarding, \texttt{sendDirect} entry).

The counters are per-node, per-message-type. Each lookup tick
resets them; a snapshot is taken at the end of every benchmark
or training cycle. Results: total messages, mean / p50 / p99 /
max per-node load, Gini coefficient over the per-node
distribution, count of nodes processing more than $10\times$ or
$100\times$ the cycle mean. Chapter~\ref{ch:eval-bandwidth}
analyzes the resulting distributions.

\textbf{What we do not model.} \emph{Bandwidth saturation
proper.} The counters tell us \emph{how many} messages each node
is processing; they do not feed back into the protocol's
behavior. A node processing $1000\times$ the network mean
incurs no additional delay, no drop, no backpressure. The
protocol's AP scoring uses static latency, not load-adjusted
delay. The implication is that we can characterize the
\emph{distribution} of load (and we have, decisively, in
Chapter~\ref{ch:eval-bandwidth}) but we cannot test whether the
protocol oscillates against actual saturation. The deferred
concern is Tier 1 Issue \#3 in the red team.

The reclassification of Issue \#3 in Chapter~\ref{ch:eval-bandwidth}
is precisely about the \emph{distribution} question, not the
\emph{saturation} question. The distribution measurement
falsifies the hypothesis that \NDHT's adaptive routing
concentrates load (it doesn't; the load is broadly distributed
at every tested scale, with fewer hot nodes than Kademlia). The
saturation question --- ``does the protocol oscillate when an
actually-saturated peer slows down?'' --- requires a load-aware
AP scoring extension and a saturation simulation, both deferred
to the production pathway (Chapter~\ref{ch:production}).

% --------------------------------------------------------------------
\section{Adversarial Behavior: Cooperative Trust Throughout}
% --------------------------------------------------------------------

\textbf{What we model.} Nothing adversarial. Every node in the
simulator obeys the protocol exactly: correct ID assignment,
correct routing-table maintenance, correct pub/sub handler
behavior, correct dead-peer signaling.

\textbf{What we do not model.} \emph{Sybil attacks.} An attacker
generates many forged identities all bidding for routing-table
positions in a target region. The simulator does not produce or
defend against Sybils.

\textbf{What we do not model.} \emph{Eclipse attacks.} An
attacker controls all peers in a target node's synaptome,
isolating the target from the rest of the network. The simulator
does not produce or defend against eclipses.

\textbf{What we do not model.} \emph{Prefix forgery.} A node
falsifies its S2 prefix to land in a region where it does not
physically reside. The simulator's nodes have correct prefixes
by construction. The locality benefits we report assume
cooperative trust; an adversarial deployment with a fraction of
forged-prefix nodes would experience proportionally degraded
locality (not catastrophically degraded routing, but the
qualitative claims about regional latency would weaken).

\textbf{What we do not model.} \emph{Routing-deviation attacks.}
A peer that selectively drops, corrupts, or misroutes messages.
The simulator's nodes route honestly.

The full catalog of adversarial threats and the mitigations
identified for each is in Chapter~\ref{ch:redteam}, Tier 2 (Issues
5--8: Sybil, churn-class heterogeneity, Byzantine pub/sub,
asymmetric reachability) and Tier 3 (Issues 9--13: hardening
items). The protocol's correctness does not depend on any of
the threat models being defended against in the simulator;
performance claims in the presence of these threats require
the corresponding mitigations in place.

% --------------------------------------------------------------------
\section{The Honest Closure}
% --------------------------------------------------------------------

\newthought{The protocol's results} measured in this volume are
\emph{accurate for the conditions modeled}. The conditions
modeled are not the full conditions a real deployment will
face. The gap between the two is documented above and ranked in
the red team (Chapter~\ref{ch:redteam}); the pathway to closing
it is in Chapter~\ref{ch:production}.

The honest framing: the brain is production-ready; the body is
not. The protocol's algorithmic core --- the substrate, the
five operations, the consolidation around vitality, the
axonal-tree pub/sub primitive --- is well-characterized and
well-measured at scales up to $50{,}000$ nodes. The transport
layer --- connection-setup time, RPC timeouts, jitter,
bandwidth saturation, adversarial peer behavior --- is the next,
\emph{measurable, scopeable} problem. The simulator extensions
described in Chapter~\ref{ch:production} are how that work
gets done.

What this means for the document's claims: every quantitative
claim is conditional on the simulator's modeling choices. A
careful reader can take any number in this document and ask
``does this depend on something the simulator does not model?''
Where the answer is yes (e.g., ``does the $94\%$ Slice World
recovery rate hold under realistic connection-setup time?''), the
honest answer is ``no, it would be lower.'' We do not hide that.
The simulator's value is that it lets us decompose the protocol's
behavior into separable components, each measurable on its own,
each addressable in turn.


% =====================================================================
\chapter{The Adversarial Test Suite}
\label{ch:sim-tests}
% =====================================================================

\newthought{Each test} in the simulator's suite was added with a
specific hypothesis to falsify. This chapter walks the suite,
giving for each test: \emph{what it does}, \emph{what hypothesis
it is designed to break}, \emph{what the result actually is}, and
\emph{what counts as a failure} (so a future test run that
\emph{does} fail will be unambiguous).

Tests are organized by what they stress: routing correctness
under benign conditions; routing correctness under partition;
robustness under churn; bandwidth distribution under load;
ablation tests that isolate specific mechanisms.

% --------------------------------------------------------------------
\section{Global and Regional Routing}
% --------------------------------------------------------------------

\textbf{What it does.} The benign baseline. Uniformly-random
source and destination; measure mean latency, mean hop count,
success rate over $500$ lookups per cell. Variants restrict the
destination to a regional disk: $500$\,km, $1{,}000$\,km,
$2{,}000$\,km, $5{,}000$\,km radius around the source.

\textbf{Hypothesis under test.} \emph{The protocol does not
reach within a small constant of the analytical $3\delta$
floor on global lookups, and it does not approach single-hop
performance on regional lookups.} If either claim fails, the
protocol's value proposition versus Kademlia is gone.

\textbf{What counts as failure.} Global mean latency more than
$2\times$ the $3\delta$ floor. Regional mean hop count
indistinguishable from global. Success rate below $99\%$ on any
benign workload.

\textbf{Result.} The hypothesis is falsified --- by the
\emph{floor-hugging} property of the actual measurements
(Chapter~\ref{ch:eval-perf}). NX-17 hits $1.16\times$ the floor on
global lookups at $25{,}000$ nodes; NH-1 hits $1.27\times$.
Regional routing converges to single-digit-hop performance, and
the gap between regional and global widens as $N$ grows. Success
rate is $100\%$ across all benign workloads at all tested scales.

% --------------------------------------------------------------------
\section{Concentrated Workloads (Hot Source / Hot Dest / Hot Pair)}
% --------------------------------------------------------------------

\textbf{What it does.} The realistic-traffic test. ``Hot
destination'' restricts the destination pool to $10\%$ of nodes
(the ``$10\%$ dest'' workload); the source is the remaining
$90\%$. ``Hot source'' is the converse; $10\%$ of nodes generate
all traffic. ``Hot pair'' restricts both: the same $10\%$ source
pool routes to a different (or same) $10\%$ destination pool.

The pattern models real applications: a CDN serves content from
a small set of replicas; a chat group routes among a small set
of frequently-conversing peers; a measurement system aggregates
into a small set of collectors.

\textbf{Hypothesis under test.} \emph{Adaptive routing does not
provide measurably more benefit on concentrated workloads than
on uniform workloads.} If the hypothesis holds, the
``learning'' part of \NDHT does not justify its complexity; you
might as well use \GDHT.

\textbf{What counts as failure.} The latency advantage of \NDHT
over \GDHT shrinks toward unity as the workload becomes more
concentrated. The hypothesis is intentionally counter-intuitive
to provide a sharp falsification target.

\textbf{Result.} The hypothesis is falsified --- in the opposite
direction. \NDHT's latency advantage \emph{grows} with
concentration. At $25{,}000$ nodes:

\begin{itemize}\setlength\itemsep{0.1em}
  \item Global random: NH-1 at $260$\,ms, \GDHT at $282$\,ms (8\% improvement)
  \item $10\%$ dest: NH-1 at $43$\,ms, \GDHT at $108$\,ms (60\% improvement)
  \item $10\% \to 10\%$: NH-1 at $33$\,ms, \GDHT at $109$\,ms (70\% improvement)
\end{itemize}

The mechanism producing the gap is hop caching: under
concentrated workloads, the same destinations are reached from
many sources; intermediate nodes accumulate shortcuts to those
destinations; subsequent lookups bypass entire prefix paths. \GDHT
has no such mechanism.

% --------------------------------------------------------------------
\section{Slice World: Single-Bridge Partition Recovery}
% --------------------------------------------------------------------

\textbf{What it does.} The sharpest stress test in the
suite.\sidenote{Named after the Tobias Werner short story
``Sliceworld,'' in which Earth is bisected. The reference is
loose; the test is sharper.} The
network is partitioned almost in half: Eastern hemisphere on one
side, Western hemisphere on the other, connected only through a
\emph{single bridge node} placed near Hawaii. Every other
cross-hemisphere edge is removed before the test begins.

The protocol is then asked to route between hemispheres: $500$
lookups whose source is in one hemisphere and destination is in
the other.

\textbf{Hypothesis under test.} \emph{The learning mechanisms
cannot recover from a network partition with only a single
bridge.} The hypothesis is plausible because the bridge is a
geometric bottleneck --- if every cross-hemisphere lookup must
route through the bridge, the bridge would become saturated,
the partition would not be recovered in any meaningful sense,
and the success rate would be far below $100\%$.

\textbf{What counts as failure.} Cross-partition success rate
below $80\%$. The threshold is generous; even a partial recovery
(say, $50\%$) would be a substantial improvement over Kademlia's
$0\%$, but a protocol whose learning genuinely worked should hit
much higher.

\textbf{Result.} The hypothesis is falsified at $\sim 94\%$
cross-partition success. The mechanism is the bridge becoming
a \emph{seed crystal}: hop caching deposits cross-hemisphere
synapses in every intermediate node along each successful
bridge crossing; triadic closure introduces direct edges
between observed transit pairs; lateral spread propagates the
new edges to geographic neighbors. After the first hundred
lookups, hundreds of cross-hemisphere synapses exist that
bypass the bridge entirely. The bridge is no longer special.

For comparison: Kademlia achieves $0\%$ (no learning, the
partition is permanent); \GDHT achieves $4.6\%$ (geographic
prefix points a few peers at the bridge but no mechanism builds
on the discovery). NX-17 and NH-1 both achieve $\sim 94\%$.

% --------------------------------------------------------------------
\section{Pub/Sub Delivery Under Churn}
% --------------------------------------------------------------------

\textbf{What it does.} A topic with a configurable number of
subscribers (default $32$) is created. The publisher sends a
test message; the simulator counts how many subscribers
received it. The test is then repeated under increasing churn
rates: $5\%$ instantaneous, $10\%$, $20\%$, up to $30\%$.

A separate variant runs continuous churn: $1\%$ of live nodes
per $5$ ticks, with three refresh intervals between churn
events to allow the tree to heal. This produces a
sustained-load curve.

\textbf{Hypothesis under test.} \emph{Axonal-tree self-healing
cannot maintain $100\%$ pub/sub delivery under realistic
churn.} The hypothesis is plausible because the tree is built
by routed subscribe; under churn, parts of the routing
topology change between when the tree is built and when a
publish flows; the message could miss subscribers attached to
branches that were severed.

\textbf{What counts as failure.} Steady-state baseline delivery
below $99\%$. Recovery from $5\%$ instantaneous churn taking
more than three refresh intervals to return to $100\%$.
Continuous-churn delivery dropping below $90\%$ at $5\%$
sustained churn rate.

\textbf{Result.} Steady-state baseline: $100\%$ delivery at
$25{,}000$ nodes with $79$ topic groups of $32$ subscribers
each. After $5\%$ instantaneous churn, immediate delivery
falls to $99.9\%$ and recovers to $100\%$ after three refresh
intervals.

The continuous-churn variant produces a smooth degradation
curve rather than a cliff: $98.7\%$ delivered at $5\%$
cumulative churn, $91.2\%$ at $10\%$, $86.5\%$ at $20\%$,
$70.0\%$ at $25\%$, $52.4\%$ at $30\%$. Crucially, no cliff:
delivery degrades \emph{linearly} with sustained churn,
indicating the protocol finds a steady state where new
re-subscribes match losses. The dominant residual failure mode
at high churn is not broken routing but \emph{subscribers
temporarily captured at relays that have lost their delivery
path to the root} --- the failure mode the temporal replay
cache (Chapter~\ref{ch:pubsub}) was designed to close.

% --------------------------------------------------------------------
\section{Bandwidth Distribution: The Concentration Audit}
% --------------------------------------------------------------------

\textbf{What it does.} Every per-node send and receive event is
counted, by message type. Per cycle (a unit of $500$ lookups),
the counters are aggregated into a distribution: total messages,
mean / p50 / p99 / max per-node load, Gini coefficient over the
per-node totals, count of nodes whose per-cycle traffic exceeds
$10\times$ or $100\times$ the network mean.

The instrumentation is added at every routing chokepoint,
including the direct-nodeMap-mediated paths in NX-6, NX-17,
and NH-1 that bypass \texttt{SimulatedNetwork.send}. Without
explicit instrumentation, the simulator cannot see this traffic
and the protocol's true bandwidth profile is invisible.

\textbf{Hypothesis under test.} \emph{Adaptive routing
concentrates load on a small number of nodes more aggressively
than non-adaptive routing.} The hypothesis was raised in the
red-team analysis (Chapter~\ref{ch:redteam}) as the
\emph{success-disaster} failure mode: LTP reinforces a fast peer
until it saturates, AP scoring routes around it when it slows,
LTP re-reinforces it when it recovers, oscillation ensues.

\textbf{What counts as failure.} \NDHT producing more
$>100\times$-mean nodes than Kademlia at any tested scale.
\NDHT's Gini coefficient growing faster with $N$ than
Kademlia's. The max/mean ratio worsening monotonically with $N$.
Any of these would suggest the adaptive mechanism is amplifying
concentration.

\textbf{Result.} The hypothesis is falsified, decisively. At
$50{,}000$ nodes:

\begin{itemize}\setlength\itemsep{0.1em}
  \item Kademlia: $56$ nodes processing $>100\times$ the mean,
    Gini $0.940$, max/mean $208\times$.
  \item \GDHT: $62$ nodes processing $>100\times$ the mean,
    Gini $0.932$, max/mean $213\times$.
  \item NX-17 default: $\mathbf{0}$ nodes $>100\times$, Gini $0.694$,
    max/mean $61\times$.
  \item NH-1 default: $\mathbf{0}$ nodes $>100\times$, Gini $0.661$,
    max/mean $47\times$.
\end{itemize}

The classical DHTs produce dozens of catastrophic-bandwidth
nodes at $50{,}000$ scale; the adaptive DHTs produce zero.
Chapter~\ref{ch:eval-bandwidth} reports this in full detail
with the $4 \times 4$ sweep and the rate-knee curve.

The episode is the canonical example of falsification-driven
development. We added the bandwidth instrumentation
specifically to test the red team's hypothesis. The test broke
the hypothesis, not the protocol.

% --------------------------------------------------------------------
\section{Ablation: The geoBits = 0 Test}
% --------------------------------------------------------------------

\textbf{What it does.} The ablation removes the geographic
prefix entirely. Node IDs are pure $H_{64}(\mathit{publicKey})$
--- random. Otherwise the protocol runs unchanged: same
synaptome capacity, same vitality, same five operations.

\textbf{Hypothesis under test.} \emph{All of \NDHT's
performance benefit comes from the geographic prefix; the
neuromorphic learning is gravy.} A skeptic's hypothesis,
because if it were true the protocol's contribution beyond
\GDHT would be marginal.

\textbf{What counts as failure.} \NDHT's latency at
\textsc{geoBits} $= 0$ being indistinguishable from Kademlia's.
The performance attribution would belong entirely to the
prefix.

\textbf{Result.} The hypothesis is falsified. With \emph{zero}
geographic information, NX-17 still routes $26\%$ faster than
Kademlia at $25{,}000$ nodes, and NH-1 routes $8\%$ faster. The
learning mechanisms do real work, attributable cleanly to the
synaptome shaping itself by traffic.

The smaller advantage for NH-1 vs NX-17 in this ablation
($8\%$ vs $26\%$) is itself informative: NX-17's two-tier
machinery and Markov pre-learning give it more starting
structure to exploit when the geographic prefix is removed.
NH-1, having consolidated those into the unified vitality
model, has less independent structural redundancy. The
ablation shows the cost of the consolidation: a small
performance drop in the no-locality regime, consistent with
the broader claim that NH-1 trades $\sim 5$--$10\%$ peak
performance for $\sim 30\%$ parameter-count reduction.

% --------------------------------------------------------------------
\section{Ablation: Highway Percentage}
% --------------------------------------------------------------------

\textbf{What it does.} A fraction of nodes (the
\texttt{highwayPct} parameter) are designated as
``highway'' --- they have unrestricted connection caps,
modeling server-class deployments. The remaining nodes are
browser-class with $\texttt{maxConnections} = 100$.

The sweep runs at $0\%$ (uniform browser-class), $5\%$, $15\%$,
$30\%$, and $100\%$ (uniform server-class) highway.

\textbf{Hypothesis under test.} \emph{Real-world deployment will
require server-class infrastructure for acceptable performance,
making the protocol effectively centralizing.} If the hypothesis
holds, the protocol fails its mission.

\textbf{What counts as failure.} Latency at $0\%$ highway being
much worse than at $100\%$. A linear performance scaling with
highway percentage --- showing that infrastructure is doing the
work, not the protocol.

\textbf{Result.} The hypothesis is falsified.
$15\%$ server-class nodes captures $70\%$ of the latency
improvement available from going to $100\%$ server-class.
Real-world deployments don't need a uniform server fleet; a
modest fraction of capable nodes (residential desktops with no
mobile/NAT constraints, plus the operator's own infrastructure)
is sufficient.

The shape of the curve is concave: most of the benefit comes
from getting the first $5$--$15\%$ above zero. Pure
browser-class deployment ($0\%$ highway) is workable, just
slower. Pure server-class deployment is not necessary.

% --------------------------------------------------------------------
\section{The Production Test: 50,000 Nodes End-to-End}
% --------------------------------------------------------------------

\textbf{What it does.} The full benchmark suite at $50{,}000$
nodes: global, regional ($500$\,km, $1{,}000$\,km, $2{,}000$\,km,
$5{,}000$\,km), concentrated workloads (source/dest/srcdest at
$10\%$), pub/sub baseline and pub/sub under $5\%$ churn. Each
test cell runs $500$ lookups; the full suite takes a few hours
on a laptop.

\textbf{Hypothesis under test.} \emph{The protocol's results at
$25{,}000$ nodes do not extend to $50{,}000$.} Some
quadratic-cost mechanism could be hiding in the smaller-scale
measurements; some concentration effect could appear at scale
that does not appear at half-scale.

\textbf{What counts as failure.} Global latency growing faster
than $\log N$. Regional latency \emph{growing} with $N$. Success
rate dropping. Pub/sub baseline below $99\%$. Bandwidth max/mean
ratio worsening.

\textbf{Result.} The hypothesis is mostly falsified. Global
latency tracks the $\log N$ curve cleanly; NH-1 holds its
$1.27\times$ floor multiplier. Regional latency continues to
\emph{decrease} with $N$ on the $500$\,km radius (more nodes per
region, more shortcut opportunities). Success rate stays at
$100\%$ on benign workloads. Pub/sub baseline stays at $100\%$.
Bandwidth max/mean ratio actually \emph{decreases} with $N$ for
all four protocols (because the per-cycle work is spread over
more nodes), and \NDHT continues to produce zero
$>100\times$-mean hot nodes while Kademlia produces $56$.

The qualifier ``mostly'' is for one observation: the
\emph{absolute} count of $>10\times$-mean hot nodes in \NDHT
\emph{grows} with $N$ (1{,}419 at $50{,}000$ vs $204$ at
$25{,}000$). The increase is sub-linear in $N$ ($7\times$
growth for $2\times$ scale) and not catastrophic, but it is the
one curve we have not flattened. The mitigation, if needed, is
the \texttt{annealRateScale} knob discussed in
Chapter~\ref{ch:eval-bandwidth}.

% --------------------------------------------------------------------
\section{What the Suite Does Not Yet Cover}
% --------------------------------------------------------------------

\newthought{Honest closure.} The test suite catches a wide range
of failure modes but is not exhaustive. Tests we have not yet
implemented:

\begin{itemize}\setlength\itemsep{0.15em}
  \item \emph{Adversarial nodes.} Sybils, eclipses, prefix
    forgery, routing-deviation attacks. Discussed in
    Chapter~\ref{ch:redteam} (Tier 2).
  \item \emph{Friction-realistic latency.} Connection setup,
    RPC timeouts, jitter. Tier 1 of the red team; closure path
    in Chapter~\ref{ch:production}.
  \item \emph{Zipf-distributed pub/sub.} A small number of
    very popular topics with high publish rate; tests the
    hot-axon-root saturation question. Currently not in the
    suite; addresses one residual of the bandwidth question
    not closed by Chapter~\ref{ch:eval-bandwidth}.
  \item \emph{Variable-rate churn.} Continuous churn at a
    rate that varies over time (peak-hour pattern). Inspired
    by Ghinita-Teo's adaptive-stabilization work
    (Chapter~\ref{ch:lineage}).
  \item \emph{Heterogeneous-capacity heterogeneity.} The
    current highway sweep is binary. A continuous distribution
    of node capabilities is more realistic.
\end{itemize}

Each gap is named, located in the larger road map (red team or
production pathway), and assigned to a future iteration. The
fact that the gaps are named is the point: a test we
\emph{know} we need to add is a deferred concern, not an
unknown unknown. The simulator-as-laboratory commitment makes
the gaps part of the engineering plan, not a footnote.


% #####################################################################
\part{Evaluation}
% #####################################################################
% All benchmarks re-run under the NH-1 generation. NX-10 era numbers are
% retired.


% =====================================================================
\chapter{Methodology}
\label{ch:eval-method}
% =====================================================================

\newthought{Every measurement} reported in this Part was produced
by the same simulator, run with the same default parameters,
unless explicitly noted. This chapter records those parameters and
the procedures we follow so that any number can be reproduced from
the open-source code base by a third party.

% --------------------------------------------------------------------
\section{Default Run Configuration}
% --------------------------------------------------------------------

Unless a chapter says otherwise, every benchmark in
Part~V uses these settings:

\begin{itemize}\setlength\itemsep{0.15em}
  \item \textbf{Node count.} Reported per measurement; typically
    $25{,}000$ or $50{,}000$. The full benchmark suite has been
    run at $5{,}000$, $10{,}000$, $25{,}000$, and $50{,}000$.
  \item \textbf{Bilateral cap (max connections).} $100$. This is
    the WebRTC-realistic deployment target: a typical desktop
    browser sustains $100$--$200$ data channels comfortably; the
    cap at $100$ leaves headroom for application-level WebRTC
    traffic outside the DHT.
  \item \textbf{Synaptome capacity.} $50$. Half the connection
    cap; the routing table cannot use all available connection
    slots, leaving room for bootstrapping new joiners.
  \item \textbf{Highway percentage.} $0\%$. Uniform browser-class
    deployment unless we are running the highway-percentage
    sweep (\S\ref{ch:eval-perf}).
  \item \textbf{geoBits.} $8$. The standard S2 cell prefix
    width. Ablations at $0$ and $16$ are reported separately.
  \item \textbf{Lookups per cell.} $500$. Each measurement is the
    mean over $500$ lookups within a single workload pattern.
    Standard deviation is reported when relevant; on stable
    workloads it is consistently small ($\pm 5\%$ of mean).
  \item \textbf{Bidirectional bootstrap.} Yes. Both sides of
    every bootstrap edge are recorded.
  \item \textbf{Web limit.} Yes. The bilateral cap is enforced.
  \item \textbf{Random seed.} The benchmarks use a fixed seed
    per run so results are deterministic. Different seeds produce
    different absolute numbers but the same qualitative
    conclusions; we have not seen a measurement where the
    qualitative ordering of protocols changes with seed.
\end{itemize}

The full parameter list, including every implementation constant
not exposed for routine tuning, is in
Appendix~\ref{app:params}.

% --------------------------------------------------------------------
\section{The Test Suite Cells}
% --------------------------------------------------------------------

A \emph{cell} is one (workload, protocol, parameter) tuple. The
benchmark sweep runs a list of cells per protocol and aggregates
into a CSV per benchmark. Workloads:

\begin{itemize}\setlength\itemsep{0.1em}
  \item \texttt{global} --- uniformly-random source and
    destination
  \item \texttt{r500}, \texttt{r1000}, \texttt{r2000},
    \texttt{r5000} --- regional lookups within radius (km) of source
  \item \texttt{src} --- $10\%$ source pool, uniform destinations
  \item \texttt{dest} --- $10\%$ destination pool (XOR-nearest
    selection), uniform sources
  \item \texttt{srcdest} --- both pools restricted to $10\%$,
    non-overlapping
  \item \texttt{pubsub} --- pub/sub baseline delivery; topic
    coverage and group size from the run config
  \item \texttt{churn} --- pub/sub delivery under $5\%$
    instantaneous churn after warm-up
\end{itemize}

For neuromorphic protocols, each test runs a warm-up phase first
(typically $5{,}000$ lookups for the destination pattern) so the
synaptome is shaped for the test workload before measurement
begins. \KDHT and \GDHT have no warm-up; their routing tables are
static after bootstrap.

% --------------------------------------------------------------------
\section{The 3\boldmath$\delta$ Floor}
\label{sec:eval-floor}
% --------------------------------------------------------------------

\textbf{What it is.} Dabek et al.\sidenote{Dabek, Li, Sit,
Robertson, Kaashoek, Morris. ``Designing a DHT for low latency
and high throughput.'' NSDI 2004.} proved a lower bound on lookup
latency in any recursive DHT: the median end-to-end lookup time
is at least $3\delta$, where $\delta$ is the median one-way RTT
between random pairs of nodes on the underlying network.

The proof is geometric. Each hop in any DHT must, on average,
cover half the remaining distance. The hop latencies sum to
$\delta + \delta/2 + \delta/4 + \ldots = 2\delta$. One additional
hop is needed for delivery, contributing $\delta$. Total:
$3\delta$.

\textbf{Where we are.} The simulator's $\delta$ at $25{,}000$
uniformly-distributed nodes on land is $68$\,ms (median),
$70$\,ms (mean), $123$\,ms (p95). The floor is therefore
$204$\,ms (median basis) or $209$\,ms (mean basis). Every
``how close to the floor are we'' analysis in this Part is
relative to that number.

The floor is a lower bound, not a target. A protocol that beats
the floor at small $N$ is not contradicting the theorem: the
theorem describes the asymptotic geometric average. A protocol
that hugs the floor at large $N$ is in some sense
\emph{unimprovable} on this metric without changing the underlying
network topology.

% --------------------------------------------------------------------
\section{Reporting Conventions}
% --------------------------------------------------------------------

Three conventions worth being explicit about.

\textbf{Latency means RTT.} All reported latencies are round-trip
times in milliseconds. A ``$240$\,ms global lookup'' means the
total wall-clock time from when the source initiates the lookup
to when it receives confirmation; both directions of every hop
are accounted for.

\textbf{Hop counts include the destination hop.} A lookup that
reaches the target in two intermediate hops is a $3$-hop lookup:
source $\to$ peer $\to$ peer $\to$ target. This matches the
$3\delta$-floor convention (the floor's third $\delta$ is the
delivery hop).

\textbf{Success rate is per-lookup, not per-cell.} A success
rate of $99.9\%$ on $500$ lookups means $499$ succeeded and one
failed; ``$99.5\%$'' means at most three failed; etc. We report
$100\%$ when no failures were observed, which is the case on all
tested benign workloads at all tested scales.


% =====================================================================
\chapter{Performance Results}
\label{ch:eval-perf}
% =====================================================================

\newthought{This chapter} reports the headline numbers for the
four-way protocol comparison at $25{,}000$ nodes. We compare
\KDHT (Kademlia, the historical baseline), \GDHT (Kademlia plus
geographic prefix, our prior work), NX-17 (the prior-generation
state-of-the-art neuromorphic protocol), and NH-1 (the
consolidated current generation). All numbers are from a
single benchmark sweep run on simulator v0.70.05 with the
standard methodology of Chapter~\ref{ch:eval-method}.

% --------------------------------------------------------------------
\section{The Headline: Four Protocols, Ten Workloads, 25,000 Nodes}
% --------------------------------------------------------------------

Table~\ref{tab:perf-headline} is the headline result. Each cell
is the mean latency in milliseconds over $500$ lookups; the
3$\delta$ floor is $204$\,ms (median basis), $209$\,ms (mean basis).

\begin{table*}[h]
\centering
\small
\caption{Mean latency (ms) at $25{,}000$ nodes, $500$ lookups per
cell. The 3$\delta$ floor at this scale is $204$\,ms (median) /
$209$\,ms (mean). \textbf{Bold} marks the winning protocol per
workload.}
\label{tab:perf-headline}
\begin{tabular}{@{}lrrrr@{}}
\toprule
\textbf{Workload} & \textbf{\KDHT} & \textbf{\GDHT} &
\textbf{NX-17} & \textbf{NH-1} \\
\midrule
global random           & $511.6$ & $282.7$ & $\mathbf{237.3}$ & $260.8$ \\
$500$\,km regional      & $499.1$ & $153.4$ & $\mathbf{80.0}$  & $93.9$ \\
$1{,}000$\,km regional  & $505.1$ & $162.0$ & $\mathbf{91.6}$  & $107.8$ \\
$2{,}000$\,km regional  & $514.6$ & $178.5$ & $\mathbf{103.2}$ & $126.0$ \\
$5{,}000$\,km regional  & $496.9$ & $205.4$ & $\mathbf{147.1}$ & $166.6$ \\
$10\%$ source pool      & $513.0$ & $292.2$ & $\mathbf{243.9}$ & $254.3$ \\
$10\%$ dest pool        & $242.7$ & $108.0$ & $\mathbf{40.6}$  & $43.2$ \\
$10\% \to 10\%$ pair    & $244.3$ & $109.1$ & $\mathbf{31.9}$  & $33.3$ \\
$5\%$ churn (post-warm) & $477.6$ & $279.2$ & $\mathbf{239.6}$ & $279.5$ \\
\bottomrule
\end{tabular}
\end{table*}

The overall picture: NX-17 wins every cell, NH-1 trails NX-17 by
$5$--$15\%$ across the board, and both adaptive protocols
dominate \GDHT, which dominates \KDHT.

The $5\%$-churn row is the resilience test. \KDHT and \GDHT
performance under churn is essentially indistinguishable from
their global baselines, because they don't actually adapt to
churn beyond replacing dead bucket entries when next they're
queried. NX-17 holds at $240$\,ms; NH-1 holds at $280$\,ms ---
NH-1 takes a slightly larger relative hit because its
consolidated mechanisms for repair are slightly less aggressive
than NX-17's two-tier highway maintenance, and the ablation
cost is most visible under churn.

% --------------------------------------------------------------------
\section{The 3\boldmath$\delta$ Floor: How Close Are We?}
% --------------------------------------------------------------------

\textbf{The headline.} On global lookups at $25{,}000$ nodes:
\KDHT runs at $2.45\times$ the median floor; \GDHT at
$1.35\times$; NX-17 at $\mathbf{1.16\times}$; NH-1 at
$\mathbf{1.27\times}$.

\NDHT (in both NX-17 and NH-1 forms) hugs the analytical floor.
The remaining $\sim 20\%$ overhead is structural: \NDHT takes
$\sim 4.4$ hops where an ideal protocol would take $3$; each
``extra'' hop costs roughly $\delta/2$, exactly as the geometric
series predicts.

\textbf{The scaling shape.} \KDHT's floor multiplier
\emph{grows} with $N$: $2.01\times$ at $5{,}000$ nodes,
$2.45\times$ at $25{,}000$, $2.65\times$ at $50{,}000$. This
is the wrong scaling: more nodes means more hops, each at full
average latency, each contributing fully to the total. \KDHT's
hop count grows as $\log N$ but its hop latency does not
shrink, so total latency grows.

\NDHT's multiplier is roughly stable. NX-17 holds at
$1.13$--$1.18\times$ across the $5{,}000$--$50{,}000$ range;
NH-1 at $1.18$--$1.30\times$. The mechanism is clear:
\NDHT's hop count grows logarithmically (same as Kademlia), but
each hop's average latency \emph{shrinks} with $N$ because the
synaptome's denser regional coverage gives shorter hops on
average. The two effects roughly cancel; total latency stays
near the floor.

\textbf{Regional scaling.} On the $500$\,km workload at
$25{,}000$ nodes, NH-1 is at $94$\,ms --- below the global
median 3$\delta$ floor of $204$\,ms. The floor is not a lower
bound on \emph{regional} latency, only on global. Regional
lookups have a different geometric basis: they cover only the
$\delta_{\mathrm{regional}}$ which is much smaller than the
global $\delta$. A $500$\,km lookup at $25{,}000$ nodes
benefits from $4.5$ hops at an average of $20$\,ms each --- a
total of $94$\,ms, dominated by the per-hop processing
constant rather than great-circle distance.

% --------------------------------------------------------------------
\section{The Locality Story: Where the Headlines Come From}
% --------------------------------------------------------------------

\textbf{Three regimes.} Real-world workloads are dominated by
regional traffic, but global routing is not negligible. The
benchmark cells separate three regimes.

\emph{Regional ($500$--$2{,}000$\,km).} The dominant real-world
workload pattern. \NDHT cuts latency by roughly $5$--$6\times$
versus \KDHT and $1.5$--$2\times$ versus \GDHT. The mechanism
is hop caching plus AP-scoring's latency penalty: regional
shortcuts accumulate in synaptomes; regional lookups exploit
them.

\emph{Cross-continental ($5{,}000$\,km+).} Less common in
day-to-day workloads but critical when it matters. \NDHT runs
at roughly $3\times$ the speed of \KDHT on $5{,}000$\,km, and
$1.2$--$1.4\times$ the speed of \GDHT. The mechanism is
two-hop AP lookahead: cross-continental routes are constructed
by chaining two regional hops with a transcontinental hop in
between, and lookahead lets the protocol find the right
intermediary.

\emph{Concentrated ($10\% \to 10\%$ pair).} The strongest
mechanism for \NDHT. NH-1 routes a $10\% \to 10\%$ pair in $33$\,ms
versus \KDHT's $244$\,ms --- a $\mathbf{7\times}$ speedup. The
mechanism is hop caching's tributary effect: each lookup
through an intermediate node deposits a synapse for the
destination at that node; subsequent lookups from any source
that pass through the intermediate use the shortcut.
Repeated access to the same destination set turns a $5$-hop
journey into a $1$-hop direct shortcut.

% --------------------------------------------------------------------
\section{Pub/Sub Performance and Robustness}
% --------------------------------------------------------------------

\textbf{Baseline delivery.} At $25{,}000$ nodes with $79$ topic
groups of $32$ subscribers each, NH-1 delivers $100.0\%$ of
publishes in steady state. Under $5\%$ instantaneous churn,
immediate delivery falls to $99.9\%$ and recovers to $100.0\%$
within three refresh intervals.

The publisher-to-relay path averages $4.8$ hops at $241$\,ms.
The broadcast tree depth averages $2.9$ hops at $198$\,ms with a
maximum fan-out of $31$ subscribers per relay.

\textbf{Sustained-churn graceful degradation.} Continuous
churn at $1\%$ of alive nodes per $5$ ticks, with three refresh
passes between kills, produces a smooth degradation curve:
$98.7\%$ delivered at $5\%$ cumulative churn, $91.2\%$ at
$10\%$, $86.5\%$ at $20\%$, $70.0\%$ at $25\%$, $52.4\%$ at
$30\%$. The curve is linear --- there is no cliff --- which
indicates the protocol finds a steady state where new-subscriber
recruitment matches loss.

\textbf{Slice World: the partition-recovery test.}
Table~\ref{tab:slice} reports the cross-hemisphere delivery rate
when the network is partitioned into Eastern and Western
hemispheres connected only through a single bridge node near
Hawaii.

\begin{table}[h]
\centering
\caption{Slice World single-bridge partition recovery. ``Slice
success~\%'' is the lookup success rate over a $500$-lookup
cross-partition workload after partition setup.}
\label{tab:slice}
\small
\begin{tabular}{@{}lrr@{}}
\toprule
\textbf{Protocol} & \textbf{Slice success~\%} & \textbf{Slice ms} \\
\midrule
\KDHT  & $0.0\%$ & --- (no recovery) \\
\GDHT  & $4.6\%$ & $525$ (no recovery) \\
NX-17  & $94.8\%$ & $423$ \\
NH-1   & $\mathbf{94.4\%}$ & $\mathbf{462}$ \\
\bottomrule
\end{tabular}
\end{table}

\KDHT has no learning; the partition is permanent. \GDHT routes
the geographic prefix toward the bridge for a few percent of
lookups but no mechanism amplifies the discovery. NX-17 and NH-1
both achieve $\sim 94\%$ via hop caching and triadic closure
converting the bridge into a seed crystal for new
cross-hemisphere edges.

The mechanism is the bridge as \emph{seed crystal}. After just
$10$ lookups through the partition, hop caching has installed
$\sim 20$ cross-hemisphere synapses (an average of $\sim 3$ per
successful crossing). By $500$ lookups, hundreds of
cross-hemisphere synapses exist; the partition has effectively
dissolved. The $94\%$ success is the integral of recovery, not
a snapshot of bridge-finding.

% --------------------------------------------------------------------
\section{The Highway-Percentage Sweep}
% --------------------------------------------------------------------

\textbf{The motivating question.} Real deployments will not be
uniformly browser-class. A fraction of nodes will run as
server-class peers (residential desktops with no NAT, organizational
infrastructure, deliberate ``contributing'' deployments). What
fraction is needed before the network functions well?

\textbf{The answer: surprisingly little.} The latency curve from
$0\%$ to $100\%$ highway is concave. The first $5$--$15\%$ of
server-class capacity captures roughly $70\%$ of the latency
benefit available from going to fully-server-class. Past
$15\%$, the marginal return diminishes; by $30\%$ the
performance is essentially indistinguishable from $100\%$.

\textbf{The architectural implication.} The protocol is
\emph{deployable today} on browser-class infrastructure with
modest server-class augmentation: a small constellation of
server-class peers (whether self-hosted by the network operator,
contributed by users, or run by community infrastructure
projects) is sufficient. We do not need to convince every user to
provision a fully-uncapped node; the network functions well at
$5$--$15\%$ heterogeneity.

% --------------------------------------------------------------------
\section{The geoBits = 0 Ablation}
% --------------------------------------------------------------------

\textbf{What it asks.} If we remove the S2 geographic prefix
entirely (\texttt{geoBits} $= 0$), does \NDHT still beat
\KDHT? The hypothesis (Chapter~\ref{ch:sim-tests}) is the
skeptic's: ``all the performance benefit comes from the
geographic prefix; the learning is gravy.''

\textbf{What it shows.} Removing the prefix degrades \NDHT's
performance, but the protocol still beats \KDHT by a
substantial margin.

\begin{table}[h]
\centering
\caption{Global lookup latency at $25{,}000$ nodes,
\textsc{geoBits} $= 8$ (default) vs \textsc{geoBits} $= 0$
(ablation). The drop with no prefix is what the geographic
mechanism contributes; the residual gap to Kademlia is what
the learning contributes.}
\label{tab:geobits-ablation}
\small
\begin{tabular}{@{}lrrr@{}}
\toprule
\textbf{Protocol} & \textbf{geoBits=8} & \textbf{geoBits=0} & \textbf{Δ vs Kad} \\
\midrule
\KDHT   & $511$ & $511$ & --- \\
NX-17   & $237$ & $377$ & $26\%$ faster than Kademlia \\
NH-1    & $260$ & $470$ & $8\%$ faster than Kademlia \\
\bottomrule
\end{tabular}
\end{table}

The ablation falsifies the hypothesis. With \emph{zero}
geographic information, both adaptive protocols still beat
Kademlia. The learning is doing real work, attributable
cleanly to the synaptome shaping itself by traffic.

The smaller residual advantage for NH-1 vs NX-17 in this
ablation ($8\%$ vs $26\%$) is informative: NX-17's two-tier
machinery and Markov pre-learning give it more starting
structure to exploit when the geographic prefix is removed. NH-1,
having consolidated those into the unified vitality model, has
less independent structural redundancy. The ablation shows the
\emph{cost} of the consolidation: a $\sim 5$--$10\%$ peak
performance reduction in the no-locality regime, traded for
$\sim 30\%$ parameter-count reduction and substantially better
maintainability.

% --------------------------------------------------------------------
\section{Heading Toward 50,000 and Beyond}
% --------------------------------------------------------------------

\textbf{What the 50,000-node measurements show.} The full sweep
at $50{,}000$ nodes is reported in Chapter~\ref{ch:eval-bandwidth}
in the context of bandwidth distribution. The headline
observation: every protocol's latency multiplier on the
$3\delta$ floor stays roughly constant from $25{,}000$ to
$50{,}000$, but \KDHT's gets slightly worse and \NDHT's stays
flat. The qualitative ordering does not change.

\textbf{What we have not yet measured.} $100{,}000$ nodes and
$250{,}000$ nodes. Preliminary runs at $100{,}000$ in earlier
generations indicated the protocol continues to scale; we have
not yet done a full benchmark sweep at that scale on the current
NH-1 generation. The simulator can run it; the limitation is
analyst time, not compute. Future work
(Chapter~\ref{ch:future}).


% =====================================================================
\chapter{Bandwidth Concentration: An Empirical Resolution}
\label{ch:eval-bandwidth}
% =====================================================================

\newthought{The companion red-team analysis} ranks
\emph{bandwidth saturation and congestion collapse} as a
critical deploy blocker (Issue \#3, Tier 1). Of the four Tier-1
issues, three (connection setup, RPC timeouts, latency jitter)
are textbook engineering problems with known mitigations.
Issue~\#3 was different: the simulator's uniform-cost-per-hop
model (Chapter~\ref{ch:sim-fidelity}) made it
\emph{architecturally unfalsifiable} whether \NDHT's adaptive
routing concentrated load (the success-disaster oscillation
hypothesized by the red team) or distributed it. Until that
question was answered, every other claim about the protocol
was contingent.

This chapter answers it. The instrumentation, the $4 \times 4$
sweep, the rate-knee, the churn validation, and the resulting
reclassification are reported in detail.

% --------------------------------------------------------------------
\section{Instrumentation}
% --------------------------------------------------------------------

We instrument every routing chokepoint in the simulator with
per-node send/receive counters. Two paths to cover:

\textbf{Kademlia and \GDHT} use \texttt{SimulatedNetwork.send}
for all on-the-wire traffic. The instrumentation lives there:
each call increments the sender's \texttt{msgsSent} counter and
the receiver's \texttt{msgsReceived} counter, plus a per-type
subtotal in \texttt{msgsByType}. One chokepoint, one
modification, complete coverage.

\textbf{NX-6, NX-17, and NH-1} access peers via
\texttt{nodeMap.get} directly --- the $50{,}000$-node simulator
benefits from skipping the message-dispatch layer for the
intra-node-graph operations the neuromorphic protocols
perform. Direct access bypasses
\texttt{SimulatedNetwork.send}, so per-node counters at that
chokepoint never fire for neuromorphic routing. We added a
parallel \texttt{\_msg(from, to, type)} helper invoked from each
conceptual wire site:

\begin{itemize}\setlength\itemsep{0.1em}
  \item Each lookup-hop transition (\texttt{find\_node})
  \item Each two-hop AP probe (\texttt{lookahead\_probe})
  \item Each LTP reinforcement walk step (\texttt{reinforce})
  \item Each hop caching and lateral spread write
    (\texttt{hop\_cache}, \texttt{hop\_cache\_lateral},
    \texttt{lateral\_spread})
  \item Each triadic-closure introduction
    (\texttt{triadic\_introduce}, \texttt{triadic\_probe})
  \item Each bilateral bootstrap handshake (\texttt{bootstrap})
  \item Each 2-hop neighbour scan in
    \texttt{\_localCandidate} (\texttt{local\_probe})
  \item Each pub/sub \texttt{routeMessage} forward
    (\texttt{route\_msg}) and \texttt{sendDirect} entry
    (\texttt{direct\_*})
\end{itemize}

Per-cycle counters are reset after each training session so the
captured numbers are deltas. The instrumentation is opt-in via
the simulator config and adds negligible overhead.

% --------------------------------------------------------------------
\section{The Four-by-Four Headline}
% --------------------------------------------------------------------

The $4 \times 4$ sweep: four protocols (\KDHT, \GDHT, NX-17,
NH-1) at four sizes ($5{,}000$, $10{,}000$, $25{,}000$,
$50{,}000$), $10$ training sessions of $500$ lookups each.
Steady-state averages over sessions $1$--$10$. Default parameters
throughout.

\begin{table*}[t]
\centering
\small
\caption{Steady-state per-cycle traffic distribution at
$50{,}000$ nodes. ``Total/cyc'' counts each wire message twice
(sender + receiver, the per-type counter convention). ``Hot
$>100\times$'' is the count of nodes processing more than
$100\times$ the network mean per cycle --- our operational
definition of a bandwidth-saturation hazard.}
\label{tab:bw-50K}
\begin{tabular}{@{}lrrrrr@{}}
\toprule
\textbf{Protocol} & \textbf{Total/cyc} & \textbf{Gini} &
\textbf{max/mean} & \textbf{Hot $>10\times$} & \textbf{Hot $>100\times$} \\
\midrule
\KDHT  & $14{,}463$ & $\mathbf{0.940}$ & $208\times$ & $551$ & $\mathbf{56}$ \\
\GDHT  & $17{,}633$ & $\mathbf{0.932}$ & $213\times$ & $423$ & $\mathbf{62}$ \\
NX-17  & $359{,}459$ & $0.694$ & $61\times$ & $1{,}524$ & $\mathbf{0}$ \\
NH-1   & $416{,}071$ & $0.661$ & $47\times$ & $1{,}419$ & $\mathbf{0}$ \\
\bottomrule
\end{tabular}
\end{table*}

The headline result, in three observations:

\textbf{1. The hypothesis is inverted.} Kademlia and \GDHT
develop $56$--$62$ catastrophic-bandwidth nodes at $50{,}000$
with max/mean ratios above $200\times$. \NDHT (in either form)
produces \emph{zero} such nodes at any tested scale. The
classical DHTs, not the adaptive ones, are the protocols with
the bandwidth-concentration problem at scale.

\textbf{2. Gini scaling is the cleanest signal.} Kademlia's
Gini grows monotonically with $N$: $0.876 \to 0.940$ from
$5{,}000$ to $50{,}000$. \NDHT's Gini grows more slowly:
$0.549 \to 0.661$ over the same range. Past some threshold,
classical DHTs concentrate \emph{worse} as you scale; \NDHT
tends toward a flatter floor.

\textbf{3. Volume is the cost.} \NDHT's total per-cycle
traffic is $25$--$30 \times$ Kademlia's. The cost of broad load
distribution is broader sustained activity. The next two
sections decompose where the cost goes and how to control it.

% --------------------------------------------------------------------
\section{Where the Volume Goes}
% --------------------------------------------------------------------

The per-type instrumentation lets us decompose NH-1's traffic at
$50{,}000$ nodes:

\begin{table}[h]
\centering
\caption{Per-type traffic decomposition for NH-1 at $50{,}000$
nodes. Routing-hop \texttt{find\_node} traffic is $1.3\%$ of
all traffic; the other $98.7\%$ is learning side-effects.}
\label{tab:bw-types}
\small
\begin{tabular}{@{}lrr@{}}
\toprule
\textbf{Message type} & \textbf{Count/cyc} & \textbf{\%} \\
\midrule
\texttt{local\_probe} (2-hop annealing) & $370{,}566$ & $\mathbf{89.1\%}$ \\
\texttt{lookahead\_probe} (two-hop AP) & $20{,}525$ & $4.9\%$ \\
\texttt{lateral\_spread} & $7{,}474$ & $1.8\%$ \\
\texttt{hop\_cache\_lateral} & $6{,}404$ & $1.5\%$ \\
\texttt{find\_node} (actual routing hop) & $5{,}538$ & $1.3\%$ \\
\texttt{hop\_cache} & $3{,}737$ & $0.9\%$ \\
\texttt{reinforce} (LTP feedback) & $1{,}823$ & $0.4\%$ \\
\bottomrule
\end{tabular}
\end{table}

A single mechanism --- the 2-hop neighbour scan in
\texttt{\_localCandidate}, used by simulated annealing and
dead-peer replacement --- accounts for $89\%$ of \NDHT's
wire traffic. This is the lever for reducing the volume cost
without disturbing the learning.

% --------------------------------------------------------------------
\section{The annealRateScale Knee}
% --------------------------------------------------------------------

We added a single parameter, \texttt{annealRateScale}, that
multiplies the per-hop annealing trigger probability without
disturbing temperature decay or the dead-peer-detection reheat
mechanism. Default $1.0$ preserves prior behavior. A sweep at
$50{,}000$ nodes:

\begin{table*}[t]
\centering
\small
\caption{The \texttt{annealRateScale} knee for NH-1 at
$50{,}000$ nodes. Routing quality (hops, time, success) is
unchanged across the entire range. The knee is at rate $\approx
0.10$: $5\times$ volume reduction at the lowest absolute
hot-node count.}
\label{tab:rate-knee}
\begin{tabular}{@{}rrrrrr@{}}
\toprule
\textbf{rate} & \textbf{Total/cyc} & \textbf{local\_probe} &
\textbf{Gini} & \textbf{Hot $>10\times$} & \textbf{vs default} \\
\midrule
$0.05$ & $63{,}847$ & $18{,}224$ & $0.866$ & $852$ & $6.5\times$ cut \\
$\mathbf{0.10}$ & $\mathbf{82{,}650}$ & $\mathbf{36{,}989}$ &
$\mathbf{0.834}$ & $\mathbf{575}$ & $\mathbf{5.0\times}$ \textbf{cut} \\
$0.25$ & $137{,}363$ & $92{,}453$ & $0.770$ & $735$ & $3.0\times$ cut \\
$0.50$ & $232{,}749$ & $187{,}366$ & $0.713$ & $1{,}215$ &
$1.8\times$ cut \\
$1.00$ & $414{,}929$ & $369{,}552$ & $0.661$ & $1{,}428$ & (default) \\
\bottomrule
\end{tabular}
\end{table*}

The shape is clean. \texttt{local\_probe} scales linearly with
the rate (perfect $10\times$ reduction at rate $0.10$).
\textbf{Routing quality is unchanged across the entire range:}
hops $5.48$--$5.54$, time $268$--$271$\,ms, success rate
$100\%$. The throttle does not cost performance.

The \emph{knee} is at rate $\approx 0.10$ --- the setting at
which the absolute count of $>10\times$-mean hot nodes drops to
its minimum ($575$, versus $1{,}428$ at default and $852$ at
overshoot). Past the knee, traffic concentrates on the
established hot paths (Gini rises from $0.834$ to $0.866$);
before the knee, more nodes do mild amounts of work and trip
the $10\times$ threshold. The middle is the sweet spot.

% --------------------------------------------------------------------
\section{The Robustness-Under-Churn Validation}
% --------------------------------------------------------------------

The throttle would be a Pyrrhic victory if it broke
churn-recovery. We tested. At $25{,}000$ nodes with $5\%$
per-cycle node turnover:

\begin{table*}[t]
\centering
\small
\caption{Steady-state distribution at $25{,}000$ nodes with
$5\%$ per-cycle churn. Throttled NH-1 has the smallest
churn-induced Gini increase of any tested protocol.}
\label{tab:bw-churn}
\begin{tabular}{@{}lrrrrr@{}}
\toprule
\textbf{Protocol} & \textbf{Total/cyc} & \textbf{Gini} &
\textbf{Hot $>100\times$} & \textbf{Δ Gini} & \textbf{Success} \\
\midrule
\KDHT  & $12{,}725$ & $0.919$ & $\mathbf{7}$ & $+0.006$ & $100\%$ \\
\GDHT  & $15{,}817$ & $0.910$ & $\mathbf{7}$ & $+0.006$ & $99.9\%$ \\
NX-17 (default) & $496{,}256$ & $0.678$ & $0$ & $+0.028$ & $100\%$ \\
NH-1 (default) & $356{,}358$ & $0.654$ & $0$ & $+0.037$ & $100\%$ \\
\textbf{NH-1 (rate=0.10)} & $\mathbf{71{,}695}$ & $\mathbf{0.786}$ &
$\mathbf{0}$ & $\mathbf{+0.012}$ & $\mathbf{100\%}$ \\
\bottomrule
\end{tabular}
\end{table*}

\textbf{Throttled NH-1 holds.} $100\%$ routing success,
zero $>100\times$-mean hot nodes, Gini $0.79$ (still better
than Kademlia's $0.92$), and the smallest Gini-under-churn
increase of any tested protocol ($+0.012$). The
\texttt{T\_REHEAT} mechanism (\S\ref{ch:operations}) fires on
dead-peer detection regardless of base annealing rate, so
churn-recovery is uncompromised by the throttle.

\textbf{Classical DHTs degrade.} Both Kademlia and \GDHT
develop $7$ catastrophic-bandwidth nodes under the same churn
load. The $5\%$ instantaneous turnover \emph{worsens} their
already-concentrated load distribution.

% --------------------------------------------------------------------
\section{Reclassification of Red-Team Issue \#3}
% --------------------------------------------------------------------

The red-team analysis (Chapter~\ref{ch:redteam}) ranks
bandwidth saturation as Critical / deploy blocker / 6--10
weeks of work. The proposed mitigation plan was Phase 3A--3E:
per-node load reporting, load-aware AP scoring, hot-axon-root
migration, MaxDisjoint replication, replay-cache load
awareness. The plan only made sense if the foundational
hypothesis (\emph{adaptive routing concentrates load}) was
correct.

The hypothesis is not correct. The instrumentation falsified it
empirically: at every tested scale, \NDHT distributes load
broadly (Gini $0.55$--$0.69$) while Kademlia and \GDHT
concentrate it catastrophically (Gini $0.87$--$0.94$, with
$56$--$62$ nodes processing $>100\times$ the mean at
$50{,}000$ nodes). The protocol's adaptive routing does not
amplify the success-disaster failure mode; it actively prevents
it.

The remaining concern is the Zipf-distributed pub/sub case ---
a single popular topic's root saturating from publish volume.
This is a bounded sub-problem whose mitigation (hot-axon-root
migration plus MaxDisjoint replication) was already specified
by the red team. The deploy-blocker framing of the broader issue
is retired.

This is the canonical example of falsification-driven
development (Chapter~\ref{ch:sim-philosophy}). The simulator's
job was to break the protocol; in this case, it broke the
\emph{red team's hypothesis}, not the protocol. The
reclassification was earned by the measurement, not by the
argument.

\textbf{Operational consequence.} A defensible operating point
emerged: NH-1 with \texttt{annealRateScale} $= 0.10$. At
$25{,}000$ nodes that is $5\times$ the bandwidth of Kademlia
(not $25\times$ as default), Gini $0.79$ (lower than Kademlia's
$0.92$), $100\%$ routing success, $0$ catastrophic hot nodes
even under churn. The operational argument for this protocol
no longer requires a 6--10-week mechanism build. It requires
setting a parameter.


% =====================================================================
\chapter{Architectural Wins Beyond the Headline}
\label{ch:eval-wins}
% =====================================================================

\newthought{The headline numbers} (Chapter~\ref{ch:eval-perf})
do not exhaust the protocol's value. Several architectural
properties show up indirectly --- in the protocol's behavior
under stress, in its scaling shape, in what it does \emph{not}
need --- and these properties matter as much for production
viability as the latency multipliers. This chapter walks them.

% --------------------------------------------------------------------
\section{Iterative Fallback as Graceful Degradation}
% --------------------------------------------------------------------

\textbf{The mechanism.} When a routing decision encounters no
synapse with forward XOR progress, NH-1's iterative fallback
takes the closest unvisited peer regardless of strict progress
(\S\ref{sec:navigate}). The lookup keeps moving, even at a small
cost in optimality.

\textbf{Why it matters.} Without iterative fallback, a network
partition or a churn-induced gap in the routing structure causes
\emph{permanent} lookup failure: the protocol arrives at a local
minimum and stops. With iterative fallback, the protocol takes
a sub-optimal step and continues; if a path exists, it will be
found, even if longer than ideal.

The Slice World result depends entirely on this. NH-1's $94\%$
cross-partition success rate is impossible without iterative
fallback --- the local routing structure cannot give a
forward-progress synapse across the partition until the bridge
has been used and shortcuts have been deposited, which requires
\emph{some lookup} to find the bridge first. Iterative fallback
is what lets that first lookup succeed.

The cost is small: in steady-state benign workloads, the
fallback fires on a small fraction of lookups
(typically $< 5\%$), and each fallback adds at most one extra
hop. The benefit is qualitative: \emph{the protocol degrades
gracefully under conditions the routing table does not yet
know how to handle.}

% --------------------------------------------------------------------
\section{Hop Caching as Implicit Replication}
% --------------------------------------------------------------------

\textbf{The mechanism.} On every successful lookup, every
intermediate node along the path installs a synapse pointing to
the destination (\S\ref{sec:learn}). After a path
$A \to B \to C \to D$, both $B$ and $C$ now know how to reach
$D$ in one hop.

\textbf{Why it matters.} The mechanism is \emph{implicit
replication}: a destination ID, after enough lookups, becomes
reachable from many intermediate nodes via direct synapse, not
just from its XOR-nearest neighbors. The pattern is structurally
identical to Tapestry's soft-state location pointers
(Chapter~\ref{ch:lineage}); we make it active on every successful
path rather than only on explicit \texttt{publishObject} events.

The empirical result is the $7\times$ latency advantage on
$10\% \to 10\%$ concentrated workloads (NH-1 at $33$\,ms vs
Kademlia at $244$\,ms). Repeated traffic to the same
destination set turns multi-hop paths into single-hop shortcuts.
The protocol \emph{learns the shape of the workload} and
configures itself accordingly.

% --------------------------------------------------------------------
\section{Triadic Closure as Structural Plasticity}
% --------------------------------------------------------------------

\textbf{The mechanism.} A node observing the same transit pair
$A \to (\mathit{this}) \to C$ multiple times introduces $A$
directly to $C$ (\S\ref{sec:learn}). The triangle
``$A \mathrm{-this-} C$'' is closed into the direct edge
``$A \mathrm{-} C$.''

\textbf{Why it matters.} Triadic closure is the mechanism by
which the network's \emph{communication graph} reshapes itself
to reflect the actual traffic patterns. Two peers that
consistently route through the same intermediate are revealed
as a peer pair that should know each other directly.

The mechanism is the network-level analogue of the brain's
NMDA AND-gate: NMDA detects coincidence at single synapses;
triadic closure detects coincidence at the network level, between
groups of peers that keep flowing through the same intermediate.
Both produce structural changes: NMDA strengthens existing
synapses; triadic closure creates new ones.

The empirical signature is in the Slice World recovery: the
$94\%$ cross-partition success rate would not be reachable by
hop caching alone. Triadic closure is what introduces the
peer pairs that converge on the bridge into direct relationship,
and once direct, those pairs no longer need the bridge as
intermediary. The partition dissolves through both mechanisms
working together.

% --------------------------------------------------------------------
\section{Temperature Reheat as Event-Driven Repair}
% --------------------------------------------------------------------

\textbf{The mechanism.} When a node detects a dead peer (a
synapse pointing at a node that no longer responds), the node
\emph{reheats}: temperature is raised to
$T_{\mathrm{reheat}} = 0.5$. Subsequent hops anneal more
aggressively until the temperature cools again
(\S\ref{sec:explore}).

\textbf{Why it matters.} Most adaptive DHT designs maintain
their routing table on a schedule: stabilization timers, bucket
refreshes, periodic checks. Ghinita-Teo
(Chapter~\ref{ch:lineage}) made the case that fixed-schedule
maintenance is wrong almost everywhere --- too low produces
failure spikes during churn bursts; too high wastes bandwidth
during quiet periods.

NH-1's temperature reheat solves the same problem differently.
The trigger is the event itself: a dead-peer encounter says
``the synaptome content is stale here, fix it now.'' Recovery
is local, opportunistic, and proportional to actual damage. No
periodic schedule, no global parameter to tune, no overhead
during quiet operation.

The mechanism is part of why \NDHT can sustain $5\%$ churn
indefinitely with $100\%$ routing success: the protocol doesn't
\emph{wait} for a maintenance interval to repair churn damage;
it repairs at the moment of damage, with effort proportional to
damage.

% --------------------------------------------------------------------
\section{Latency-Aware AP Scoring as Speed-First Priority}
% --------------------------------------------------------------------

\textbf{The mechanism.} AP scoring's hop selection formula
(equation~\ref{eq:ap}) combines XOR progress with observed
latency. The latency factor is exponential with a $100$\,ms
half-life: a synapse with $200$\,ms latency loses one factor
of $2$ in its score; a synapse with $400$\,ms latency loses
two factors.

\textbf{Why it matters.} Pure XOR-progress routing (Kademlia)
treats a $20$\,ms hop and a $200$\,ms hop as equivalent. AP
scoring weights the $20$\,ms hop ten-fold higher. The
protocol's regional-traffic latencies emerge from this:
candidates that happen to be physically close get strongly
preferred even when their XOR distance is slightly worse.

The mechanism interacts with the geographic prefix to produce
the locality story: the prefix biases XOR distance toward
physical proximity (so most candidates are nearby anyway), and
AP scoring's latency penalty refines the choice within the
biased candidate set. Either alone gives some locality
benefit; together they give the $5$--$6\times$ speedup of
\NDHT over \KDHT on regional workloads.

The mechanism's elegance is that it is \emph{pure routing
logic}, not a separate module. There is no
``geographic routing layer'' in NH-1; locality emerges from a
single line of the AP score formula plus the structural seed of
the prefix.

% --------------------------------------------------------------------
\section{Unified Vitality as Principled Pruning}
% --------------------------------------------------------------------

\textbf{The mechanism.} Every admission decision in NH-1 ---
when a new synapse competes for entry into a full synaptome ---
is governed by a single equation: vitality = weight × recency.
Twelve rules, twelve parameters, one admission gate.

\textbf{Why it matters.} Earlier generations had eighteen rules
across many ad-hoc admission decisions: should this hop-cache
entry evict that highway-tier entry? Should this triadic
introduction evict that bootstrap synapse? Each pairwise
decision had its own rule, its own parameters, its own
trade-offs.

The unification collapses the bookkeeping. Every entry has a
vitality score; the lowest-vitality non-LTP-locked entry is
the eviction target. The bookkeeping is uniform across all
mechanisms that introduce new synapses (hop caching, triadic,
lateral spread, anneal, dead-peer replacement, incoming
promotion). The behavior is uniform too: peers that the
network's actual traffic uses accumulate weight and recency;
peers that don't, lose their slot to peers that do.

This is the consolidation NH-1 is named for. The benchmark
numbers (Chapter~\ref{ch:eval-perf}) show that the consolidation
is essentially behavior-preserving: NH-1 trails NX-17 by
$5$--$15\%$ on absolute metrics but matches it on every
\emph{qualitative} property --- learning, locality, partition
recovery, churn resistance, pub/sub delivery.

% --------------------------------------------------------------------
\section{What \NDHT Does Not Need}
% --------------------------------------------------------------------

\newthought{The architectural wins} listed above describe
what NH-1 \emph{has}. There is a complementary list: things the
protocol does \emph{not} need that other adaptive DHT designs
do.

\begin{itemize}\setlength\itemsep{0.15em}
  \item \emph{No periodic stabilization.} Unlike Chord, Pastry,
    CAN. Repair is event-triggered.
  \item \emph{No two-tier routing structure.} Unlike NX-17 (and
    Pastry's R/L/M). One synaptome.
  \item \emph{No K-closest replication for pub/sub.} Unlike
    NX-15. One root per topic, healed by re-subscribe.
  \item \emph{No parent tracking in pub/sub.} Children
    re-subscribe upstream; that's the only state needed.
  \item \emph{No explicit failure detection.} Dead-peer detection
    is a routing primitive, used opportunistically.
  \item \emph{No global parameters.} Every parameter is per-node;
    no synchronization required.
  \item \emph{No protocol negotiation.} All nodes run the same
    code; specialization emerges from traffic.
  \item \emph{No central coordinator.} The protocol is
    \emph{strictly} peer-to-peer at every layer.
\end{itemize}

The list is itself a contribution. Each item is a piece of
machinery that some other DHT design carries because it was
needed there; \NDHT does not carry it because the unified
vitality model and the five-operation frame make it
unnecessary.

% --------------------------------------------------------------------
\section{What This Adds Up To}
% --------------------------------------------------------------------

The \emph{quantitative} picture is in
Chapter~\ref{ch:eval-perf}. The \emph{qualitative} picture, the
one this chapter has tried to make visible, is that \NDHT
produces these properties from a small set of mechanisms in
combination. No single mechanism delivers a property by
itself; they interact. Iterative fallback alone wouldn't
recover Slice World; you need iterative fallback plus hop
caching plus triadic closure plus lateral spread, all firing
together. AP scoring alone wouldn't give regional latency; you
need the prefix plus AP plus learning.

The combination is what we report. Chapter~\ref{ch:redteam}
catalogues what we still don't have, and Chapter~\ref{ch:production}
describes the closure path. The brain is production-ready;
the body is a measurable, scopeable problem.


% #####################################################################
\part{Deployment and Hardening}
% #####################################################################


% =====================================================================
\chapter{Red Team Analysis}
\label{ch:redteam}
% =====================================================================

\begin{quote}
  \itshape What would break if we shipped this tomorrow?
\end{quote}

\newthought{The simulator results} of Part~V show \NDHT working
under the conditions the simulator models. The conditions the
simulator does not model
(Chapter~\ref{ch:sim-fidelity}) are where the questions live.
The companion red-team analysis\sidenote{NH1-RedTeam-v0.3.38.md
in the project archive. Three independent red-team passes were
performed against the protocol over the course of its
development; the v0.3.38 ranking is the composite, ordered by
occurrence $\times$ severity $\times$ detectability $\times$
time-to-fix $\times$ deploy-blocker status.} is the
adversarial audit: thirteen ranked deployment-relevant gaps,
each one a way real-world deployment can fail in a way the
simulator cannot currently observe.

This chapter summarizes the audit and documents the changes
that followed from the bandwidth-concentration analysis of
Chapter~\ref{ch:eval-bandwidth}.

% --------------------------------------------------------------------
\section{Methodology of the Audit}
% --------------------------------------------------------------------

The audit identifies deployment risks through three lenses:

\textbf{1.\ What the simulator does not model.} Connection
setup cost, RPC timeouts, bandwidth limits, jitter, Byzantine
peers --- absent or trivially modeled. Anything the protocol
depends on for correctness that the simulator removes is a
candidate failure mode (Chapter~\ref{ch:sim-fidelity}).

\textbf{2.\ What prior art warns about.} Castro et al.\ (OSDI
2002), Harvesf-Blough (P2P 2007), Makris et al.\ (BIGDATA 2017),
Ghinita-Teo (IPDPS 2006), Dabek et al.\ (NSDI 2004) each
catalogue a class of distributed-system failure mode. Issues
that map onto known failure classes from this literature are
weighted heavily.

\textbf{3.\ What asymmetric or adversarial workloads expose.}
The simulator runs uniform-rate lookups on uniformly-distributed
nodes against uniformly-distributed targets. Real networks have
Zipf-distributed traffic, heterogeneous churn, asymmetric
reachability, adversarial participants. Anything whose analysis
assumes uniformity is suspect.

The audit ranks issues into three tiers: \emph{deploy blockers}
(the protocol's claimed properties cannot be verified in
production until these are addressed), \emph{correctness under
stress} (the protocol works in nominal conditions but degrades
under realistic adversarial or heterogeneous conditions), and
\emph{operational hardening} (known gaps that should be closed
before scale-out, but where staged rollout can manage the
risk).

% --------------------------------------------------------------------
\section{Tier 1: Deploy Blockers}
% --------------------------------------------------------------------

Issues 1--4 in the original ranking. Three remain; one was
reclassified following the bandwidth analysis.

\textbf{Issue 1: Connection-setup friction.} Real WebRTC
requires ICE + STUN/TURN + DTLS, taking $1.5$--$3$\,seconds per
new synapse under typical NAT conditions, with a $10$--$15\%$
ICE-failure rate in the wild. \NDHT's most important
behavioral results --- Slice World re-stitching, annealing
replacement, hop-caching cascade --- all depend on synapses
becoming usable rapidly. Production recovery times will be
substantially worse than simulator measurements suggest until
\texttt{CONNECTION\_SETUP\_MS} is modeled.

The mitigation is mostly engineering, not protocol redesign:
new synapses sit in \texttt{PENDING} state, excluded from AP
scoring until setup elapses; browser-aware concurrency cap
($\sim 4$ ICE in flight on Chrome desktop, lower on
Safari/mobile); GC-pause tolerance to suppress spurious
timeout suspicion. Estimated $4$--$6$ weeks. Priority: highest;
this is the gating issue for any production deploy claim.

\textbf{Issue 2: RPC timeouts and asynchronous black holes.}
RPCs to dead nodes return instantly in the simulator because
the simulator knows. No timeouts, no dropped packets, no
asymmetric reply paths (request goes through, reply doesn't).
The $100\%$ pub/sub recovery under $5\%$ churn assumes instant
detection; in the wild, every churn round costs multi-second
timeout windows during which messages are lost.

The mitigation is the request/reply RPC refactor:
\texttt{RPC\_TIMEOUT\_MS} (default $3{,}000$\,ms) sets a stall
timeout before iterative-fallback / next-AP-hop;
\texttt{routeMessage} traces forward and reverse paths
independently; either failure fails the RPC; reply may take a
different path back. Estimated $4$--$6$ weeks.

\textbf{Issue 3: Bandwidth saturation. \emph{Reclassified.}}

The original Tier-1 hypothesis was that AP routing concentrates
load on overloaded nodes (the success-disaster failure mode).
The empirical resolution
(Chapter~\ref{ch:eval-bandwidth}) showed the opposite: \NDHT
distributes load broadly at every tested scale ($0$
$>100\times$-mean nodes); the classical DHTs concentrate it
catastrophically at scale (\KDHT produces $56$ such nodes at
$50{,}000$ nodes). The $25\times$ volume tax is a single
tunable parameter (\texttt{annealRateScale}) with a clean
knee at rate $\approx 0.10$.

The residual concern is the \textbf{Zipf-distributed pub/sub
case}: a single popular topic's root saturating from publish
volume rather than from routing traffic. This is bounded
sub-problem; mitigation by hot-axon-root migration
(Makris et al.\ DFE pattern) plus MaxDisjoint topic
replication (Harvesf-Blough placement) is sketched in
\S\ref{sec:production-zipf}. The deploy-blocker framing of
Issue \#3 is retired.

\textbf{Issue 4: Latency jitter in AP routing.} Real RTTs
fluctuate $\pm 30\%$ from bufferbloat, asymmetric paths,
queueing. Our simulator's distance-derived latency is monotone
and clean. The LTP exponential moving average has an easy job
in the simulator; high-frequency noise could prematurely decay
good synapses or promote lucky-but-unstable ones.

Mitigation: jitter injection in the simulator (Gaussian noise
on every hop's RTT, $\sigma$ calibrated against measured
production RTT distributions); validate that LTP-EMA does not
oscillate; tune the EMA constant if it does. Estimated
$2$--$3$ weeks.

\textbf{Tier 1 summary after reclassification.} Three engineering
issues with known mitigations and bounded cost. Total
estimated work: $10$--$15$ person-weeks. The protocol's
algorithmic claims are not gated on these; they are gated on
the simulator extensions that close the
simulator/reality gap.

% --------------------------------------------------------------------
\section{Tier 2: Correctness Under Stress}
% --------------------------------------------------------------------

Issues 5--8 in the original ranking. The protocol works in
nominal conditions; under realistic adversarial or heterogeneous
workloads it degrades. Staged rollout can manage these --- but
each must be measured, not assumed.

\textbf{Issue 5: Sybil and cell-eclipse hardening.} The S2
prefix is self-declared (Chapter~\ref{ch:locality}). A node can
claim any prefix it wants. The benign failure mode is degraded
routing; the adversarial failure mode is a Sybil swarm in a
target cell (cell eclipse). Mitigations: Vivaldi RTT
integration (Sybil-resistant locality without trusting peer
self-claims); geographic proof-of-work / IP-ASN binding (raises
the cost of cell eclipse drastically); trusted-witness schemes.
Estimated $6$--$10$ weeks for an initial defense layer.

\textbf{Issue 6: Heterogeneous-churn convergence.} The
benchmark suite uses uniform-rate churn. Real networks churn
unevenly --- corporate workdays, time-of-day peaks, regional
outages. Mitigation: variable-churn benchmark inspired by
Ghinita-Teo (alternate periods of high churn with periods of
quiet, measure adaptation time at transitions); adaptive
anneal cooling rate driven by locally-observed
$(\mu, \lambda)$; liveness vs accuracy decoupled into
independent threshold-triggered channels (Ghinita-Teo style);
Patchwork-style background liveness probes (Tapestry 2004).
Estimated $4$--$6$ weeks.

\textbf{Issue 7: Byzantine pub/sub mitigation.} A malicious
peer in the axon tree can drop, duplicate, or reorder
publishes. Mitigations: MaxDisjoint topic replication
(Harvesf-Blough): replicate every axon-tree root at $d$
disjoint locations so no single node or single S2-cell eclipse
can silence a topic; Zipf-distributed publish workload
benchmark + hot-axon-root migration (Makris et al.):
threshold-triggered topic migration with DFE-style redirect at
the original location, closing the gateway-concentration
failure mode under content popularity. Estimated $6$--$8$ weeks.

\textbf{Issue 8: Asymmetric reachability and bilateral
failure.} The simulator assumes that if $A$ connects to $B$,
$B$ can reach $A$ on the same edge. Carrier-grade NAT,
asymmetric firewalls, and asymmetric bandwidth break this. Every
learning rule (LTP, hop caching, lateral spread, triadic) silently
encodes the bilateral assumption. Mitigations: per-direction
synapse fields (\texttt{forwardReachable}, \texttt{reverseReachable},
\texttt{forwardLatencyEMA}, \texttt{reverseLatencyEMA},
\texttt{asymmetryFlag}); direction-specific AP scoring; loop
detection in iterative fallback; bidirectional eviction agency
where both sides independently run admission. Estimated $4$--$6$ weeks.

% --------------------------------------------------------------------
\section{Tier 3: Operational Hardening}
% --------------------------------------------------------------------

Issues 9--13. Known gaps that should be closed before scale-out
but where staged rollout can manage the risk.

\textbf{Issue 9: Identity provisioning and key rotation.} Node
identity is a 64-bit ID derived from a public key. The
production system needs key generation, key storage, and key
rotation policy --- none of which the simulator addresses.

\textbf{Issue 10: Storage durability.} The DHT stores ephemeral
routing state, but applications layered on top will store
persistent data. The whitepaper covers routing only; the storage
layer (replication, durability, garbage collection) is its own
design problem.

\textbf{Issue 11: Replay-cache semantics under load.} The
bounded ring buffer (default 100 messages per topic) is sized
for typical workloads. High-rate topics can overflow the cache
between subscriber refresh intervals, leaving gaps. Mitigation:
dynamic cache sizing as a function of publish rate plus
explicit ``replay window expired'' responses to subscribers.

\textbf{Issue 12: Diagnostic observability.} Per-node metrics,
network-aggregate dashboards, pub/sub-specific metrics, alerting
rules, diagnosis playbooks. The simulator instruments the
counters; production needs the dashboards.
Chapter~\ref{ch:operability} is the production-grade
specification.

\textbf{Issue 13: Bootstrap rendezvous.} New nodes need at least
one known peer to start. The current simulator uses a fixed
sponsor list; production needs a published rendezvous mechanism
(DNS-based seed list, on-chain DID registry, or similar). Standard
peer-to-peer engineering; not a protocol design problem but a
deployment one.

% --------------------------------------------------------------------
\section{Total Estimated Work}
% --------------------------------------------------------------------

\textbf{To clear deploy blockers (Tier 1):} $10$--$15$
person-weeks after the bandwidth reclassification. Realistic
timeline: $4$--$6$ weeks with two engineers; $8$--$12$ weeks
with one. Issues 1, 2, 4 are independent and can be parallelized.

\textbf{To clear all 13 issues:} $30$--$50$ person-weeks total
(reduced from the pre-reclassification estimate of $65$--$90$
person-weeks because the bandwidth Phase 3 mechanism build is no
longer required). With staged rollout, Tier 2 and Tier 3 can be
addressed during early production runs without blocking initial
deployment.

\textbf{What this means for the deployment plan.} The audit's
output is not ``the protocol is unfit for production'' but
``the protocol is fit for staged production with the following
extensions in flight.'' Chapter~\ref{ch:production} is the
plan for the extensions.


% =====================================================================
\chapter{Production Pathway}
\label{ch:production}
% =====================================================================

\newthought{The red team} catalogues the gaps. This chapter
sketches the extensions that close them: a friction-realistic
simulator, learned locality via Vivaldi, adaptive parameter
tuning, MaxDisjoint replication, hot-axon-root migration, and
the deployment-staging plan that ties them together. Each
extension is described at the level of ``what it does and where
it goes,'' with detailed implementation deferred to the
engineering phase.

% --------------------------------------------------------------------
\section{Friction-Realistic Simulator Extensions}
\label{sec:production-friction}
% --------------------------------------------------------------------

The simulator is missing four classes of friction
(Chapter~\ref{ch:sim-fidelity}). Each is closed by a
parametrized extension to the simulator that produces
production-realistic versions of the existing benchmarks.

\textbf{Connection-setup time.} Add
\texttt{CONNECTION\_SETUP\_MS} ($1{,}500$--$3{,}000$\,ms with a
Gaussian distribution); new synapses enter a \texttt{PENDING}
state on creation, excluded from AP scoring until setup
elapses; ICE-failure rate of $10$--$15\%$ is modeled as a
Bernoulli flip per attempt with retry semantics.
Acceptance criteria: Slice World re-stitching time degrades
from $\sim 500$ lookups to $\sim 5{,}000$ lookups in the
extension run, but cross-partition success holds at $\geq
85\%$.

\textbf{RPC timeouts.} Add \texttt{RPC\_TIMEOUT\_MS}
($3{,}000$\,ms default); requests to dead peers stall for the
timeout window before being declared dead.
\texttt{routeMessage} is refactored into separate forward and
reverse path tracing; reply may take a different path back.
Acceptance criteria: pub/sub recovery from $5\%$ instantaneous
churn lengthens from $3$ refresh intervals to $5$--$7$, but
final delivery still reaches $100\%$.

\textbf{Bandwidth saturation.} Per-node load tracking
(\texttt{outgoingMsgRate}, \texttt{bandwidthCap},
\texttt{loadFactor}) becomes part of the AP score: effective
delay $= \mathrm{base\_delay} \times (1 + \mathrm{msg\_rate} /
\mathrm{bandwidth\_cap})$, in the form Makris et al.\ describe.
A relay at $50\%$ load gets a $0.71\times$ penalty in AP scoring;
at $80\%$, $0.57\times$; at $100\%$, $0.50\times$. Smooth
gradient away from saturated nodes without the binary
``in/out'' cliff that drives oscillation. Acceptance criteria:
Zipf-distributed pub/sub at $\alpha = 1.0$ does not produce
oscillation in any topic root's load over a $1$-hour window.

\textbf{Latency jitter.} Add Gaussian noise to per-hop RTT,
$\sigma = 0.30 \times$ mean (calibrated to production RTT
measurements). Verify that LTP-EMA does not oscillate; tune
the EMA constant if it does. Acceptance criteria: lookup
latency variance grows by $\sim 30\%$ but mean is unchanged;
LTP convergence on the optimal-path corpus does not slow by
more than $20\%$.

% --------------------------------------------------------------------
\section{Vivaldi-Style Learned Locality}
% --------------------------------------------------------------------

The S2 prefix is structural, self-declared, and a Tier-2 risk.
Vivaldi (Chapter~\ref{ch:lineage}) replaces it with a learned
synthetic coordinate that approximates RTT geometry.

\textbf{The migration.} Each peer maintains a $3$D position
vector + height. On every message exchange, the peer observes
the actual RTT and applies a small spring-force update to its
own coordinate. Confidence-weighted, decentralized, no
coordinator needed. After convergence (typically $\sim 1{,}000$
exchanges), pair-wise Euclidean distance approximates pair-wise
RTT.

The AP score's latency factor (equation~\ref{eq:ap}) is then
fed by the predicted Euclidean distance instead of the
S2-derived haversine. Existing routing logic is unchanged; the
\emph{source} of the latency estimate changes.

\textbf{Coordinate-as-prefix.} The high bits of the node ID can
be derived from the learned coordinate, replacing the S2 prefix.
This is more involved (it requires a coordinate-to-bit-prefix
discretization scheme) but it gives Sybil-resistant locality:
an attacker that misclaims must also fake RTTs, which is
expensive at scale.

\textbf{Estimated work.} $4$--$6$ weeks for the basic Vivaldi
implementation as a \emph{latency-prediction layer} on top of
the existing protocol; $4$--$8$ additional weeks for the
coordinate-as-prefix migration that retires the S2 dependence.

% --------------------------------------------------------------------
\section{Adaptive Parameter Tuning (Metaplastic NH-1)}
% --------------------------------------------------------------------

NH-1's twelve parameters are currently hand-picked. Following
Ghinita-Teo (Chapter~\ref{ch:lineage}), the parameters
themselves can adapt based on local observation.

\textbf{The local observables.} Each peer maintains a rolling
estimate of node-failure rate $\mu$ and node-join rate $\lambda$
in its synaptome neighborhood. From these plus an estimate of
network size (from successor-list density), the peer computes
the probability that a given pointer's target has died and the
probability that a new joiner now sits between this pointer and
its ideal target.

\textbf{The adaptive parameters.} Annealing cooling rate, reheat
amount, staleness threshold for synapse eviction, the LTP-lock
window --- all become functions of $(\mu, \lambda)$ rather than
fixed constants. A peer in a stable region anneals slowly; a
peer in a high-churn region anneals aggressively.

\textbf{The QoS knob.} The protocol exposes a single tunable
knob: \emph{target lookup-failure rate} $P_f$. Each peer adjusts
its own thresholds to hit it. The operational interface is the
target failure rate; the underlying parameters self-tune.

This is the natural endpoint of the path NH-1 lays out. The
metaplastic NH-1 takes the consolidation a step further: not
just collapsing eighteen rules into twelve, but eliminating the
need to manually tune the twelve. The shape would converge on
something like ``one operator dial, twelve self-tuned
parameters underneath.''

\textbf{Estimated work.} $6$--$10$ weeks for a working
implementation; longer for the calibration and stability work to
make the self-tuning actually produce the QoS-knob promise.

% --------------------------------------------------------------------
\section{MaxDisjoint Replication}
% --------------------------------------------------------------------

The pub/sub layer's residual concern is the Zipf-distributed
case: a single popular topic's root saturating from publish
volume. Two complementary mitigations are specified in the
literature, and we adopt both.

\textbf{MaxDisjoint topic replication} (Harvesf-Blough P2P 2007).
For topics whose subscriber count exceeds a threshold (say,
$1{,}000$ subscribers), replicate the root across $d = 3$
disjoint locations. Subscribers select their root by a
deterministic hash of their own ID (load-balancing across
replicas). Publishers send to all replicas in parallel.

The cost: $d \times$ the publish cost, but each replica handles
$1/d$ of the subscriber base. The benefit: scales pub/sub
linearly with replica count rather than concentrating on a
single root. Architectural escape from the success-disaster
trap.

\textbf{Hot-axon-root migration} (Makris et al.\ BIGDATA 2017,
DFE pattern). When an axon root's load factor exceeds $0.8$
for more than $60$\,seconds, the root \emph{migrates} to a
less-loaded peer in its synaptome neighborhood. The original
root installs a redirect entry; new subscribes hitting the
original root are redirected to the new one. After
$30$\,seconds the redirect is dropped and the topic is fully
migrated.

The redirect is the DFE --- a ``directory for exceptions''
indicating that the canonical placement (publisher's S2 cell)
has been overridden for hotspot mitigation.

\textbf{Estimated work.} $3$--$4$ weeks for MaxDisjoint;
$3$--$5$ weeks for hot-axon-root migration; $2$ weeks for
end-to-end testing under a Zipf workload benchmark.

\label{sec:production-zipf}

% --------------------------------------------------------------------
\section{Deployment Staging}
% --------------------------------------------------------------------

\newthought{The deploy plan} is staged. Each stage closes a
specific concern; failures at each stage block the next.

\textbf{Stage 0: simulator extensions.} Implement the four
friction-realistic extensions of \S\ref{sec:production-friction}.
Re-run the benchmark suite. The numbers from
Chapter~\ref{ch:eval-perf} will degrade --- this is intentional,
the new numbers are production-realistic. Any benchmark cell
where the degradation is more than $\sim 30\%$ is investigated
before progressing.

\textbf{Stage 1: small-scale staging.} Deploy a $5{,}000$-node
network on real hardware (cloud VMs, or a federated testbed of
contributing partner nodes). Validate that the friction-realistic
simulator predictions match. Tier-1 alerts (network partition,
lookup collapse, pub/sub failure) below threshold for $30$ days
of operation.

\textbf{Stage 2: medium-scale public.} Deploy a $50{,}000$-node
public network. Add MaxDisjoint replication for hot topics
($> 1{,}000$ subscribers). Add hot-axon-root migration. Run
under realistic workload (i.e., real applications, not
synthetic). Operability dashboards (Chapter~\ref{ch:operability})
in production.

\textbf{Stage 3: large-scale public + Vivaldi.} Migrate to
Vivaldi-style learned locality. Add adaptive parameter tuning
(metaplastic NH-1). Validate at $250{,}000$ and $1{,}000{,}000$
nodes. The protocol at this stage is the
``one-operator-dial, twelve-self-tuned-parameters'' end state.

\textbf{Total timeline.} $4$--$6$ months from clean checkout to
Stage 1 with two engineers full-time. Stage 2 within
$12$--$18$ months. Stage 3 is a research direction, not a
fixed schedule.

% --------------------------------------------------------------------
\section{What the Deployment Will Look Like}
% --------------------------------------------------------------------

The protocol's intended deployment shape is informed by the
mission (Chapter~\ref{ch:intro}): a privacy-first peer-to-peer
substrate carrying the routing and pub/sub workload of a
decentralized internet. Three properties shape the deployment.

\textbf{Browser-first.} Most peers are browsers, running
\NDHT in WebAssembly with JavaScript glue, communicating over
WebRTC data channels. The bilateral cap is $50$--$100$
connections. Average uptime is short; churn is high.

\textbf{Server-augmented.} A modest fraction
($\sim 5$--$15\%$) of peers are server-class: residential
desktops with no NAT constraints, organizational
infrastructure, contributing community nodes. The
highway-percentage analysis (Chapter~\ref{ch:eval-perf})
shows this is enough.

\textbf{Federated bootstrap.} New nodes need at least one
known peer to start. The bootstrap rendezvous is a federated
list of seed peers, published via DNS or on-chain DID
registry. No central authority; multiple operators can publish
seed lists.

The operability requirements that flow from this deployment
shape are described in Chapter~\ref{ch:operability}.


% =====================================================================
\chapter{Operability}
\label{ch:operability}
% =====================================================================

\newthought{This chapter} describes the operational surface a
production deployment of \NDHT presents: what an operator
measures, what they alert on, and what playbooks they run when
something goes wrong. The simulator instruments most of the
relevant counters already
(Chapter~\ref{ch:sim-fidelity}); this chapter promotes them to
operational dashboards and ties them to specific diagnostic
procedures.

% --------------------------------------------------------------------
\section{Per-Node Metrics}
% --------------------------------------------------------------------

Each node exposes the following metrics, exported in the standard
Prometheus / OpenTelemetry format and scraped by the operator's
monitoring stack.

\textbf{Routing metrics.}

\begin{itemize}\setlength\itemsep{0.1em}
  \item \texttt{ndht\_lookup\_total} --- counter; increments
    on every lookup originating at this node.
  \item \texttt{ndht\_lookup\_hops} --- histogram; per-lookup
    hop count.
  \item \texttt{ndht\_lookup\_latency\_ms} --- histogram;
    per-lookup wall-clock RTT.
  \item \texttt{ndht\_lookup\_failures\_total} --- counter;
    failed lookups.
  \item \texttt{ndht\_iterative\_fallback\_total} --- counter;
    lookups that triggered iterative fallback.
  \item \texttt{ndht\_synaptome\_size} --- gauge; current
    synaptome capacity used.
  \item \texttt{ndht\_synaptome\_evictions\_total} --- counter,
    labelled by mechanism (anneal / hop\_cache / triadic /
    promote).
\end{itemize}

\textbf{Bandwidth metrics.}

\begin{itemize}\setlength\itemsep{0.1em}
  \item \texttt{ndht\_msgs\_sent\_total} --- counter, labelled
    by message type (\texttt{find\_node} /
    \texttt{lookahead\_probe} / \texttt{local\_probe} / etc.).
  \item \texttt{ndht\_msgs\_received\_total} --- counter,
    labelled by type.
  \item \texttt{ndht\_msgs\_per\_second} --- derived rate
    (sender + receiver).
  \item \texttt{ndht\_temperature} --- gauge.
\end{itemize}

\textbf{Pub/sub metrics.}

\begin{itemize}\setlength\itemsep{0.1em}
  \item \texttt{ndht\_pubsub\_subscriptions} --- gauge; active
    subscriptions originating here.
  \item \texttt{ndht\_pubsub\_published\_total} --- counter;
    publishes originating here.
  \item \texttt{ndht\_pubsub\_delivered\_total} --- counter;
    publishes delivered to this node as subscriber.
  \item \texttt{ndht\_pubsub\_dropped\_total} --- counter;
    publishes that hit a full replay cache.
  \item \texttt{ndht\_axon\_children} --- gauge per topic the
    node serves as axon for; current direct-child count.
\end{itemize}

% --------------------------------------------------------------------
\section{Network-Aggregate Dashboards}
% --------------------------------------------------------------------

Operators need network-wide aggregates, not just per-node. Four
top-level dashboards.

\textbf{Routing health.} P50, p95, p99 of lookup latency across
the network. Lookup-success rate. Iterative-fallback rate.
Anomaly detection: deviation from the rolling baseline. The
operational target: p95 lookup latency $< 1.5\times$ rolling
baseline, success rate $> 99.5\%$, iterative-fallback rate
$< 5\%$. Sustained breach for $> 5$ minutes triggers an alert.

\textbf{Bandwidth distribution.} Gini coefficient of per-node
message rates across the live population. Count of nodes
exceeding $10\times$ and $100\times$ the network mean. Per-type
breakdown. The operational target: Gini $< 0.85$ (above this is
classical-DHT territory; if it climbs there, something has
forced \NDHT into a degenerate regime). Hot $>100\times$ count
$< 10$ network-wide. Persistent breaches indicate either a
parameter misconfiguration or a workload pattern that needs
hot-axon-root migration to handle.

\textbf{Pub/sub health.} Per-topic delivery rate, subscriber
count, axon tree depth. Replay-cache overflow rate. Per-topic
churn (subscriber join/leave rate). The operational target:
delivery rate $> 99\%$ on any topic with $\geq 100$
subscribers; replay-cache overflow rate $< 1\%$.

\textbf{Churn signal.} Network-wide node-failure rate, node-join
rate, average synaptome dead-peer count. The operational signal:
churn rate above $5\%$ per cycle is unusual; sustained periods
of $> 10\%$ require investigation.

% --------------------------------------------------------------------
\section{Alerting Rules}
% --------------------------------------------------------------------

\textbf{Critical alerts (page operator immediately).}

\begin{enumerate}\setlength\itemsep{0.1em}
  \item \emph{Network partition.} Sustained $> 5\%$ lookup
    failure rate across the network for $\geq 60$ seconds.
    Diagnosis playbook: \S\ref{sec:playbook-partition}.
  \item \emph{Lookup latency collapse.} P95 lookup latency more
    than $5\times$ rolling baseline. Indicates either a
    high-bandwidth-saturation event, a bug in the protocol,
    or an attacker.
  \item \emph{Pub/sub delivery failure.} Per-topic delivery rate
    drops below $90\%$ on a high-traffic topic for $\geq 5$
    minutes. Diagnosis playbook: \S\ref{sec:playbook-pubsub}.
\end{enumerate}

\textbf{Warning alerts (slack channel, follow up within 1 hour).}

\begin{enumerate}\setlength\itemsep{0.1em}
  \item \emph{Bandwidth concentration.} Hot $>100\times$-mean
    node count exceeds $10$ for $\geq 5$ minutes. Indicates
    workload skew (Zipf-distributed pub/sub) requires
    hot-axon-root migration intervention.
  \item \emph{Synaptome thrashing.} Eviction rate per node above
    a threshold (operationally, $> 1$ eviction per second
    sustained). Indicates parameters need adjustment, or churn
    rate has spiked.
  \item \emph{Latency creep.} P95 lookup latency between $1.5\times$
    and $5\times$ rolling baseline. Pre-collapse signal.
    Diagnosis playbook: \S\ref{sec:playbook-creep}.
\end{enumerate}

% --------------------------------------------------------------------
\section{Diagnosis Playbooks}
% --------------------------------------------------------------------

Three playbooks for the most-common operational failures. Each
is a structured checklist: symptoms, diagnosis steps, likely
causes, mitigations.

\subsection*{Playbook 1: Network partition}
\label{sec:playbook-partition}

\textbf{Symptom.} Lookup failure rate $> 5\%$ for $\geq 60$
seconds, distributed across many sources. Pub/sub delivery
drops on multiple topics simultaneously.

\textbf{Diagnosis steps.}
\begin{enumerate}\setlength\itemsep{0.05em}
  \item Check the network-wide \emph{churn signal} dashboard.
    Is node-failure rate elevated?
  \item Check geographic distribution of failed lookups. Are
    failures concentrated in a single S2 region?
  \item Check the bandwidth dashboard for max/mean ratio. Is
    one node carrying disproportionate traffic?
\end{enumerate}

\textbf{Likely causes.} Real network outage in a region (cable
cut, ISP failure); coordinated attack on a single S2 cell
(eclipse); hot-axon-root saturation on a popular topic.

\textbf{Mitigations.} If outage: wait for restoration; the
protocol's iterative fallback and hop caching will rebuild
cross-region edges automatically (Slice World property).
If attack: apply Tier-2 hardening (MaxDisjoint replication,
prefix-binding). If pub/sub saturation: trigger manual
hot-axon-root migration on the affected topic.

\subsection*{Playbook 2: Pub/sub delivery failure}
\label{sec:playbook-pubsub}

\textbf{Symptom.} Per-topic delivery rate below $90\%$ on a
high-traffic topic.

\textbf{Diagnosis steps.}
\begin{enumerate}\setlength\itemsep{0.05em}
  \item Check the topic's axon tree depth. Has the tree grown
    deep ($> 5$ levels) due to overflow?
  \item Check the topic's subscriber-count history. Has the
    subscriber count grown faster than the recruit-on-overflow
    rate can keep up?
  \item Check the topic's replay-cache overflow rate.
  \item Check the topic root's load factor.
\end{enumerate}

\textbf{Likely causes.} Subscriber surge (viral content);
replay-cache overflow on high-rate publish; root saturation.

\textbf{Mitigations.} For subscriber surge: ensure
recruit-on-overflow is firing (not an obvious bug); manually
trigger MaxDisjoint replication if available. For replay-cache
overflow: increase \texttt{replayCacheSize} for the affected
topic, or shorten cache time-to-live and instruct subscribers to
expect window-expired responses. For root saturation: trigger
hot-axon-root migration.

\subsection*{Playbook 3: Latency creep}
\label{sec:playbook-creep}

\textbf{Symptom.} P95 lookup latency between $1.5\times$ and
$5\times$ rolling baseline. Pre-collapse signal.

\textbf{Diagnosis steps.}
\begin{enumerate}\setlength\itemsep{0.05em}
  \item Check per-protocol per-region latency breakdowns. Is
    the creep concentrated regionally or network-wide?
  \item Check bandwidth distribution. Is one node or a small
    cluster of nodes processing disproportionate traffic?
  \item Check anneal-rate parameters; has someone recently
    changed \texttt{annealRateScale} or temperature defaults?
  \item Check the overall churn signal.
\end{enumerate}

\textbf{Likely causes.} Workload pattern shift (e.g., a new
viral topic creating concentrated traffic); parameter
misconfiguration; rising churn outpacing repair rate; transient
saturation.

\textbf{Mitigations.} Identify the workload pattern via topic
trends; trigger targeted hot-axon-root migration if a topic is
the cause. Review recent parameter changes; revert if needed. If
churn-driven, increase reheat threshold temporarily; the
protocol will catch up.

% --------------------------------------------------------------------
\section{The Production Readiness Checklist}
% --------------------------------------------------------------------

\textbf{Before Stage 1 deployment} (small-scale staging,
$\sim 5{,}000$ nodes):

\begin{itemize}\setlength\itemsep{0.05em}
  \item Friction-realistic simulator (Tier 1 issues 1, 2, 4 closed)
  \item Full benchmark suite re-run; baseline established
  \item Per-node metrics exported and dashboards built
  \item All Tier-1 critical alerts wired up with paging
  \item Three diagnosis playbooks tested (failure injection)
  \item Bootstrap rendezvous mechanism deployed
\end{itemize}

\textbf{Before Stage 2 deployment} (medium-scale public,
$\sim 50{,}000$ nodes):

\begin{itemize}\setlength\itemsep{0.05em}
  \item MaxDisjoint replication implemented for hot topics
  \item Hot-axon-root migration implemented and tested
  \item Tier-2 hardening (Sybil resistance via Vivaldi or
    proof-of-location) at least partial
  \item Operational runbook documented; on-call rotation set
  \item Operator dashboards stable and alerting-ready
\end{itemize}

\textbf{Before Stage 3 deployment} (large-scale public,
$\sim 250{,}000$ nodes):

\begin{itemize}\setlength\itemsep{0.05em}
  \item Vivaldi locality (or other learned-locality)
    implemented
  \item Adaptive parameter tuning (metaplastic NH-1) deployed
    or in flight
  \item All Tier-2 issues addressed
  \item All Tier-3 hardening complete
\end{itemize}

The progression is the production pathway. The current
\NDHT codebase, post-bandwidth-resolution, sits between
``ready for Stage 0'' (the simulator extensions) and ``ready
for Stage 1'' (small-scale real-world). The work catalogued in
Chapter~\ref{ch:redteam} is what closes that gap.


% #####################################################################
\part{Applications}
% #####################################################################


% =====================================================================
\chapter{Existing Applications: A Deep Dive}
\label{ch:apps}
% =====================================================================

\newthought{This chapter} examines each of the dozen applications
that depend on a DHT today. For each: what it does, the current
DHT layer it uses, the performance profile of that layer, where
it falls short of user expectations, and what changes if it
migrates to \NDHT. The improvement projections are back-of-envelope
from Chapter~\ref{ch:eval-perf}'s measurements, applied to the
application's typical workload pattern.

\textit{Streaming-class applications --- those that become
possible with the latency budget \NDHT provides but that don't
exist today --- are out of scope for this volume. They are the
subject of a separate forward-looking analysis.}

% --------------------------------------------------------------------
\section{BitTorrent (Mainline DHT)}
% --------------------------------------------------------------------

\textbf{What it does.} BitTorrent's Mainline DHT is the
trackerless peer-discovery layer for the BitTorrent file-sharing
protocol. A torrent's metadata (the \texttt{info\_hash}) is
mapped into the DHT keyspace; peers participating in the swarm
register themselves under that key; new peers query the DHT to
discover the swarm.

\textbf{Current implementation.} Kademlia, $160$-bit IDs, $K=8$
buckets, $\alpha = 3$ parallel queries.\sidenote{The original
BEP-0005 specification:
\url{https://www.bittorrent.org/beps/bep_0005.html}. Mainline
DHT is the largest production DHT in the world by node count,
estimated at $10$--$20$ million nodes at peak.}

\textbf{Performance.} Real-world measurements: lookup latency
from cold start is typically $5$--$30$\,seconds for a
\texttt{get\_peers} query, dominated by the iterative
$\alpha$-parallel walk through random-ID-positioned peers
spread globally. A single torrent's bootstrap can take
$10$--$60$ seconds before transfer begins. Mainline DHT's
churn rate is high (median session length under one hour) and
its routing tables drift continuously.

\textbf{Where it falls short.} The lookup latency is the most
visible. Modern torrent clients hide it behind a tracker (a
centralized fallback) for the first peer-discovery round; the
DHT is used as a backup for tracker-resistance, not the primary
path. The $5$--$30$\,second cold-start latency is unacceptable
for a primary-path discovery system, and the trackers exist
specifically to paper over it.

\textbf{What \NDHT changes.} Mainline DHT lookups translate
directly to \NDHT lookups (the same operation, the same data
model). Latency drops from $5$--$30$\,seconds to
$\sim 250$--$300$\,ms at $25{,}000$ nodes. Cold-start swarm
discovery goes from a tracker-dependent round to a primary-path
DHT round. Trackers become obsolete for peer discovery (they
remain useful as analytics endpoints, but not as required
infrastructure).

The qualitative shift: \NDHT makes the DHT \emph{the} primary
discovery layer, not a fallback. The $20$--$100\times$ latency
improvement collapses the user-visible delay below the threshold
where tracker fallback adds value.

% --------------------------------------------------------------------
\section{IPFS / Filecoin / libp2p}
% --------------------------------------------------------------------

\textbf{What it does.} IPFS (the InterPlanetary File System) is
a content-addressable peer-to-peer storage and retrieval
protocol. Each piece of content is referenced by a hash (its
content identifier, CID); the libp2p Kad-DHT routes lookups to
peers that have announced ``I have this CID.'' Filecoin is the
incentive-aligned storage layer built on the same substrate.

\textbf{Current implementation.} libp2p Kad-DHT, a Kademlia
variant. Default $K = 20$, $\alpha = 3$. Production deployment
runs on tens of thousands of nodes; published numbers vary, but
the IPFS team's measurements indicate cold-start CID lookup
latencies in the $1$--$10$\,second range under good conditions
and minutes-long failures under
poor.\sidenote{See the IPFS team's measurement studies and the
libp2p Kad-DHT specification:
\url{https://github.com/libp2p/specs/tree/master/kad-dht}.}

\textbf{Performance.} CID lookup typically averages $1$--$5$
seconds; provider-record retrieval (the second-stage operation
that locates a peer holding the content) adds another
$1$--$3$ seconds. End-to-end first-byte-of-content latency from
cold cache is in the $5$--$15$ second range.

\textbf{Where it falls short.} The latency profile makes IPFS
unviable as a CDN replacement. Centralized HTTP CDNs deliver
first-byte latency in the $50$--$200$\,ms range; IPFS is
$25$--$300\times$ slower from cold cache. The IPFS gateway
infrastructure exists to provide HTTP-fronted access for the
common case, paying for the gap with centralized servers.

There is a separate problem: IPFS pub/sub (the
\texttt{libp2p-pubsub} layer for IPNS updates and other
mutable-name resolution) uses gossipsub, which is mesh-based
and has its own scaling challenges.

\textbf{What \NDHT changes.} The DHT-layer lookups drop to
$\sim 250$--$300$\,ms (Chapter~\ref{ch:eval-perf}), comparable
to a centralized HTTP CDN's first-byte for far-away content.
End-to-end first-byte for IPFS becomes
$\sim 500$\,ms--$1$\,second --- a $10$--$20\times$
improvement. CDN-replacement viability becomes a real
question rather than a hypothetical.

The pub/sub side: \NDHT's axonal trees would replace gossipsub
for IPNS updates, with the load-distribution properties of
Chapter~\ref{ch:eval-bandwidth}. The combined effect makes
IPFS deployable on browser-class infrastructure for both
storage discovery and mutable-name updates, which it currently
is not.

% --------------------------------------------------------------------
\section{Ethereum (discv5)}
% --------------------------------------------------------------------

\textbf{What it does.} discv5 is Ethereum's peer-discovery
layer --- the system by which a new Ethereum node finds peers
to join the network. It uses a Kademlia DHT variant for
identity-aware routing, with TLS-like authenticated handshakes
on top.

\textbf{Current implementation.} Kademlia base, with extensions
for verifiable-identity ENRs (Ethereum Node Records). Lookup
operations are similar to Mainline DHT but with stricter
identity requirements.

\textbf{Performance.} Peer discovery from cold start is
typically $30$\,seconds to several minutes for a new node to
acquire enough peers to be operational. The latency is
dominated by the Kademlia walk plus the per-peer ENR
verification round trips.

\textbf{Where it falls short.} For a node joining the network
on demand (a fresh validator coming online, a new MEV searcher),
the multi-minute join time is operational friction. Mempool
gossip via gossipsub adds another scaling dimension --- transaction
propagation times in the $10$--$50$\,ms range across the
mempool require careful peer selection that the static
Kademlia table cannot adapt to.

\textbf{What \NDHT changes.} Discovery latency drops by
$10$--$50\times$ for the DHT layer. The sponsor-chain join
mechanism (\S\ref{sec:nh1-bootstrap-tick}) gives new nodes
locality-aware initial peer sets, which matters more for
Ethereum than for BitTorrent because mempool propagation is
latency-sensitive.

For mempool: \NDHT's axonal trees are a natural fit for
``broadcast a transaction to all interested validators in a
region.'' The combination of locality-aware routing plus
self-healing pub/sub trees would address the gossipsub scaling
limitation directly.

% --------------------------------------------------------------------
\section{Bitcoin Block / Transaction Propagation}
% --------------------------------------------------------------------

\textbf{What it does.} Bitcoin nodes propagate blocks and
transactions through a peer-to-peer mesh. The bootstrap is via
DNS seeds; subsequent peer discovery is via the
\texttt{addr} message exchange; propagation is via gossip.

\textbf{Current implementation.} Bitcoin does not use a DHT for
peer discovery in the strict sense; it uses a flooded peer-table
exchange and assumes that random-walk discovery will produce
adequate coverage. Block and transaction propagation use a
gossip protocol --- ``send to all peers, deduplicate via
inventory exchange.''

\textbf{Performance.} Block propagation across the global
network averages $1$--$3$ seconds for full blocks;
\emph{compact-block} propagation reduces this to a few hundred
milliseconds for the announcement plus a re-request for missing
transactions. Mempool transaction propagation latencies are in
the $50$--$200$\,ms range across the network.

\textbf{Where it falls short.} The flooded gossip is not
locality-aware: a transaction propagating from a US-east miner
to an Asian node may take a long path through Europe.
Compact-block propagation requires both peers to have the
transactions in mempool already; for novel transactions
(non-standard opcodes, large transactions) the propagation is
slow. MEV concentrates around peers with low-latency access to
miners.

\textbf{What \NDHT changes.} A locality-aware DHT for peer
discovery solves the regional-clustering problem; nodes find
nearby peers by default, transactions propagate along
short-path edges. Compact-block propagation becomes
\emph{routed-pubsub} per chain: every node subscribes to the
``block-announcement'' topic; the publisher (miner) publishes;
the axonal tree fans out with $\sim 100$\,ms latency at $50{,}000$
nodes. The structural advantage over Bitcoin's current gossip is
that the tree is built once and reused, rather than the inventory
exchange per peer per round.

The MEV concentration argument also changes: with locality-aware
routing and pub/sub, peer-to-miner latency becomes a property of
geography, not of peering relationships. The privileged-peer
problem partially dissolves.

% --------------------------------------------------------------------
\section{Tor Onion Services v3}
% --------------------------------------------------------------------

\textbf{What it does.} Tor onion services (formerly hidden
services) provide self-authenticating, location-hidden
endpoints. The descriptor lookup (which Tor relays should be
contacted to reach the service) is mediated by a hash ring of
hashing servers, called the HSDir.

\textbf{Current implementation.} A Kademlia-like DHT, but with
strong privacy constraints. Each onion-service descriptor is
stored on six HSDir servers chosen from the consensus directory.
The lookup latency is dominated by the multi-hop Tor circuit
that wraps the DHT query, not by the DHT's own performance.

\textbf{Performance.} Onion-service rendezvous typically takes
$5$--$15$\,seconds end-to-end, of which the Tor circuit setup
dominates. The descriptor lookup itself adds maybe $1$--$3$
seconds.

\textbf{Where it falls short.} The HSDir nodes are
infrastructure-tier (operated by trusted directory authorities),
which is a centralization concession. The descriptor caching is
short-lived; popular services experience descriptor-fetch storms.

\textbf{What \NDHT changes.} The descriptor lookup itself
becomes faster ($\sim 300$\,ms versus $1$--$3$\,seconds), but
the dominant cost is still the Tor circuit. The structural
benefit is that the HSDir centralization can be reduced:
\NDHT's load-distribution properties (Chapter~\ref{ch:eval-bandwidth})
mean that descriptor caching can be done at any node along the
lookup path (the hop-caching mechanism), reducing the load on
designated HSDir relays. The directory authority centralization
is not removable without a different trust model, but the
descriptor-fetch tier could be widened.

% --------------------------------------------------------------------
\section{Storj / Sia (Decentralized Storage)}
% --------------------------------------------------------------------

\textbf{What it does.} Storj and Sia are decentralized
filecoin-like storage networks. Each user's file is sharded,
encrypted, and replicated across a set of provider nodes; the
DHT mediates shard placement and retrieval.

\textbf{Current implementation.} Storj uses a Kademlia-derived
DHT for node-id and bucket discovery. Sia uses a chain-based
discovery layer; the DHT plays a smaller role.

\textbf{Performance.} Storj's shard-retrieval latency is
typically $200$--$800$\,ms for a $4$\,KB shard, dominated by
the network connection to a remote storage provider rather than
by the DHT lookup itself. The provider-discovery latency at
network edge is $\sim 1$\,second.

\textbf{Where it falls short.} Locality-aware shard placement
is the obvious next step. Storj's current placement does not
account for the user's geographic location, so a user in Berlin
may have shards stored in Sao Paulo. The retrieval pattern
exposes the random-IDs latency tax.

\textbf{What \NDHT changes.} Locality-aware placement: shards
are placed on providers in the user's S2 cell with high
probability; retrieval latency drops to the regional-routing
budget ($\sim 100$--$200$\,ms). The repair-coordination pub/sub
channel (when a shard provider goes offline, the network needs
to coordinate replacement) becomes an axonal-tree topic with
$100\%$ delivery under churn.

% --------------------------------------------------------------------
\section{Tox / Briar (Peer-to-Peer Messengers)}
% --------------------------------------------------------------------

\textbf{What it does.} Tox and Briar are peer-to-peer
messengers; users communicate directly without any
central server. Tox uses a Kademlia DHT for contact discovery
(``Bob is online and reachable at this address''); Briar uses
multiple transports including Bluetooth and a DHT-like overlay.

\textbf{Current implementation.} Tox: Kademlia, $32$-byte
public-key IDs as DHT keys. Briar: custom, hybrid local +
overlay routing.

\textbf{Performance.} Tox contact-availability check takes
$2$--$10$ seconds typically. Initial contact establishment can
take longer if the peer is behind aggressive NAT.

\textbf{Where it falls short.} Group-chat membership is
expensive: every message must be sent to every member, which
makes large groups infeasible. Push notifications (``a contact
came online'') require periodic polling because the protocol has
no proper pub/sub.

\textbf{What \NDHT changes.} Contact-discovery latency drops to
$\sim 250$\,ms (\NDHT's global lookup baseline), comparable to
centralized messengers' latency for an offline-status check.
Group chat becomes a topic in the axonal-tree pub/sub layer:
$100\%$ delivery to all subscribers with $5$\,ms-per-hop tree
fan-out, scaling to thousands of group members with bounded
cost per message. ``Contact came online'' becomes a pub/sub
event delivered automatically. The combination is the first
peer-to-peer messenger architecture that can scale to large
groups with the latency profile of a centralized service.

% --------------------------------------------------------------------
\section{YaCy (Decentralized Search)}
% --------------------------------------------------------------------

\textbf{What it does.} YaCy is a peer-to-peer search engine.
Each peer crawls a portion of the web; a custom DHT mediates
shared indexing and query distribution.

\textbf{Current implementation.} A custom DHT layer roughly
analogous to Kademlia but with index-specific extensions. Per-peer
search latency varies widely; published numbers indicate
$2$--$30$ seconds for a query depending on word count and
match coverage.

\textbf{Performance.} The latency cost is mostly index lookup
plus query distribution, both of which scale with $\log N$ on
top of a high constant (per-peer query takes $200$--$500$\,ms).

\textbf{Where it falls short.} Compared to centralized search
(Google: $\sim 200$\,ms median), YaCy is unusable for
interactive search. The user-experience gap is what has kept
YaCy a research curiosity rather than a competitive search
engine.

\textbf{What \NDHT changes.} The DHT-layer query distribution
drops to $\sim 250$\,ms; ranked-result aggregation across
shards is a separate algorithmic problem, but the DHT is no
longer the bottleneck. Combined with shard-level result caching
via hop caching, an interactive YaCy-like search engine becomes
plausible. Whether this is enough to replace Google is a
different question (centralized search has decades of ranking
infrastructure, not just speed); but the latency floor stops
being the disqualifying factor.

% --------------------------------------------------------------------
\section{Radicle (Decentralized Git Forge)}
% --------------------------------------------------------------------

\textbf{What it does.} Radicle is a peer-to-peer git
collaboration platform. Repositories are addressed by their
peer-IDs; pull requests are routed via a custom protocol; the
DHT layer handles repository discovery.

\textbf{Current implementation.} A custom DHT (gossiping libp2p
underneath); $\sim 1{,}000$-node production network.

\textbf{Performance.} Repository discovery: $1$--$3$\,seconds
typical. Change-set propagation: depends on workflow but
typically $5$--$30$\,seconds for a new commit to be visible
to all interested watchers.

\textbf{Where it falls short.} Discovery latency makes the
``find a project, browse it'' workflow feel slow compared to
GitHub. Change propagation latency makes ``watch this repo for
new commits'' poll-based rather than push-based.

\textbf{What \NDHT changes.} Repository-discovery latency drops
to $\sim 300$\,ms. Change-set propagation becomes an axonal-tree
topic (``commits-on-this-repo''), with $\sim 200$\,ms fan-out
to all subscribers. The platform becomes interactive in the
GitHub-like sense, not just browsable.

% --------------------------------------------------------------------
\section{Handshake / ENS (Decentralized DNS)}
% --------------------------------------------------------------------

\textbf{What it does.} Handshake and ENS provide decentralized
name resolution. Domain ownership is recorded on a blockchain;
a DHT-like layer caches recent lookups for performance.

\textbf{Current implementation.} ENS uses Ethereum's discv5 plus
a separate caching tier. Handshake uses a custom Kademlia-like
DHT for cache lookups.

\textbf{Performance.} Cached lookups are fast (under
$100$\,ms when hot); cold lookups go to the underlying chain
and take seconds to minutes. The cache hit rate is the dominant
performance variable.

\textbf{Where it falls short.} Cold-cache latency makes
first-time visits to a domain slow. The cache tier is itself
operated by infrastructure peers (a centralization tradeoff)
because random-Kademlia caching does not scale.

\textbf{What \NDHT changes.} Hop caching makes the ``cache
tier'' implicit: every routed lookup deposits the result at
intermediate nodes; popular domains accumulate shortcuts
naturally; the explicit caching infrastructure becomes
unnecessary. Cold-cache latency drops because the routing
itself is faster.

% --------------------------------------------------------------------
\section{WebTorrent / WebRTC P2P}
% --------------------------------------------------------------------

\textbf{What it does.} WebTorrent is BitTorrent for browsers,
running on WebRTC data channels. It uses a bridge to Mainline
DHT for peer discovery (because most browsers cannot speak the
Mainline DHT protocol directly).

\textbf{Current implementation.} A Mainline-DHT-bridge plus
WebRTC overlay. The bridge introduces a centralization point
and adds latency.

\textbf{Performance.} Peer discovery: $5$--$30$\,seconds for a
new torrent (the same as Mainline DHT, plus bridge overhead).
Browser-imposed limits restrict the swarm size and the WebRTC
connection count.

\textbf{Where it falls short.} The Mainline-DHT bridge is a
centralization concession driven entirely by the protocol's
inability to run in browsers. The bridge operators are a
single point of failure and surveillance.

\textbf{What \NDHT changes.} \NDHT \emph{is} a browser-native
DHT --- it was designed for the WebRTC environment from the
start. WebTorrent on \NDHT eliminates the bridge: peer discovery
is direct, latency is the regional-routing budget
($100$--$300$\,ms), and the browser becomes a first-class
peer-to-peer participant. The connection-cap analysis
(Chapter~\ref{ch:eval-perf}) shows this is sustainable on
browser-class infrastructure with the highway-percent
augmentation.

% --------------------------------------------------------------------
\section{Mastodon-Style Federations (No DHT Today)}
% --------------------------------------------------------------------

\textbf{What it does.} Mastodon (and the broader ActivityPub
fediverse) is a federation of independent servers. Users
follow accounts on other servers; the home server fetches and
serves the federated content.

\textbf{Current implementation.} No DHT. Discovery is via
WebFinger (a centralized handle resolution protocol per
domain). Cross-server delivery is via direct HTTP push from
publisher's home server to each subscriber's home server.

\textbf{Performance.} Per-user latency depends on home-server
performance; typical Mastodon home servers run on small VPS
hardware and exhibit $200$--$2{,}000$\,ms response times under
load. Cross-server federation can lag by minutes for popular
posts due to the per-subscriber-server fan-out.

\textbf{Where it falls short.} The home-server-as-bottleneck
pattern: a viral post on a small instance can take down the
whole instance because every other instance's followers must
fetch from the publisher's home. Discovery is brittle: WebFinger
queries hit the home server on every cold-cache lookup.

\textbf{What \NDHT changes.} A proposed \NDHT layer underneath
the fediverse would decentralize both discovery and delivery.
User-handle resolution becomes an \NDHT lookup (independent of
home-server load). Status-update delivery becomes an
axonal-tree topic per follower-graph cluster: a viral post fans
out via the tree at a bounded per-relay cost, never saturating
any single home server. Home servers continue to provide
durable storage and the WebFinger origin authority, but the
fan-out load moves to the DHT layer.

This would make Mastodon-style federations \emph{actually}
federation-friendly: small servers stop being single points of
failure for their users' follower-base reach.

% --------------------------------------------------------------------
\section{Summary of Improvements}
% --------------------------------------------------------------------

\begin{table*}[h]
\centering
\small
\caption{Cross-application summary of \NDHT improvements.
Latency numbers in this table are projections based on
Chapter~\ref{ch:eval-perf}; production deployments will need
the friction-realistic simulator extensions
(Chapter~\ref{ch:production}) to confirm.}
\label{tab:apps-summary}
\begin{tabular}{@{}lp{0.30\linewidth}p{0.40\linewidth}@{}}
\toprule
\textbf{Application} & \textbf{Current bottleneck} & \textbf{\NDHT improvement} \\
\midrule
BitTorrent (Mainline) & $5$--$30$\,s cold lookup & $\sim 300$\,ms; primary path possible \\
IPFS / Filecoin & $5$--$15$\,s first-byte & $\sim 1$\,s; CDN-replacement viable \\
Ethereum discv5 & Multi-minute join, mempool latency & $\sim 1$\,s join, locality-aware mempool \\
Bitcoin propagation & Region-blind, MEV-prone & Locality-aware, axonal-tree announce \\
Tor onion services & DHT minor; HSDir centralization & Wider HSDir tier, faster lookups \\
Storj / Sia & Region-blind shard placement & Locality-aware shards + repair pubsub \\
Tox / Briar messengers & Slow group chat, polling status & Pubsub group chat, push status \\
YaCy & $5$--$30$\,s queries & $\sim 300$\,ms DHT layer; interactive plausible \\
Radicle & $5$--$30$\,s changeset propagation & $\sim 200$\,ms via axonal pub/sub \\
Handshake / ENS & Cold-cache slow & Implicit caching via hop caching \\
WebTorrent & Mainline-DHT bridge & Bridge eliminated; browser-native \\
Mastodon-style federation & Home-server bottleneck & Decentralized fan-out via axonal trees \\
\bottomrule
\end{tabular}
\end{table*}

The pattern across applications: \emph{the DHT layer is rarely
the only bottleneck, but it is consistently a bottleneck}, and
moving to \NDHT consistently moves it out of the way of the
user-visible latency. The applications that get the most
benefit are the interactive ones (messengers, search, dev
forges) where the DHT-layer latency is in the user's
critical path. The applications that get the least benefit are
the ones where the DHT is already a small fraction of total
latency (Tor, Storj retrieval).

The applications that become \emph{architecturally different}
on \NDHT are the federation patterns (Mastodon-style),
because \NDHT's pub/sub primitive replaces the home-server
fan-out that those patterns suffer under. This is the closest
this chapter comes to a forward-looking claim, and we keep it
modest: the architecture is enabled, not necessarily preferable
in every dimension.


% #####################################################################
\part{Reference}
% #####################################################################


% =====================================================================
\chapter{Future Work}
\label{ch:future}
% =====================================================================

\newthought{The future-work agenda} follows the red-team
prioritization (Chapter~\ref{ch:redteam}) and integrates the
comparative-literature mechanisms identified in
Chapter~\ref{ch:lineage}. We distinguish \emph{near-term} (work
required for production deployment), \emph{medium-term}
(extensions that improve the protocol but are not blocking),
and \emph{long-term} (research directions that change the
protocol's operating model).

% --------------------------------------------------------------------
\section{Near-Term: The Production Pathway}
% --------------------------------------------------------------------

The work that turns NH-1 from a simulator-validated protocol
into a production system. Detailed plans in
Chapter~\ref{ch:production}; itemized here for completeness.

\textbf{Friction-realistic simulator.} Connection setup time
(\texttt{CONNECTION\_SETUP\_MS}) with new synapses in a
\textsc{pending} state until setup elapses; RPC timeouts
(\texttt{RPC\_TIMEOUT\_MS}) for silently-dropped peers, with
request/reply path tracing; load-dependent AP scoring
(Makris-style); Gaussian jitter injection on every hop. The
resulting measurements are production-realistic versions of the
results in Chapter~\ref{ch:eval-perf}.

\textbf{Operational hardening.} Per-node metrics export, network
aggregate dashboards, alerting rules, the three diagnosis
playbooks, the production readiness checklist. Most of this is
implementation work; the design is in
Chapter~\ref{ch:operability}.

\textbf{Bootstrap rendezvous.} Federated DNS-based seed lists
plus on-chain DID registry as the primary discovery mechanism for
new nodes. Standard peer-to-peer engineering, not a protocol
design issue, but it must exist before deployment.

% --------------------------------------------------------------------
\section{Medium-Term: Architectural Extensions}
% --------------------------------------------------------------------

Extensions that improve specific properties of the protocol
without changing its operating model.

\textbf{Vivaldi-style learned locality.} Replace the
self-declared S2 prefix with synthetic Euclidean coordinates
(Dabek et al., SIGCOMM 2004) that converge from observed RTTs.
The migration path is described in
\S\ref{ch:production}: a Vivaldi-style coordinate system slots
in under the AP score's latency factor without disturbing the
rest of the protocol. Closes the cell-eclipse gap
(Chapter~\ref{ch:locality}) by removing the cooperative-trust
assumption: an attacker cannot falsify a coordinate that is
verified against measured RTT.

\textbf{Byzantine pub/sub via MaxDisjoint replication.}
Replicate each axon-tree root at $d$ disjoint locations using
Harvesf-Blough's placement formula. With $d = 8$ replicas this
delivers $90\%$ pub/sub success at $50\%$ malicious nodes per
their measurement. Combined with surrogate routing on the
multicast tree, this closes the Byzantine-resistance gap
without abandoning the axonal-tree primitive.

\textbf{Adaptive synaptome capacity.} The current synaptome is
fixed at $50$. A dynamically-sized synaptome --- bumping
capacity during join-heavy periods, pruning in steady state ---
is an engineering refinement not yet attempted. The mechanism
is straightforward (capacity is just a parameter); the question
is what trigger to use and how to ensure capacity changes don't
disturb the steady state's behavior.

\textbf{Global-pool annealing.} Currently, annealing's
2-hop neighborhood scan is the entire candidate pool. A
periodic ``replace the lowest-vitality synapse with a globally
sampled candidate'' would close the gap between the
bootstrap-synaptome and the omniscient-synaptome
(Chapter~\ref{ch:eval-wins}). Tradeoff: more aggressive
exploration costs more bandwidth (the bandwidth question of
Chapter~\ref{ch:eval-bandwidth} reappears, though now with the
\texttt{annealRateScale} knob to tune it).

% --------------------------------------------------------------------
\section{Long-Term: Metaplastic NH-1}
% --------------------------------------------------------------------

The natural endpoint of the path NH-1 lays out: a protocol where
the parameters \emph{themselves} self-tune. The brain analogue
is metaplasticity --- ``plasticity of plasticity'' --- in which
the rules governing learning change based on the brain's
activity level (Abraham \& Bear, 1996). A neuron that has been
very active becomes harder to strengthen further (otherwise
everything saturates). A neuron that has been quiet becomes
easier. The learning rules adapt.

NH-1's parameters --- decay rate, protection window,
exploration rate --- are currently hand-picked constants. A
metaplastic version would let the network self-tune them based
on local conditions:

\begin{itemize}\setlength\itemsep{0.1em}
  \item Each peer maintains rolling estimates of local failure
    rate $\mu$ and join rate $\lambda$ (Ghinita-Teo style).
  \item Annealing cooling rate, reheat amount, and the
    staleness threshold for synapse eviction become functions
    of $(\mu, \lambda)$ rather than fixed constants.
  \item A peer in a stable region adapts slowly; a peer in a
    high-churn region adapts aggressively.
  \item The user / operator sets one knob: target lookup
    failure rate $P_f$. The network self-tunes the underlying
    constants.
\end{itemize}

This is the next layer of brain-inspired self-organization, and
the natural endpoint of the path the protocol lays out. The
deployment would still be NH-1's twelve-rule architecture; the
operational interface would shrink from twelve parameters to one.

% --------------------------------------------------------------------
\section{Larger-Scale Evaluation}
% --------------------------------------------------------------------

The current benchmark evaluation tops out at $50{,}000$ nodes.
We expect the protocol's qualitative behavior to extend to
larger scales but have not measured. Specific sizes worth
running:

\begin{itemize}\setlength\itemsep{0.1em}
  \item $100{,}000$ nodes --- the next doubling.
    Bootstrap behavior is visibly different at $100{,}000$
    in earlier-generation work; we should measure NH-1 there.
  \item $250{,}000$ nodes --- the production target's lower
    bound. A serious deployment should expect to scale here.
  \item $1{,}000{,}000$ nodes --- the production target's
    aspirational upper bound. We have not run NH-1 at this scale;
    the simulator can run it (with patience), but the analyst time
    to drive it has not been allocated.
\end{itemize}

The hypothesis to test at each scale: the latency multiplier on
the $3\delta$ floor stays roughly constant; the Gini coefficient
of the per-node bandwidth distribution does not climb; the
hot-$>100\times$-mean count remains zero. If any of these
hypotheses fails, the protocol's claim to ``scale-invariant''
behavior is challenged and an architectural review is required.

% --------------------------------------------------------------------
\section{Open Research Directions}
% --------------------------------------------------------------------

Beyond the medium- and long-term extensions, three open
research questions worth raising.

\textbf{Proof-of-location.} The S2 prefix is currently
cooperative trust. A verifiable location primitive (a TURN-server
handshake, an ASN-bound IP attestation, RTT-triangulation
via trusted witnesses) would prevent prefix forgery. Each
candidate has tradeoffs; none has been deeply investigated by
this project. We catalog them as a research direction rather
than a planned extension.

\textbf{Privacy under DHT routing.} The protocol's routing
discloses ``the source is asking for this target'' to every
node along the path. Tor solves this with onion encryption at
high latency cost; \NDHT does not address it directly. A privacy
extension --- onion-style relay encryption, or
Garlic-routing-style multi-message encryption, or a pure
information-theoretic technique like PIR --- is a research
question larger than this protocol's scope.

\textbf{Federation across non-uniform networks.} The protocol
assumes a single network with consistent S2 cell space. In a
real-world deployment, multiple operators might run separate
\NDHT networks (bridged through specific peers). The federation
semantics ---  how lookup requests cross between networks, what
identity-resolution model spans them --- is not currently
specified.

% --------------------------------------------------------------------
\section{The Through-Line}
% --------------------------------------------------------------------

\newthought{The work catalogued} in this chapter falls into
three groups, and the groups have a natural sequence.

The \emph{near-term} group is engineering: implement the
simulator extensions, set up the operational tooling, finish the
production-readiness checklist. The protocol design is fixed;
the work is to make it deployable.

The \emph{medium-term} group is principled extension: Vivaldi
locality, Byzantine pub/sub, adaptive synaptome capacity. Each
can be evaluated against the existing benchmark suite; each adds
a measurable property to the protocol's repertoire.

The \emph{long-term} group is metaplastic: the protocol learns
to tune itself, the operator's interface shrinks to one knob,
the protocol becomes \emph{self-administering}. This is the
direction the consolidation work in NH-1 was always pointing
toward, and it is the direction the deployment trajectory in
Chapter~\ref{ch:production} is preparing for.

The mission --- a privacy-first, decentralized internet
substrate --- is realized progressively, not in a single
release. Each stage of this work shrinks the gap between what
the protocol does and what the mission requires.


% =====================================================================
\chapter*{References}
\label{ch:refs}
% =====================================================================
\addcontentsline{toc}{chapter}{References}

\section*{Foundational DHT Work}

\begin{itemize}\setlength\itemsep{0.2em}
  \item Maymounkov, P., \& Mazi\`eres, D. ``Kademlia: A
    Peer-to-Peer Information System Based on the XOR Metric.''
    \emph{Proceedings of IPTPS}, 2002.
  \item Stoica, I., Morris, R., Karger, D., Kaashoek, M.~F., \&
    Balakrishnan, H. ``Chord: A Scalable Peer-to-Peer Lookup
    Service for Internet Applications.'' \emph{SIGCOMM},
    2001.
  \item Ratnasamy, S., Francis, P., Handley, M., Karp, R., \&
    Schenker, S. ``A Scalable Content-Addressable Network.''
    \emph{SIGCOMM}, 2001.
  \item Rowstron, A., \& Druschel, P. ``Pastry: Scalable,
    decentralized object location and routing for large-scale
    peer-to-peer systems.'' \emph{Middleware}, 2001.
  \item Zhao, B.~Y., Huang, L., Stribling, J., Rhea, S.~C.,
    Joseph, A.~D., \& Kubiatowicz, J.~D. ``Tapestry: A Resilient
    Global-Scale Overlay for Service Deployment.''
    \emph{IEEE JSAC} 22(1), 2004.
\end{itemize}

\section*{Latency-Aware DHTs}

\begin{itemize}\setlength\itemsep{0.2em}
  \item Dabek, F., Li, J., Sit, E., Robertson, J., Kaashoek,
    M.~F., \& Morris, R. ``Designing a DHT for low latency and
    high throughput.'' \emph{NSDI}, 2004.
  \item Castro, M., Druschel, P., Kermarrec, A., \& Rowstron, A.
    ``SCRIBE: A large-scale and decentralized
    application-level multicast infrastructure.''
    \emph{IEEE JSAC}, 2002.
  \item Freedman, M.~J., Freudenthal, E., \& Mazi\`eres, D.
    ``Democratizing Content Publication with Coral.''
    \emph{NSDI}, 2004.
\end{itemize}

\section*{Adaptive Coordinates}

\begin{itemize}\setlength\itemsep{0.2em}
  \item Dabek, F., Cox, R., Kaashoek, M.~F., \& Morris, R.
    ``Vivaldi: A Decentralized Network Coordinate System.''
    \emph{SIGCOMM}, 2004.
  \item Ledlie, J., Gardner, P., \& Seltzer, M. ``Network
    Coordinates in the Wild.'' \emph{NSDI}, 2007.
\end{itemize}

\section*{Byzantine Resistance and Route Diversity}

\begin{itemize}\setlength\itemsep{0.2em}
  \item Castro, M., Druschel, P., Ganesh, A., Rowstron, A., \&
    Wallach, D.~S. ``Secure routing for structured peer-to-peer
    overlay networks.'' \emph{OSDI}, 2002.
  \item Harvesf, C., \& Blough, D.~M. ``The Effect of Replica
    Placement on Routing Robustness in Distributed Hash
    Tables.'' \emph{IEEE P2P}, 2007.
\end{itemize}

\section*{Load Balancing and Hotspot Mitigation}

\begin{itemize}\setlength\itemsep{0.2em}
  \item Makris, C., Tserpes, K., \& Anagnostopoulos, D. ``A novel
    DHT placement strategy for hot data using directory of
    exceptions.'' \emph{IEEE Big Data}, 2017.
  \item Godfrey, B., Lakshminarayanan, K., Surana, S., Karp, R.,
    \& Stoica, I. ``Load Balancing in Dynamic Structured P2P
    Systems.'' \emph{INFOCOM}, 2004.
\end{itemize}

\section*{Adaptive Maintenance and Churn Handling}

\begin{itemize}\setlength\itemsep{0.2em}
  \item Ghinita, G., \& Teo, Y.~M. ``An Adaptive Stabilization
    Framework for Distributed Hash Tables.'' \emph{IPDPS},
    2006.
  \item Saroiu, S., Gummadi, P.~K., \& Gribble, S.~D. ``A
    Measurement Study of Peer-to-Peer File Sharing Systems.''
    \emph{MMCN}, 2002.
\end{itemize}

\section*{Pub/Sub Substrates}

\begin{itemize}\setlength\itemsep{0.2em}
  \item Castro, M., Druschel, P., Kermarrec, A.-M., Nandi, A.,
    Rowstron, A., \& Singh, A. ``SplitStream: High-bandwidth
    multicast in cooperative environments.'' \emph{SOSP},
    2003.
  \item Ratnasamy, S., Handley, M., Karp, R., \& Shenker, S.
    ``Application-level multicast using content-addressable
    networks.'' \emph{NGC}, 2001.
\end{itemize}

\section*{Substrate: Learning, Decay, Pruning}

\begin{itemize}\setlength\itemsep{0.2em}
  \item Hebb, D.~O. \emph{The Organization of Behavior.}
    Wiley, 1949.
  \item Bliss, T.~V.~P., \& L\o mo, T. ``Long-lasting potentiation
    of synaptic transmission in the dentate area of the
    anaesthetized rabbit following stimulation of the
    perforant path.'' \emph{Journal of Physiology} 232,
    1973.
  \item Frey, U., \& Morris, R.~G.~M. ``Synaptic tagging and
    long-term potentiation.'' \emph{Nature} 385(6616), 1997.
  \item Abraham, W.~C., \& Bear, M.~F. ``Metaplasticity: the
    plasticity of synaptic plasticity.'' \emph{Trends in
    Neurosciences} 19(4), 1996.
  \item Drachman, D.~A. ``Do we have brain to spare?''
    \emph{Neurology} 64(12), 2004.
\end{itemize}

\section*{Geometry and Geography}

\begin{itemize}\setlength\itemsep{0.2em}
  \item Google. ``S2 Geometry Library.''
    \url{https://s2geometry.io/}, 2011.
  \item Hilbert, D. ``\"Uber die stetige Abbildung einer Linie auf
    ein Fl\"achenst\"uck.'' \emph{Mathematische Annalen}
    38(3), 1891.
\end{itemize}


\appendix


% =====================================================================
\chapter{NX-1 to NX-17: Detailed Evolution}
\label{app:nx-evolution}
% =====================================================================

\newthought{This appendix} traces the experimental sequence
that produced NH-1. Each NX generation either added a mechanism
that survived consolidation, added a mechanism that was
later retired, or refactored an existing mechanism into a
cleaner form. The list below is concise; the full simulator
codebase preserves each generation under
\texttt{src/dht/neuromorphic/}, and Appendix~B's parameter
tables include the configurations used.

% --------------------------------------------------------------------
\section{Capsule Per Generation}
% --------------------------------------------------------------------

\textbf{NX-1 (the routing-table Hebbian baseline).}
Single-tier synaptome with weighted entries. LTP increments
weights on successful paths. Fixed-rate decay. Bootstrap by
random sampling. \emph{Established that simple Hebbian
reinforcement on a routing table improves over Kademlia.}

\textbf{NX-3 (geographic three-layer initialization).}
Introduced geographic-prefix node IDs (\GDHT). Bootstrap
allocates entries across three layers: intra-cell local,
inter-cell structured (one per geographic-prefix bucket),
random global. \emph{Locality became part of the protocol.}

\textbf{NX-4 (iterative fallback).} When forward XOR-progress
candidates are exhausted, take the closest unvisited peer
regardless of strict progress. Enables routing through
partitions and around routing-table gaps. \emph{The protocol
stopped failing on partitioned topologies.}

\textbf{NX-5 (stratified bootstrap and incoming-synapse
promotion).} Bootstrap allocates entries with explicit per-stratum
floors; new node coverage of every XOR-distance bucket is
guaranteed. Incoming synapses (passive learning) added: peers
that reach this node are recorded and promoted on use. \emph{The
synaptome population became routing-quality balanced from the
start.}

\textbf{NX-6 (churn-resilient routing).} Temperature reheat on
dead-peer detection. Dead-synapse eviction and replacement.
Bootstrap-preference in the eviction gate. \emph{Churn recovery
became event-triggered, eliminating the need for periodic
stabilization.}

\textbf{NX-7 (dendritic pub/sub v1).} First attempt at pub/sub:
$25\%$ peel-off relay tree. \emph{Worked, but the $25\%$
threshold was an ad-hoc choice; tree balance was uneven.}

\textbf{NX-8 (dendritic pub/sub v2).} Balanced binary-split
relay tree. \emph{Better balance but the splitting logic was
brittle under churn.}

\textbf{NX-9 (geographic dendritic pub/sub).} S2-clustered relay
tree with direct 1-hop delivery. \emph{Locality benefit on pub/sub,
but the tree was still fragile.}

\textbf{NX-10 (routing-topology forwarding tree).} Pub/sub tree
follows the routing topology directly: every first-hop synapse
is a forwarder. Strong load distribution, but tree was rigid.
\emph{The routing topology and the pub/sub topology should be
the same thing.}

\textbf{NX-11 (diversified bootstrap + axonal pub/sub).}
$80/20$ stratified/random bootstrap. Axonal-tree pub/sub. The
first generation that combined good routing with good pub/sub
in a unified architecture.

\textbf{NX-12 (load-tracked AP scoring).} Experimental
load-aware AP scoring, where AP factor includes a load
penalty. \emph{Did not survive ablation: at the scales tested,
the load penalty was either redundant (vitality already
preferred under-used peers) or destabilizing (oscillation when
the load measure was noisy). Retired.}

\textbf{NX-13 (G-DHT three-layer init refinement).} Refined
three-layer bootstrap. Mostly a parameter sweep over NX-11.

\textbf{NX-15 (AxonManager abstraction).} Introduced the
\texttt{AxonManager} abstraction over pub/sub state, separating
the topic-membership logic from the routing protocol. Enabled
pluggable pub/sub policies (recruit-on-overflow, relay
selection).

\textbf{NX-17 (publisher-prefix topic IDs and routed axonal
trees).} Final pre-NH-1 generation. Topic IDs include the
publisher's S2 cell prefix; pub/sub uses pure routed-tree mode
(no K-closest replication). Eighteen rules and forty-four
parameters. \emph{State-of-the-art performance, but the rule
set was at the edge of comprehensibility.}

\textbf{NH-1 (the consolidation).} Audit of NX-17. Six rules
absorbed or retired (highway tier into vitality;
stratum-floor into bilateral cap; Markov pre-learning into
triadic; load-tracked AP retired; relay-pinning into standard
recruit; \texttt{weightScale} into removal). Twelve rules
remain; twelve parameters; one equation (vitality) governing
admissions. \emph{Same behavior as NX-17 within $5$--$15\%$
on every benchmark; substantially more maintainable.}

% --------------------------------------------------------------------
\section{Lessons from the Sequence}
% --------------------------------------------------------------------

Three patterns are visible across the NX-1 → NH-1 progression.

\textbf{Most generations added one mechanism.} Each was
ablation-tested before being committed; mechanisms that
didn't survive ablation were either reverted or marked for
retirement at the next consolidation.

\textbf{The biggest gains came from removing things, not adding
them.} The NX-17 → NH-1 audit removed six rules and
thirty-two parameters with essentially no behavioral cost. The
``learning'' is in part learning what \emph{not} to keep.

\textbf{Pub/sub took five generations to get right} (NX-7
through NX-11). Each iteration tried a different tree topology;
the winning answer (axonal trees rooted at publisher-prefix
topic IDs, with self-healing via routed re-subscribe) emerged
only after four failed attempts. This is normal in
distributed-systems research; it is documented here so that
future readers understand why the current design is the way
it is.


% =====================================================================
\chapter{Full Parameter Tables}
\label{app:params}
% =====================================================================

\newthought{This appendix} consolidates every constant in the
NH-1 implementation: the twelve tunable parameters from
Chapter~\ref{ch:nh1-rules} (recapitulated here for one-stop
reference), the routing-layer fixed constants, the pub/sub
operational tunables, and the simulator-side configuration.

% --------------------------------------------------------------------
\section{Tunable Parameters}
% --------------------------------------------------------------------

The twelve operational parameters of NH-1 (also listed in
Table~\ref{tab:nh1-params}):

\begin{table*}[h]
\centering
\small
\begin{tabular}{@{}lrlp{0.46\linewidth}@{}}
\toprule
\textbf{Parameter} & \textbf{Default} & \textbf{Legal range} & \textbf{Effect} \\
\midrule
\texttt{maxSynaptome}    & $50$    & $20$--$256$  & Capacity tradeoff vs WebRTC connection budget \\
\texttt{lookaheadAlpha}  & $5$     & $1$--$10$    & 2-hop probe breadth; cost is per-hop network traffic \\
\texttt{maxHops}         & $40$    & $20$--$100$  & Hard ceiling; rarely reached in practice \\
\texttt{epsilon}         & $0.05$  & $0$--$0.20$  & First-hop random-choice probability \\
\texttt{annealCooling}   & $0.9997$ & $0.99$--$0.9999$ & Per-hop temperature decay rate \\
\texttt{annealRateScale} & $1.0$   & $0$--$1.0$   & Master volume on annealing (the bandwidth knob) \\
\texttt{decayGamma}      & $0.995$ & $0.99$--$0.9999$ & Per-decay-tick weight multiplier \\
\texttt{vitalityFloor}   & $0.05$  & $0$--$0.5$   & Below this, freely replaceable regardless of recency \\
\texttt{inertiaDuration} & $20$    & $5$--$100$   & LTP-lock window; protects fresh reinforcements \\
\texttt{promoteThreshold} & $2$    & $1$--$10$    & Incoming-synapse promotion trigger \\
\texttt{triadicThreshold} & $2$    & $1$--$10$    & Triadic-closure firing threshold \\
\texttt{recencyHalfLife} & $50$    & $20$--$200$  & Tick half-life of post-inertia recency decay \\
\bottomrule
\end{tabular}
\end{table*}

% --------------------------------------------------------------------
\section{Routing-Layer Fixed Constants}
% --------------------------------------------------------------------

These are not intended for routine tuning; they are
implementation choices documented here for completeness.

\begin{itemize}\setlength\itemsep{0.1em}
  \item \textsc{T\_INIT} $= 1.0$ --- initial node temperature.
  \item \textsc{T\_MIN} $= 0.05$ --- temperature floor.
  \item \textsc{T\_REHEAT} $= 0.5$ --- temperature on dead-peer
    detection.
  \item \textsc{DECAY\_INTERVAL} $= 100$ --- lookups per decay
    tick.
  \item \textsc{ANNEAL\_LOCAL\_SAMPLE} $= 50$ --- per-call cap
    on \texttt{\_localCandidate} probes.
  \item \textsc{GEO\_BITS} $= 8$ --- S2 cell prefix width.
  \item \textsc{GEO\_REGION\_BITS} $= 4$ --- granularity for
    lateral spread.
  \item \textsc{STRATA\_GROUPS} $= 16$ --- bucketing for
    stratum-diversity logic.
  \item \textsc{LATERAL\_K} $= 2$ --- lateral-spread fanout.
\end{itemize}

% --------------------------------------------------------------------
\section{Pub/Sub Operational Tunables}
% --------------------------------------------------------------------

Five parameters governing the axonal-tree pub/sub layer:

\begin{itemize}\setlength\itemsep{0.1em}
  \item \texttt{maxDirectSubs} $= 32$ --- max direct children
    per axon before recruit-on-overflow fires.
  \item \texttt{minDirectSubs} $= 4$ --- min direct children
    before relinquishing role.
  \item \texttt{refreshIntervalMs} $= 30000$ --- per-member
    refresh interval.
  \item \texttt{maxSubscriptionAgeMs} $= 90000$ --- TTL for
    child entries before sweep.
  \item \texttt{replayCacheSize} $= 100$ --- ring buffer
    messages per axon-role.
\end{itemize}

% --------------------------------------------------------------------
\section{Simulator Configuration}
% --------------------------------------------------------------------

The simulator's per-run configuration:

\begin{itemize}\setlength\itemsep{0.1em}
  \item \texttt{nodeCount} --- total nodes (typically
    $5000$--$50000$ in this volume).
  \item \texttt{maxConnections} $= 100$ --- bilateral cap.
  \item \texttt{highwayPct} $= 0$ --- fraction of
    server-class peers (default $0\%$ uniform browser-class).
  \item \texttt{webLimit} $=$ true --- enforce bilateral cap.
  \item \texttt{lookupsPerCell} $= 500$ --- benchmark sample
    size per workload type.
  \item \texttt{warmupSessions} --- per-protocol warmup before
    measurement (variable; see methodology).
  \item \texttt{churnRate} --- percentage replacement per
    cycle in churn tests.
  \item \texttt{nodeDelayMs} $= 10$ --- per-hop processing
    constant.
\end{itemize}


% =====================================================================
\chapter{Extended Glossary}
\label{app:glossary}
% =====================================================================

\newthought{This appendix} repeats the glossary from
Chapter~\ref{ch:vocab}, extended with cross-references to source
code symbols and chapter sections. A reader who knows what they
are looking for and wants to find it in the codebase or in the
chapter where the concept is fully developed should start here.

The structure: term, one-line definition, source location
(implementation symbol or chapter / section), and forward
references where relevant.

\begin{description}\setlength\itemsep{0.4em}

\item[\textsc{addByVitality}] The unified admission gate for
synapses. Implementation:
\texttt{NeuromorphicDHTNH1.\_addByVitality()}. Algorithm in
Chapter~\ref{ch:nh1-routing}.

\item[\textsc{annealRateScale}] The bandwidth-tuning master
volume on annealing. Implementation: \texttt{r.annealRateScale}
in NH-1 constructor; default $1.0$. Knee analysis in
Chapter~\ref{ch:eval-bandwidth}.

\item[\textsc{AP scoring}] The hop-selection mechanism.
Equation~\ref{eq:ap}. Implementation:
\texttt{NeuronNode.bestByAP()} for single-hop;
\texttt{NeuromorphicDHTNH1.\_bestByTwoHopAP()} for two-hop
lookahead. Discussion in Chapter~\ref{ch:operations}.

\item[\textsc{axon}] A node serving as a relay for a pub/sub
topic. Implementation: \texttt{AxonManager} per node, with a
\texttt{role} record per topic. Chapter~\ref{ch:pubsub}.

\item[\textsc{axonal tree}] The complete delivery structure for
a topic. Built by \texttt{routeMessage} subscribes; healed by
periodic re-subscribe. Chapter~\ref{ch:nh1-pubsub}.

\item[\textsc{bilateral cap}] Bilateral connection-cap
enforcement. Implementation: \texttt{node.tryConnect(peer)}
returns false if either side at capacity.
Chapter~\ref{ch:sim-fidelity}.

\item[\textsc{Dabek $3\delta$ floor}] The analytical lower bound
on lookup latency in any recursive DHT. Calibrated in
methodology, Chapter~\ref{ch:eval-method}.

\item[\textsc{decay}] Weight erosion via multiplicative
$\gamma$. Implementation: \texttt{\_tickDecay()}.
Chapter~\ref{ch:operations} (\textsc{forget}).

\item[\textsc{Gini coefficient}] Per-node bandwidth
distribution inequality measure. Definition in
Chapter~\ref{ch:vocab}; analysis in Chapter~\ref{ch:eval-bandwidth}.

\item[\textsc{hop caching}] Every intermediate node along a
successful path remembers the destination. Implementation:
\texttt{\_hopCache()}. Chapter~\ref{ch:operations}
(\textsc{learn}).

\item[\textsc{incoming synapse}] A passive learning record of
peers that have reached this node. Implementation:
\texttt{node.incomingSynapses}. Chapter~\ref{ch:substrate}.

\item[\textsc{inertia}] The protection window after a synapse is
reinforced. \textsc{INERTIA\_DURATION} $= 20$ ticks.
Chapter~\ref{ch:operations}.

\item[\textsc{iterative fallback}] Routing when forward XOR
progress is exhausted. Implementation: scan unvisited synapses,
take closest. Chapter~\ref{ch:operations}.

\item[\textsc{$k$-bucket}] Kademlia distance-bucket structure.
Bucket-organized at bootstrap; entries migrate via traffic.
Chapter~\ref{ch:lineage}.

\item[\textsc{lateral spread}] Hop-cache propagation to
geographic neighbors. Implementation:
\texttt{\_hopCache()} recursion at depth 1.
Chapter~\ref{ch:operations}.

\item[\textsc{LTP / LTD}] Long-term potentiation /
depression. The biological basis for weight reinforcement and
decay. Chapter~\ref{ch:substrate}.

\item[\textsc{N-DHT}] Generic family name for the neuromorphic
DHT. Chapter~\ref{ch:intro}.

\item[\textsc{NH-1}] The current generation; consolidation of
NX-17. Twelve rules, twelve parameters. Chapter~\ref{ch:nh1-rules}.

\item[\textsc{node ID}] 64-bit identifier. High 8 bits = S2
cell prefix; low 56 bits = $H_{56}$(public key).
Chapter~\ref{ch:locality}.

\item[\textsc{publisher-prefix topic ID}] Topic ID with high 8
bits set to publisher's S2 prefix.
Chapter~\ref{ch:pubsub}.

\item[\textsc{recency}] Exponential-decay multiplier in
vitality. Within inertia: $1.0$. Outside: $\max(0.1,
e^{-\Delta t / \tau})$. Chapter~\ref{ch:operations}.

\item[\textsc{replay cache}] Bounded ring buffer of recent
publishes per topic. Chapter~\ref{ch:pubsub}.

\item[\textsc{S2 cell prefix}] Geographic encoding via Hilbert
curve. \textsc{geoBits} $= 8$ default.
Chapter~\ref{ch:locality}.

\item[\textsc{Slice World}] Single-bridge partition recovery
test. Implementation: \texttt{src/dht/sliceWorld.js}.
Chapter~\ref{ch:sim-tests}.

\item[\textsc{stratified bootstrap}] Bucket-balanced bootstrap
with explicit geographic-strata fill. Implementation:
\texttt{NeuromorphicDHTNH1.buildRoutingTables()}.
Chapter~\ref{ch:locality}.

\item[\textsc{stratum}] An XOR-distance bucket; equivalent to
``Kademlia bucket index.''

\item[\textsc{synapse}] Single peer entry in the synaptome.
Implementation: \texttt{Synapse} class.
Chapter~\ref{ch:substrate}.

\item[\textsc{synaptome}] Routing table; the set of synapses a
node maintains. Implementation: \texttt{node.synaptome}.

\item[\textsc{temperature}] Per-node scalar governing annealing
trigger probability. Cools by \textsc{ANNEAL\_COOLING};
reheats on dead-peer detection.
Chapter~\ref{ch:operations}.

\item[\textsc{triadic closure}] Network-level coincidence
detection: same transit pair $\to$ direct introduction.
Chapter~\ref{ch:operations} (\textsc{learn}).

\item[\textsc{two-hop lookahead}] AP-scoring extension that
explores two hops of candidates. Implementation:
\texttt{\_bestByTwoHopAP()}. Chapter~\ref{ch:operations}.

\item[\textsc{vitality}] \textsc{weight} $\times$
\textsc{recency}. The single equation governing every
admission decision. Equation~\ref{eq:vitality}.
Chapter~\ref{ch:operations}.

\item[\textsc{weight}] Strengthening factor in $[0, 1]$.
Reinforced by LTP, decayed multiplicatively.
Chapter~\ref{ch:operations}.

\end{description}


% =====================================================================
\backmatter
% =====================================================================

\end{document}
